% concatenated from hnn10.tex
\documentclass[a4paper,12pt]{article}

%\documentclass{amsart}


\usepackage[english]{babel}
\usepackage[T1]{fontenc}
\usepackage{amssymb}
\usepackage{amsmath}
\usepackage{amsthm}
\usepackage{mathrsfs}
\usepackage[all]{xy}
\newcounter{continua}
\usepackage{ulem}%%%riscado em latex
%%%%%%%%%%%%%%%%%%%%%%%%%%%%%%%%%%%%%%este %\usepackage{sectsty}

\usepackage{sectsty}

% Para dekfinir a fonte da seção
%%%%%%%%%%%%%%%%%%%%%%%%%%%%%%%% este \sectionfont{\fontsize{14}{15}\selectfont}

\sectionfont{\fontsize{14}{15}\selectfont}

%%%%%%%%%%%%%%%%%%%%%%%%%%%%%%%%%%%%%%%%%%%%%%%%%%%%%%%%%%%%%%%%%%%%%%%%%%%%%%%%%%%%%%%%%%%%%%%%

%\theoremstyle{thmit} % Numbered and Italic

\newtheorem{thm}{Theorem}[section]
\newtheorem{lem}[thm]{Lemma}
\newtheorem{cor}[thm]{Corollary}
\newtheorem{prop}[thm]{Proposition}
\newtheorem{question}{Question}
\newtheorem{deff}[thm]{Definition}
\newcommand{\Z}{\mathbb{Z}}
\newcommand{\Q}{\mathbb{Q}}
\newcommand{\f}{\mathbb{F}}
\newcommand{\A}{{\cal A}}
\newcommand{\IT}{{\cal IT}}
\newcommand{\Oe}{\rm \widehat{OE}}
\def\OE{{\rm OE}}
\newcommand{\B}{{\cal B}}
\newcommand{\C}{{\cal C}}
\newcommand{\F}{{\cal F}}
\newcommand{\G}{{\cal G}} 
\newcommand{\W}{{\cal W}}
\newcommand{\N}{{\cal N}}
\newcommand{\I}{{\cal I}}
\newcommand{\K}{{\cal K}}
\newcommand{\cR}{{\cal R}}
\newcommand{\g}{{\bf g}}
\def\aut {{\rm Aut}}
\def\out {{\rm Out}}
\def\inn {{\rm Inn}}
\def\hnn {{\rm HNN}}
\def\Iso {{\rm Iso}}
\def\hom {{\rm hom}}
\def\mono {{\rm Mono}}
\def\orbita{{\rm orbit}}
\def\orbitas{{\rm orbits}}
%\newcommand{\G}{\widehat{G}_{p}}
\def\f{{\mathfrak F}}
\def\g{{\bf g}}
\def\gp{{\bf g}_p}
\def\Q{{\mathbb Q}}
\def\CQ{\overline{\Q}{}}
\def\BC{{\mathbb C}}
\def\Z{{\mathbb Z}}
%\theoremstyle{thmrm} % Numbered and Roman
\newtheorem{exa}[thm]{Example}
% Roman and not numbered
%tava assim mudei \newtheorem{rem}*{Remark}
\newtheorem{rem}[thm]{Remark}
\newtheorem*{oldproof}{Proof}
%\renewenvironment{proof}[1][{}]{\begin{oldproof}[#1]}{\qed\end{oldproof}}
%\usepackage{ulem,xpatch}
\usepackage{amssymb}
\usepackage{amsmath}
\usepackage{amsthm}
\usepackage{mathrsfs}
\usepackage[all]{xy}
\usepackage{graphicx,color}
%\def\Z{{\mathbb Z}}
\usepackage{authblk}

\title{Profinite genus of HNN-extensions with finite associated subgroups}
%\author{Vagner R. de Bessa, \ Anderson L. P. Porto and \ Pavel A. Zalesskii}
%\date{July 27, 2012}
%\author{de Bessa, V. R., Porto, A. L. P., Zalesskii, P. A.  } %\address{Federal University of Viçosa}
%\author{Porto, A. L. P.}
%\address{First line of address,  Second line of address,  etc.}
%\author{Zalesskii, P. A.}
%\address{First line of address\\Second line of address\\etc.}


\author[$\dagger$]{V.R. de Bessa} \author[$\star$]{A.L.P. Porto} \author[$\ddag$]{P.A. Zalesskii\footnote{The third author was partially supported by CNPq. The second author is grateful for the financial support from FAPDF and the Post-Doctoral internship period at UFMG. }}

\affil[$\dagger$]{Federal University of Viçosa, Rio Paranaíba, Brazil}
\affil[$\star$]{ICT, Federal University of the Jequitinhonha and Mucuri Valleys, Diamantina, Brazil} 
\affil[$\ddag$]{Department of Mathematics, University of Brasilia, Brasilia, Brazil}


%\keywords{Genus for groups. Profinite groups. Profinite completions. }
%\classification{20E18 profinite, 20E08 grupos e arvores, hnn e algmama 20E06, 20E34 estrutura geral}


\begin{document}

\maketitle


\begin{abstract} %A finitely generated residually finite group is an $\Oe$-group if when it acts on a profinite tree with finite edge stabilizers it has a global fixed point.
%A finitely generated residually finite group $G$ is an $\Oe$-group if any action of its profinite completion $\widehat G$  on a profinite tree with finite edge stabilizers admits a global fixed point. 

We study the profinite genus of HNN-extensions whose associated subgroups are finite. We give precise formulas for the number of isomorphism classes of $\hnn(G,H,K,t,f)$ and of its profinite completion and compute the profinite genus of such an HNN-extension $\hnn(G,H,K,t,f)$. We also list various situations when   $\hnn(G,H,K,t,f)$ is determined by its profinite completion. 
\end{abstract}


%\begin{abstract} In this paper, we study the profinite genus of  free products with finite amalgamation of finitely generated  groups. We find also sufficient conditions for such amalgamated products to have only one element in its genus.
%\end{abstract}


\section{Introduction} \label{intro}


There has been much recent study of whether residually finite groups, or classes of residually finite groups of combinatorial nature may be distinguished from each other by their sets of finite quotient groups.

In group theory the study in this direction started in 70-th of the last century when Baumslag \cite{BAU74},  Stebe \cite{Ste72} and others found examples of finitely generated residually finite groups having the same set of finite quotients. The general question  addressed in this study can be formulated as follows: 
		
		\begin{question}\label{question}
			To what extent a finitely generated residually finite group $\Gamma$ is  determined by its finite quotients?
		\end{question} 
		
		The study leaded to the notion of genus $\mathfrak{g}(G)$ of a residually finite group $G$, the set of isomorphism classes of finitely generated residually finite groups having the same set of finite quotients as $G$. Equivalently, $\mathfrak{g}(G)$ is the set of isomorphism classes of finitely generated residually finite groups having the  profinite completion isomorphic to the profinite completion $\widehat G$ of $G$. 
		
		This study mostly was concentrated to establish whether the cardinality of  the  genus $\mathfrak{g}(G)$ is finite or  1 (see  \cite{Aka, BCR16, BMRS18, BMRS20, BPZ2, BPZ, BZ, Dav, GPS, GS, GZ, J, J2, J3, M, Sam, Sam2, Wil17} for example) (we use the same term {\it genus} for $g(G)$ from now on). There are only few papers where exact numbers or estimates of the genus appear due to the difficulties of such calculation (see \cite{BPZ2, BPZ, BZ, GZ, Ner19, Ner20, Ner24}).
However, the following principle question of Remeslennikov is still open \cite[Question 5.48]{K}.
		
		\begin{question} (V.N. Remeslennikov)
		Is the genus $g(F)$ of a free group $F$ of finite rank equals 1?
		
		\end{question}

There are two basic constructions in combinatorial group theory: free product with amalgamation and HNN-extension. As many groups of combinatorial and geometric nature split as an amalgamted free product or HNN-extension, it has sense  to study the genus of these constructions.  It is natural to start with splitting over finite groups, and for free products with amalgamation this was done in \cite{BPZ}.  However, since Remeslennikov's question is not answered, it is not clear whether the profinite completion of an one-ended group does not split as an amalgamated free product or an HNN-extension over a finite group. This naturally gives restriction to the family in which one does considerations.



A finitely generated residually finite group $G$ with less than 2 ends, i.e. a finite or one ended group, will be called an $\OE$-group in the paper.   It follows from the famous Stallings theorem that $G$ is an $\OE$-group  if and only if whenever it acts on a tree with finite edge stabilizers
it  has a global fixed point. A finitely generated profinite group will be called $\OE$-group if it has the same property:   whenever it acts   on a  profinite tree  with finite edge stabilizers, then it fixes a vertex.  Note that  finite and one-ended groups acting on a tree have a global fixed point and we need to extend this property to the profinite completion.  A finitely generated residually finite group $G$ will be called $\Oe$-group if  $\widehat G$  is a profinite $\OE$-group, i.e. $\widehat G$  can not act on a profinite tree with finite edge stabilizers without a global fixed point (see Definition \ref{Profinite tree} for the definition of a profinite tree). 

Note that the class $\Oe$-groups is quite large. It contains all finitely generated residually finite not infinite virtually cyclic soluble groups, virtually surface  groups, indecomposable (into free products) 3-manifold groups,  indecomposable RAAGs, as well as all arithmetic groups of rank $>1$ and finitely generated residually finite Fab groups (see Propositions \ref{virtually cyclic} and \ref{examples}). 

 

The subject of the present paper is to complement the work that was done in \cite{BPZ2} and  \cite{BPZ} by studying the profinite genus of HNN-extensions of an $\Oe$-group whose associated subgroups are finite. Note that 
in \cite{BZ} examples of HNN-extensions with finite base group whose profinite completions are isomorphic were given and in this paper we give more examples with not obligatory finite base group. 



We first investigate the isomorphism problems for (profinite) abstract HNN-extensions that are of independent interest and in fact calculate the number of isomorpism classes of such abstract HNN-extensions (see Sections \ref{iso abstract HNN} and \ref{profinitehnnextensions}). 

Let $G=\hnn(G_1, H, K, f, t)$ be an abstract HNN-extension with finite associated subgroups $H$ and $K$.  Denote by $S$ the set of all finite subgroups of $G_1$. There is a natural action of $G_1\rtimes \aut(G_1)$ on the set $S$ namely \begin{equation} (g_1,\alpha)\cdot X= g_1\alpha(X)g_1^{-1} \in S, \mbox{ where } (g_1,\alpha) \in G_1\rtimes \aut(G_1), X \in S. \end{equation}



Denote by  $\aut_{G_1}(H)$  the set of automorphisms of $G_1$ that leaves $H$ invariant. %and by $\widetilde{\aut}_{G_1}(H)$ the image of it in the group of outer automorphisms $\out(H)$ of $H.$ 
Then the stabilizer of the conjugacy class of $H$ in $G_1\rtimes \aut(G_1)$ is the subgroup
\begin{equation} \Gamma_H:=G_1 \rtimes \aut_{G_1}(H)\end{equation}  of the holomorph $G_1\rtimes \aut(G_1)$. 
Denote by $\Gamma_{HK}$  the stabilizer of $K\in S$ in $\Gamma_H$. 
 
 Let $\Iso(H, K)$ be the set of all isomorphisms of $H$ in $K.$ Then $\Gamma_{HK}$ acts naturally on the set of isomorphisms $\Iso(H,K)$  whose action is given by $$(g_1,\alpha) \cdot f=\tau_{g_1} \alpha f \alpha^{-1},$$ 
 where $(g_1,\alpha) \in \Gamma_{HK}$,  $f \in \Iso(H, K)$ and $\tau_{g_1}$ it the inner automorphism corresponding to $g_1$.   
 
 
 We denote by $\widetilde\Gamma_{HK}$ the image of $\Gamma_{HK}$ in the symmetric group $Sym(\Iso(H,K)).$ Depending on the choice  of $H$ and $K$  it might exist $\psi\in \aut(G_1)$  that swaps the conjugacy classes of $H$ and $K$, say $\psi(H)=K$ and $\psi(K)=H^{g_1}$ for some $g_1\in G_1$. In this case the "invertion" of elements of $\Iso(H,K)$ followed by the action by the element $( g', \psi)\in G_1\rtimes \aut(G_1)$ gives an element $\iota\in Sym(\Iso(H,K))$ that normalizes $\widetilde\Gamma_{HK}$ and so $\overline\Gamma_{HK}:=\widetilde\Gamma_{HK} \cdot \langle\iota\rangle$ is a group that does not depend in fact upon a choice of $\psi$ (see Lemma \ref{carmo}). If such $\psi$ does not exist we put $\overline\Gamma_{HK}:=\widetilde\Gamma_{HK}$. 


\begin{thm} Let $G_1$ be an $\OE$-group and $$G=\hnn(G_1, H, K, f, t), B=\hnn(G_1, H, K, f_1, t_1)$$ be  abstract HNN-extensions with fixed finite associated  subgroups $H,K$ of $G_1$. Then $G$ and $B$ are isomorphic if and only if $f$ and $f_1$ belong to the same $\overline\Gamma_{HK}$-orbit in $\Iso(H,K)$.    
\end{thm}

\begin{cor}\label{general iso} The number of isomorphism classes of HNN-extensions \break
$\hnn(G_1,H,K,f,t)$  with $G_1,H,K$ fixed is 

$$|\overline{\Gamma}_{HK}\backslash \Iso(H,K)|.$$

\end{cor}

\bigskip
If the associated subgroups $H$ and $K$ are conjugate in $G_1$, we can make then more explicit the formular of Corollary \ref{general iso}. 
Denote by $\widetilde{\aut}_{G_1}(H)$ and by $\widetilde{N}_{G}(H)$ the natural images of $\aut_{G_1}(H)$ and $N_G(H)$ in the group of outer automorphisms $\out(H)$ of $H$.


\begin{cor} Let $H$ and $K$ be conjugate finite subgroups of an $\OE$-group $G_1$. Then the number  of isomorphism classes of HNN-extensions \break $\hnn(G_1, H, K, f,t)$ 
 is $$|\widetilde{N}_{G_1}(H) \backslash \out(H)/  ( \widetilde{\aut}_{G_1}(H) \times C_2) |,$$
 where $C_2$ acts on $\out(H)$ by inversion. 
\end{cor}



Now similarly, denote by $P$ the set of all finite subgroups of $\widehat G_1.$ There is a natural action of $\widehat G_1\rtimes \aut_{\widehat G_1}(H)$ on the set $P$, namely \begin{equation} (g_1,\alpha)\cdot X= g_1\alpha(X)g_1^{-1} \in P, \mbox{ where } X \in P. \end{equation}
 We denote by $\Gamma_{\widehat{HK}}$  the stabilizer of $K\in P$ of this action. 
  Then $\Gamma_{\widehat{HK}}$ also acts naturally on the set of isomorphisms $\Iso(H,K)$. 
  
  We denote by $\widetilde\Gamma_{\widehat{HK}}$ the image of $\Gamma_{\widehat{HK}}$ in the symmetric group $Sym(\Iso(H,K)).$ Depending on the choice   $H$ and $K$  it might exist $\psi\in \aut(\widehat G_1)$  that swaps the conjugacy classes of $H$ and $K$ in $\widehat G_1$, say $\psi(H)=K$ and  $\psi(K)=H^{g_1}$ for some $g_1\in \widehat G_1$. In this case the "invertion" of elements of $\Iso(H,K)$ followed by the action by the element $( g', \psi)\in \widehat G_1\rtimes \aut(\widehat G_1)$ %$( g_1^{-1}, \psi)\in \widehat G_1\rtimes \aut(\widehat G_1)$ 
  gives an element $\iota\in Sym(\Iso(H,K))$ that normalizes $\widetilde\Gamma_{\widehat{HK}}$ 
  and so $\overline\Gamma_{\widehat{HK}}:=\widetilde\Gamma_{\widehat{HK}} \cdot \langle\iota\rangle$ is a group that does not depend in fact upon a choice of $\psi$ (see Lemma \ref{carmoprof}). If such $\psi$ does not exist we put $\overline\Gamma_{\widehat{HK}}:=\widetilde\Gamma_{\widehat{HK}}$.
  

In the profinite case we can only give an estimate for the number of isomorphism classes of profinite HNN-extensions with fixed associated subgroups, as follows.
\begin{thm}  Let $H, K$ be fixed finite subgroups of a profinite $\OE$-group $G_1$. Then the number of isomorphism classes of profinite HNN-extensions \\$\hnn(G_1, H, K, f, t)$ is less than or equal to $$|\overline{\Gamma}_{\widehat{HK}} \backslash \Iso(H, K)|.$$\end{thm}  

%\begin{thm} Let $G_1$ be a profinite $\OE$-group and $G=\hnn(G_1, H, K, f, t)$ and $B=\hnn(G_1, H, K, f_1, t_1)$ be  profinite HNN-extensions with finite associated  subgroups. Then $G$ and $B$ are isomorphic if an only if $f$ and $f_1$ belong to the same $\overline\Gamma_{\widehat{HK}}$-orbit in $\Iso(H,K)$.    \end{thm}}

%Thus the number of isomorphism classes of profinite HNN-extensions $\hnn(G_1,H,K,f,t)$ with $G_1,H,K$ fixed is $$|\overline\Gamma_{\widehat{HK}}\backslash \Iso(H,K)|.$$

If the associated subgroups are conjugate in $\widehat G_1$  then in contrast to abstract case we get only an upper bound in terms of $H$.

\begin{thm} Let $H$ and $K$ be fixed finite subgroups that are conjugate in the profinite $\OE$-group $G_1$. Then the number of the isomorphism classes of profinite HNN-extensions $G=\hnn(G_1, H, K, f, t)$ is less or equal than $$\left| \widetilde N_{G_1}(H) \backslash \out(H)/ \left( \widetilde{\aut}_{G_1}(H) \times C_2\right) \right|,$$ where $C_2$ acts on $\out(H)$ by inversion. 
\end{thm}

After that, we study the genus of HNN-extensions where the base group $G_1$ is an $\rm \Oe$-group and the associated subgroups are finite. 

%Let $H \leq G,$ denote by $N_{G}(H)$ the normalizer of $H$ in $G,$ and $\overline{N}_{G}(H)$ and $\widetilde{N}_{G}(H)$ its natural images in $\aut(H)$ and $\out(H)$, respectively.

\medskip
Let $\mathcal{F}=\mathcal F(G_1,H,K)$ be the class of all 
abstract HNN-extensions $$G=\hnn(G_1, H, K, f, t),$$ where $H$ and $K$ are fixed finite subgroups of $G_1$. Denote by $N^+$  the following subset  
  \begin{equation*} % \label{N mais com completamento}
N^+=\left\{ gt^{-1} \in N_{\widehat G}(H) \hspace{0.1cm} \left| \hspace{0.1cm}  \overline{\left\langle g, \widehat G_1 \right\rangle}=\widehat G \right. \right\}.   
\end{equation*} 
Let $\overline{N}^{+}$ and $\widetilde{N}^{+}$ be the natural images of $N^{+}$ in $\aut(H)$ and $\out(H)$, respectively. Given $G\in \F$ we first concentrate on calculation of the cardinality of genus $g(G,\F)$ (denoted by $|g(G, \F)|$)  within this family of HNN-extensions. 


\begin{thm} Let $H$ and $K$ be finite subgroups of an $\Oe$-group $G_1$. Then the cardinality of the genus $g(G, \mathcal{F})$ of an HNN-extension $G=\hnn(G_1,H,K,f,t)$ is equal to $$|g(G, \mathcal{F})| = \left|\overline\Gamma_{HK}  \backslash \overline\Gamma_{\widehat{HK}}  \cdot  f \overline{N}^+\right|,$$ where $\overline\Gamma_{HK}  \backslash \overline\Gamma_{\widehat{HK}}  \cdot f\overline{N}^+$ represents the $\overline\Gamma_{HK}$-orbits in the set of   $\overline\Gamma_{\widehat{HK}}$-orbits of twists $f\tau_{n^{-1}}\,(n \in N^{+})$ of $f$. 
\end{thm}


%$\overline{\aut}_{G_1}(H), \aut_{\widehat{G}_1}(H),$

An HNN-extension will be called normal if the stable letter normalizes the associated subgroup. Note that if $G$ is an HNN-extension with conjugate associated finite subgroups  in the base group then it is isomorphic to a normal HNN-extension (see Lemma \ref{conjugados da normal}), so in this case we can consider $f \in \aut(H)$. Let $\tilde f$ be the image of $f$ in $\out(H).$ It follows that $\tilde{f}\widetilde{N}^+$ is a subset of $\tilde f \widetilde{N}_{\widehat G}(H)$.

In this case  $\widetilde{\aut}_{\widehat G_1}(H)$ acts on $\out(H)$ on the right by conjugation. The group $\widetilde{N}_{\widehat G_1}(H)$ acts on $\out(H)$ on the left by multiplication. It follows that $\widetilde{N}_{\widehat G_1}(H) \cdot \tilde{f} \widetilde{N}^+ \cdot \widetilde{\aut}_{\widehat G_1}(H)$ is invariant with respect to the action of $\widetilde{N}_{G_1}(H)$ on the left and with respect to the action of $\widetilde{\aut}_{G_1}(H)$ on the right.







\begin{thm}  Let $G_1$ be an $\Oe$-group and $H$ a fixed finite subgroup of $G_1$. Let $G=\hnn(G_1,H, f,t)$ be a normal HNN-extension.  Then $$|g(G, \mathcal{F})| = | \widetilde{N}_{G_1}(H) \backslash \widetilde{N}_{\widehat G_1}(H) \cdot \tilde{f}\widetilde{N}^+ \cdot \widetilde{\aut}_{\widehat G_1}(H)/( \widetilde{\aut}_{G_1}(H) \times C_2)|,$$ where $C_2$ acts on $\out(H)$ by inversion. 
\end{thm}





%Let $\mathcal{F}(G_1,\A)$ be the class of all  abstract HNN-extensions of the form $\hnn (G_1, H_1, K_1, f_1, t_1)$, where $H_1$ and $K_1$ are finite subgroups of $G_1,$ $f_1 \in \Iso(H_1, K_1)$ and $H \cong H_1 \cong K_1 \cong K.$ In this case we give an exact formula for the profinite genus, however we leave the details for reading in Subsection \ref{non fixed}. 


%\begin{thm} \label{genus hnn nao fixos into} Let $G_1$ be an $\Oe$-group and $G=\hnn(G_1, H, K, f, t)$ be an HNN-extension with finite associated subgroups. Then the cardinality of genus $g(G,\mathcal{NF})$ of $G$  equals to $$|g(G,\mathcal{NF})|=\displaystyle\sum_{\{H_1,K_1\}} \left| g\left(\hnn(G_1, H_1, K_1, f_1, t_1), \F\right)\right|,$$ where $\{H_1,K_1\}$  ranges over unordered pairs of representatives of  $\mathscr{L}_{_{HK}}\\ (\mbox{see Remark } \ref{LHK}).$  
%\end{thm}

%Note that if $G_1$ has only finitely many conjugacy classes of finite subgroups, then $|\mathscr{L}_{_{HK}}|$ is finite. Arithmetic groups   \cite{borelarith} and many groups of geometric nature as well as  mapping class groups of any surface of finite type \cite[Theorem 6]{Brid}, the automorphism group and the outer automorphism group of a free group of finite rank (see \cite{culler}, \cite{zimm} and p. 19 in \cite{vog}),  hyperbolic groups   \cite[Theorem 3.2, p. 459]{bridcurvature},  
%groups that act properly and cocompactly by isometries on a $CAT(0)$ space \cite[Corollary 2.8 (2), p. 179]{bridcurvature}, share this property.




%\begin{thm} Suppose $G_1$ is a profinitely rigid $\Oe$-group. Then the statement of Theorem \ref{genus hnn nao fixos into} is valid for $g(G,\mathcal{B})$, where $\mathcal{B}$ is the family of all finitely generated residually
%finite not one-ended groups.
%\end{thm}

%In fact, if one-endedness is a profinite property then the formulas would be valid for absolute genus (i.e. within all finitely generated residually finite groups). 

 

We also establish a list of situations in which the genus formula is simpler (see Subsection \ref{section 5 1}), in particular, we find conditions for the genus cardinality to be  1 ( generalizing the Theorem 1.1 in \cite{BZ}). 

\begin{deff} A pair $(H,K)$ of subgroups of $G$ will be called a swapping pair if there exists an automorphism of $G$ that swaps the conjugacy classes of $H$ and $K$.  \end{deff}

\begin{thm}\label{genus 1 intr}
Let $H$ and $K$ be a finite subgroups of an $\Oe$-group $G_1$ and $G=\hnn(G_1, H, K, f, t)$ be an HNN-extension, where $f \in \Iso(H, K).$ 
Suppose the following conditions are satisfied \begin{itemize}
    \item $\widetilde{\aut}_{\widehat{G}_1}(H)=\widetilde{\aut}_{G_1}(H)$ and $[G_1:N_{G_1}(K)]=[\widehat G_1:N_{\widehat G_1}(K)]<\infty$,
%    \item $[G_1:N_{G_1}(K)]=[\widehat G_1:N_{\widehat G_1}(K)]<\infty$,
    \item  $(H,K)$ is either a swapping pair in $G_1$ or not a swapping pair in $\widehat G_1$. 

\end{itemize} Then there is only one element in the genus of $G$ in $\F,$ if one of the following conditions is satisfied:

\begin{itemize}

%\item[i)] $H$ and $K$ are conjugate in $G_1$;

\item[i)]  $\out(H)$  coincides with the image of the normalizer $N_{G_1}(H)$ in $\out(H)$;

\item[ii)] %$H$ and $K$ are conjugate in $G_1$ and $\tau f\in \overline{N}_{G_1}(H)$, for some $\tau\in \inn(G_1)$;} 
$H^{g_1}=K$ for some $g_1 \in G_1$ and $\tau_{g_1} f\in \overline{N}_{G_1}(H);$

%\item [iii)] $H \neq K$ and either $H$ or $K$ is central in $G_1;$}  
    \item [iii)] $H$ and $K$ are not conjugate in $G_1$ and $f$ extends to an automorphism of $G_1;$  
    %% $f \in \aut(G_1);$
    \item [iv)] $H$ and $K$ are not conjugate in $G_1$ and $N_{G_1}(K)=K \cdot C_{G_1}(K)$ or $N_{G_1}(H)=H \cdot C_{G_1}(H)$ (in particular, if $H \neq K$ and $H$ or $K$ is central in $G_1$).  
\end{itemize}
\end{thm}




The structure of the  paper is as follows.  Section \ref{preliminares} contains elements of the profinite version of the Bass-Serre theory used in the paper (see \cite{R} for more details) and some results on normalizers of profinite and  abstract HNN-extensions.  Futhermore it posseses  examples and properties of $\Oe$-groups. Sections \ref{iso abstract HNN} and \ref{profinitehnnextensions} dedicated to the isomorphism problem of (profinite) abstract HNN-extensions which are of independent interest. Moreover, we calculate the number of isomorphism classes of abstract and profinite HNN-extensions $G=\hnn(G_1, H, K, f, t)$ with not fixed associted subgroups $H,K$. More precisely,  the action of $\aut(G_1)$ on the set $S$ of finite subgroups of $G_1$ induces the action on the set of unordered pairs $[\overline{S}]^2$ of conjugacy classes of finite isomorphic subgroups of $G_1.$ 
The next theorem then gives a formula (resp. estimate) for the number of isomorphism classes of all abstract (resp. profinite) HNN-extensions of $G_1$ with finite associated subgroups.



\begin{thm} \label{int general finite associated} Let $G_1$ be  an $\OE$-group (resp. profinite $\OE$-group). Then the number  of isomorphism classes $n_a$ (resp. $n_p$) of  abstract (resp. profinite) HNN-extensions $\hnn(G_1, H_,K, f,t)$ with finite associated subgroups  is  $$n_a=\sum_{\{H, K\}} |\overline{\Gamma}_{HK} \backslash \Iso(H, K)|,$$
$$({\rm resp.}\  n_p\leq\sum_{\{H, K\}} |\overline{\Gamma}_{\widehat{HK}} \backslash \Iso(H, K)|) $$ where $\{H, K\}$ ranges over representatives of the $\aut(G_1)$-orbits of the set of unordered pairs $[\overline{S}]^2$ of conjugacy classes of finite isomorphic subgroups of $G_1$. 
\end{thm} 

The formula for such abstract (resp. profinite) HNN-extensions where associated subgroups are conjugate (equivalently coincide) is as follows:

\begin{thm}\label{int conjugate associated} Let $G_1$ be an $\OE$-group (resp. profinite $\OE$-group). Then the number of isomorphism classes of abstract (resp. profinite) HNN-extensions $\hnn(G_1, H, K, f, t)$ with finite conjugate associated subgroups  equals (resp. less or equal) $$\sum_{H} |\widetilde{N}_{G_1}(H) \backslash \out(H)/  ( \widetilde{\aut}_{G_1}(H) \times C_2) |,$$
 where $C_2$ acts on $\out(H)$ by inversion and $H$ ranges over representatives of the $\aut(G_1)$-orbits of finite subgroups of $G_1$.
\end{thm} 

In Section \ref{secao 5} we calculate the genus or estimate it for HNN-extensions. Sections 6 is dedicated to the calculation of genus and bounds in more general classes of groups.  Namely, Theorems \ref{int general finite associated} and \ref{int conjugate associated} combined with  \cite[Theorem 1.1]{BPZ}  can be used to compute  the genus of HNN-extension within a larger class $\A$ of finitely generated residually finite accessible groups whose vertex groups of its JSJ-decomposition are $\Oe$-groups  (see Theorem \ref{fixed H2} in Subsection  \ref{general section}).  



%Moreover,  if $G_1, G_2$ are profinitely rigid we can considerably widen the family of our consideration

%\bigskip
%Our basic reference for notations and results about profinite groups is \cite{RZ}. We refer the reader to Lyndon-Schupp \cite{LS}, Magnus-Karrass-Solitar \cite{MKS}, Serre \cite{S} or Dicks-Dunwoody \cite{DI} for an account of basic facts on amalgamated free products  and to \cite{RZ} for the profinite versions of these constructions. Our methods based on the profinite version of the Bass-Serre theory of groups acting on trees that can be found in \cite{R}. All homomorphisms of profinite groups are assumed to be continuous in this paper. We will use the standard abbreviation $x^{g}$ for $g^{-1}xg$ when $g,x$ are elements of a group $G;$ the inner automorphism of $G$ corresponding to the conjugation $gxg^{-1}$ will be denoted by $\tau_g$ (as all maps are written on the left). If $H$ is a subgroup of $G,$ we shall use $H^G$ for the  the normal closure of $H$ in $G$.

%The composition of  two maps $f$ and $g$ is often defined simply as $f\circ g = fg.$ For a subgroup $H$ of $G$ we shall denote by $\aut_G(H)$ the group of automorphisms of $G$ that leave $H$ invariant and by $\overline{\aut}_G(H),$ $\widetilde{\aut}_G(H)$ its images in $\aut(H)$ and $\out(H),$ respectively. %All HNN-extensions  $G=\hnn(G_1, H, K, f, t)$ (resp. profinite HNN-extensions) 
%will be assumed non-fictitious in the paper, i.e. $G_1\cap \left\langle t\right \rangle = \{1\}$ and 
%$H \neq G_1$. 

%All abstract HNN-extensions are residually finite and in the profinite case, they will always be proper (i.e. the base group  embeds in the HNN-extension);  by Proposition 9.4.3 ii)-iii) in \cite{RZ} this is always the case if  $H, K$ are finite.


{\subsection{Notations and terminologies} 

Our basic reference for notations and results about profinite groups is \cite{RZ}. We refer the reader to Lyndon-Schupp \cite{LS}, Magnus-Karrass-Solitar \cite{MKS}, Serre \cite{S} or Dicks-Dunwoody \cite{DI} for an account of basic facts on amalgamated free products  and to \cite{RZ} for the profinite versions of these constructions. Our methods based on the profinite version of the Bass-Serre theory of groups acting on trees that can be found in \cite{R}. 

We  give a brief description of the main notations and terminologies used throughout the paper. 

All homomorphisms of profinite groups are assumed to be continuous in this paper. The composition of  two maps $f$ and $g$ is often defined simply as $f\circ g = fg.$ We will use the standard abbreviation $x^{g}$ for $g^{-1}xg$ when $g,x$ are elements of a group $G;$ the internal automorphisms of $G,$ denoted by $\inn(G),$ is the group composed of the elements $\tau_g\, (g \in G)$ (as all maps are written on the left) such that $\tau_g(x)=gxg^{-1}, \forall x \in G.$ 
 
If $H$ is a subgroup of a group $G$, we will use $H^G$ for the normal closure of $H$ in $G,$  $N_G(H)$ for the normalizer of $H$ in $G$ and $C_G(H)$ for the centralizer of $H$ in $G.$  Furthermore, let $\aut_G(H)$ be the group of all automorphisms of $G$ that leave $H$ invariant and denote by $\overline{\aut}_G(H),$ $\widetilde{\aut}_G(H)$ its images in $\aut(H)$ and $\out(H),$ respectively. The natural image of $x \in N_G(H)$ in $\aut(H)$ will be given by $\tau_{x^{-1}},$ and let $\overline{N}_G(H)$ and $\widetilde{N}_G(H)$ denote the natural images of $N_G(H)$ in $\aut(H)$ and $\out(H),$ respectively.

 All abstract HNN-extensions are residually finite and in the profinite case, they will always be proper (i.e. the base group  embeds in the HNN-extension);  by Proposition 9.4.3 ii)-iii) in \cite{RZ} this is always the case if  $H, K$ are finite.

Throughout the paper we shall us the following families of groups. 
\begin{itemize}

\item $\mathcal{F}=\mathcal F(G_1,H,K)$ - be the class of all 
abstract HNN-extensions $$\hnn(G_1, H, K, f, t),$$ where $H$ and $K$ are fixed finite subgroups of $\Oe$-group $G_1$. 

\item $\A$ - the family of  finitely generated residually finite accessible groups whose vertex groups of its JSJ-decomposition are $\Oe$-groups. 

\item $\mathcal F(G, G_1, \A)$ - be the subclass of $\mathcal{A}$ consisting of all abstract HNN-extensions of the form $$B=\hnn(G_1, H_1, K_1, f_1, t_1),$$ where $G_1$ is an $\Oe$-group fixed  in the class $\mathcal{A}$ and the associated subgroups are finite.

\end{itemize}


\section{\large Preliminary Results} \label{preliminares}

In this section we will provide the basic and preliminary facts for the development of the main sections. 

\subsection{Basics of profinite trees} \label{basics of trees}

In this subsection we recall the necessary notions of the Bass-Serre theory for abstract and profinite graphs.
	
	\begin{deff}[Profinite graph]
		A (profinite) graph is a (profinite space) set $\Gamma$ with a distinguished closed nonempty subset $V(\Gamma)$ called the vertex set, $E(\Gamma)=\Gamma-V(\Gamma)$ the edge set and two (continuous) maps $d_0,d_1:\Gamma \rightarrow V(\Gamma)$ whose restrictions to $V(\Gamma)$ are the identity map $id_{V(\Gamma)}$. We refer to $d_0$ and $d_1$ as the incidence maps of the (profinite) graph $\Gamma$.  
	\end{deff}
	
	A morphism $\alpha:\Gamma \longrightarrow \Delta$ of profinite graphs is a continuous map with $\alpha d_i=d_i \alpha$ for $i=0,1$. By \cite[Proposition 2.1.4]{R} every profinite graph $\Gamma$ is an inverse limit of finite quotient graphs of $\Gamma$. A profinite graph $\Gamma$ is called connected if every its finite quotient graph is connected as a usual graph. 
	
	\begin{deff}\label{Profinite tree} Let $\Gamma$ be a profinite graph. Define $E^{*}(\Gamma)=\Gamma/V(\Gamma)$ to be the quotient space of $\Gamma$ (viewed as a profinite space) modulo the subspace of vertices $V(\Gamma)$. Consider the free profinite $\widehat\Z$-modules $[[\widehat\Z(E^{*}(\Gamma),*)]]$ and $[[\widehat\Z V(\Gamma)]]$ on the pointed profinite space $(E^{*}(\Gamma),*)$ and on the profinite space $V(\Gamma)$, respectively. Denote by $C(\Gamma,\widehat\Z)$ the chain complex 
	
	
	$$	\xymatrix{ 0 \ar[r] & [[\widehat\Z(E^{*}(\Gamma),*)]] \ar[r]^d &[[\widehat\Z V(\Gamma)]] \ar[r]^{\varepsilon} & \widehat\Z \ar[r] & 0}$$ of free profinite $\widehat\Z$-modules and continuous $\widehat\Z$-homomorphisms $d$ and $\varepsilon$ determined by $\varepsilon(v)=1$, for every $v \in V(\Gamma)$, $d(\overline{e})=d_1(e)-d_0(e)$, where $\overline{e}$ is the image of an edge $e \in E(\Gamma)$ in the quotient space $E^{*}(\Gamma)$, and $d(*)=0$. 
		One says that $\Gamma$ is a profinite tree if the sequence $C(\Gamma,\widehat\Z)$ is exact. 
	\end{deff}	
	
	  If $v$ and $w$ are elements of a tree (respectively profinite tree) $T$, one denotes by $[v,w]$ the smallest subtree (respectively profinite subtree) of $T$ containing $v$ and $w$.
	

%\begin{deff}
%Let $\Gamma$ be a connected finite graph. A    graph of profinite groups $({\cal G},\Gamma)$ over
%$\Gamma$ consists of  specifying a profinite group ${\cal G}(m)$ for each $m\in \Gamma$, and continuous monomorphisms
%$\partial_i: {\cal G}(e)\longrightarrow {\cal G}(d_i(e))$ for each edge
%$e\in E(\Gamma)$, $i=0,1$. We say that it is reduced if $\G(e)\neq \G(v_i)$ for all edges of $\Gamma$ which are not loops.
%\end{deff}

%In \cite[paragraph (3.3)]{ZM1},  the fundamental group
 %$G=\Pi_1(\G,\Gamma)$ is  defined explicitly in terms of generators and relations
 %associated to a chosen maximal subtree $D$. Namely the profinite presentation is as follows:
 
% \begin{eqnarray} \label{presentation} 
 % G=\Big\langle
 %\G(v), t_e, v\in V(\Gamma), e\in E(\Gamma) \mid &  \nonumber\\ t_e=1, \ {\rm for}\  e\in D, \partial_0(g)=t_e\partial_1(g)t_e^{-1}, \ {\rm for}\ g\in \G(e)\Big\rangle & 
%\end{eqnarray}
%I.e., if one takes the abstract fundamental group $G_0=\pi_1(\G,\Gamma)$,
%then $$\Pi_1(\G,\Gamma)=\varprojlim_N G_0/N,$$ where $N$ ranges over
%all normal subgroups of $G_0$  with $N\cap
%\G(v)$ open in $\G(v)$ for all $v\in V(\Gamma)$. Note that this last
%condition is automatic if $\G(v)$ is finitely generated (as a
%profinite group, see \cite{NS-07}). 
 %  It is also proved in \cite{ZM1}
%that the definition given above is independent of the choice of
%the maximal subtree $D$.

%If all vertex groups are trivial we get the definition of the profinite fundamental group $\pi_1(\Gamma)$ that is the profinite completion of the usual fundamental group.

%Associated with the profinite graph of profinite %groups $({\cal G}, \Gamma)$ there is
%a corresponding   standard profinite tree  (or %universal covering graph)
 % $$S=S(G)=\bigcup_{m\in \Gamma}
%G/\G(m)$$ (cf. \cite[Theorem 3.8]{ZM1}).  The vertices of
%$S$ are those cosets of the form
%$g\G(v)$, with $v\in V(\Gamma)$
%and $g\in G$; its edges are the cosets of the form %$g\G(e)$, with $e\in
%E(\Gamma)$; and the incidence maps of $S$ are given by the formulas: $$d_0 (g\G(e))= g\G(d_0(e)); \quad  d_1(g\G(e))=gt\G(d_1(e)) \ \ 
%(e\in E(\Gamma), t_e=1\hbox{ if }e\in D), $$
%where $D$ is a maximal subtree of 
%$\Gamma$. 

 %There is a natural  continuous action of
% $G$ on $S$, and clearly $ G\backslash S= \Gamma$. 
%%nova
\medskip 

%%%%%%%%%%%%%%%%%%%%%%%%%%%%%%%%%%
%If $G=\hnn(G_1,H,K,f,t)$ is an HNN-extension then $G=\Pi_1(\G, \Gamma)$ with $\Gamma$ having one vertice and one edge, so that $G_1$ is the vertex group and $H$ being an edge group. 

Let $G=\hnn(G_1, H, K, f, t)$ be a profinite HNN-extension. In contrast to the abstract case, $G_1$ do not always embed into $G$, i.e. a  profinite HNN-extension is not always proper. Note that $G$ is always proper  when $H$ and $K$ are finite (see Propositon $9.4.3$ item $(2)$ in \cite{RZ}) which will be assumed for starting from Subsection \ref{general results}.  In any case we shall consider only proper free profinite product with amalgamation and profinite HNN-extensions in this paper.

When $G=\hnn(G_1,H,K,f,t)$ is proper, there is
a corresponding    standard profinite tree  (or universal covering graph) 
  (cf. \cite[Theorem 3.8]{ZM1}) $S(G)=G/G_1\cup G/H$, $V(S(G))=G/G_1$ and  $d_0 (gH)= gG_1; \quad  d_1(gH)=gtG_1 $. There is a natural  continuous action of
 $G$ on $S(G)$, and clearly $ G\backslash S(G)$ is an edge with one vertex. 

%\medskip 
 %If $G=G_1\amalg_HG_2$ is an amalgamated free profinite product then $G=\Pi_1(\G, \Gamma)$ with $\Gamma$ having two vertices and one edge, $G_1,G_2$ being vertex groups and $H$ being an edge group. In this case $S(G)=G/G_1\cup G/G_2\cup G/H$ and  $d_0 (g\G(e))= g\G(d_0(e)); \quad  d_1(g\G(e))=g\G(d_1(e)) \ \ 
%(e\in E(\Gamma)).$

%\bigskip
%\bigskip
%We shall need slightly more general version of \cite[Lemma 4.4]{GZ}. 

%\begin{lem}\label{generation} Let $G=G_1*_HG_2$ be an  amalgamated free product and let $g\in G$. Then $G=\langle G_1, G_i^g\rangle$ if and only if $i=2$ and $g\in G_2G_1$.

%\end{lem}  

 %The proof of the lemma is the same as the proof of  \cite[Lemma 4.4]{GZ}  and  will be ommited.

%Note that in the profinite world the factors $G_1,G_2$ do not always embed in the free product with amalgamation, however it is always the case when $H$ is finite (see Exercise $9.2.7$ item $(3)$ in \cite{RZ}) which will be assumed for the rest of the paper.

\subsection{Normalizers}\label{normalizers}

The  description of   normalizers in this subsection is of independent interest for profinite groups.

\begin{prop} \label{normalizer}
\begin{enumerate}
\item[(i)] Let $G=G_1\amalg_H G_2$ be a proper  free profinite product with  amalgamation. Then $N_G(H)=N_{G_1}(H)\amalg_H N_{G_2}(H)$.
\item[(ii)] Let $G=\hnn(G_1, H, K, f, t)$ be a proper profinite HNN-extension with associated subgroups $H,$ $K$ and stable letter $t$. Then  
\begin{itemize}

\item[(a)]  If $H^{g_1}=K$ for some $g_1 \in G_1$ then $N_{G}(H)$ is a profinite HNN-extension  $N_{G}(H)=\hnn(N_{G_1}(H), H, \tau_{g_1} f)$. 

\item[(b)] If $H$ and $K$ are not conjugate subgroups in $G_1$ then $N_{G}(H)$ is a profinite amalgamated free product  $N_{G}(H)=N_{G_1}(H)\amalg_{H} N_{G_1^{t^{-1}}}(H)$. 
\end{itemize}

\end{enumerate}
\end{prop}

\begin{proof} %(i) Follows from Proposition 2.4 in \cite{BPZ2}.\\
Let $S(G)$ be the standard profinite tree on which $G$ acts continuously. Denote by $S(G)^{H}$ the subset of $S(G)$ of  fixed points with respect to the action of $H$. Then $S(G)^{H}$ is a profinite tree (Theorem 2.8 in \cite{ZM1} or Theorem 4.1.5 \cite{R}). Moreover $N_{G}(H)$ acts continuously on $S(G)^{H}$. Indeed, let $s \in S(G)^{H}, g \in N_{G}(H)$ and $h \in H,$ then: $$h\cdot(g\cdot s)=(hg)\cdot s=(gh')\cdot s = g\cdot (h'\cdot s)=g\cdot s,$$ for some $h'\in H,$  whence follows that $g\cdot s \in S(G)^{H},$ where $g\cdot s$ means the action of $g$ on $s$. 

Now the graph $S(G)^{H}/N_{G}(H)$ has only one edge. To see this it suffices to show that given $g\in G$ and $e\in E(S(G)^{H})$, then $ge\in E(S(G)^{H})$ implies  $g\in N_{G}(H)$. But this is obvious because both $H$ and $H^{g^{-1}}$ stabilize $ge$  which implies $H\leq  H^{g^{-1}}.$ However, in a profinite group, a conjugate of a subgroup cannot be contained in it properly; to see this, one need only look at the finite quotients and note that the conjugate subgroups have the same order. % However, in a profinite group a conjugate of a subgroup can not be contained in it properly; to see this it suffices to look ate the finite quotients and use that conjugate sugroups have the same order. 
So $H = H^{g^{-1}}$ as needed.

 Thus $S(G)/G \cong S(G)^{H}/N_{G}(H)$ is an isomorphism of graphs and both are either a loop or an edge with two vertices. Note that the first case occurs when $G$ is an HNN-extension and $H$ and $K$ are conjugate in $G_1$ and the second case occurs when $H$ and $K$ are not conjugated in $G_1$ or $G$ is a free product with amalgamation. Now, note that $S(G)$ is connected and simply connected (see Theorem 3.4 in \cite{ZM} or Theorem 6.3.5 \cite{R}), thus  we can apply Proposition 4.4 in \cite{ZM} or Theorem 6.6.1 together with Proposition 6.5.3 in \cite{R} to deduce that $N_{G}(H)$ has the desired structures  of this proposition.
\end{proof}

In the following proposition gives a description of the normalizer of a  finite associated subgroup $H$ of an HNN-extension in the discrete and profinite categories. 


\begin{prop}\label{cp1} 
Let $G = \hnn(G_{1}, H, K, f, t)$ be an HNN-extensions of a residually finite group $G_1$ with finite associated subgroups $H$ and $K$ and stable letter $t$. Then the following hold: \begin{itemize}

\item[(i)] If $H^{g_1}=K$ for some $g_1 \in G_1$ then $N_{G}(H)=\hnn(N_{G_1}(H), H, \tau_{g_1} f)$. 
If $H$ and $K$ are not conjugate subgroups in $G_1$ then $N_{G}(H)=N_{G_1}(H)*_{H} N_{{G_1}^{t^{-1}}}(H)$.


\item[(ii)] If $H^{g_1}=K$ for some $g_1 \in \widehat{G_1}$ then $N_{\widehat{G}}(H)$ is a profinite HNN-extension (i.e. the profinite completion of the usual HNN-extension) $N_{\widehat{G}}(H)=\widehat{\hnn}(N_{G_1}(H), H, \tau_{g_1} f)=\hnn(N_{\widehat{G_1}}(H), H, \tau_{g_1} f)$. If $H$ and $K$ are not conjugate subgroups in $\widehat{G_1}$ then $N_{\widehat{G}}(H)$ is a profinite amalgamated free product (i.e. the profinite completion of the usual amalgamated free product) $N_{\widehat{G}}(H)=N_{\widehat{G_1}}(H)\amalg_{H} N_{\widehat{G_1}^{t^{-1}}}(H)$. 
\end{itemize}
\end{prop}

\begin{proof} (ii)   Follows from Proposition \ref{normalizer}.  

\medskip
(i) The abstract version of the proof is analogous to the proof of  Proposition \ref{normalizer} and will be omitted;  the abstract version of results quoted in that proof  can be found in \cite{S} or \cite{DI}. 
\end{proof}

%\begin{cor}\label{8} If $G_1$ is a polycyclic-by-finite group, then $N_{G}(H)$ is dense in $N_{\widehat{G}}(H)$. Furthermore, if $H$ is finite, then $$\overline{N}_{\widehat{G}}(H)=\overline{N}_{G}(H)$$ where $\overline{N}_{G}(H)$ and $\overline{N}_{\widehat{G}}(H)$ are the natural images (by restriction) of $N_{G}(H)$ and $N_{\widehat{G}}(H)$ in the automorphism group $\aut(H),$ respectively. 
%\end{cor}

\begin{cor}\label{8} If $G_1$ is a polycyclic-by-finite group, then $N_{G}(H)$ is dense in $N_{\widehat{G}}(H)$. Furthermore,  $$\overline{N}_{\widehat{G}}(H)=\overline{N}_{G}(H)$$ where $\overline{N}_{G}(H)$ and $\overline{N}_{\widehat{G}}(H)$ are the natural images (by restriction) of $N_{G}(H)$ and $N_{\widehat{G}}(H)$ in the automorphism group $\aut(H),$ respectively. 
\end{cor}

\begin{proof} First observe that $H$ and $K$ are conjugate in $G_1$ if and only if they are conjugate in $\widehat G_1$  (see \cite[ Chapter 4, Theorem 7]{Se}).

Suppose that $H$ and $K$ are not conjugate in $\widehat{G_1}$. Then by Proposition \ref{cp1} we have   $$N_{G}(H)=N_{G_1}(H)*_{H} N_{G_{1}^{t^{-1}}}(H) \mbox{ and } N_{\widehat{G}}(H)=N_{\widehat{G_1}}(H)\amalg_{H} N_{\widehat{G_{1}}^{t^{-1}}}(H),$$ and so  $\left\langle N_{\widehat{G_1}}(H), N_{\widehat{G_{1}}^{t^{-1}}}(H)\right\rangle$ is dense in $N_{\widehat{G}}(H)$. As $G_1$ is polycyclic-by-finite it follows from Proposition 3.3 in \cite{RSZ} that $N_{G_{1}}(H)$ is dense in $N_{\widehat{G}_{1}}(H)$ and $N_{G_{1}^{t^{-1}}}(H)$ is dense in $N_{\widehat{G_{1}}^{t^{-1}}}(H)$, therefore  $N_{G}(H)$ is dense in $N_{\widehat{G}}(H)$. 

If $H$ and $K$ are conjugate subgroups in $\widehat{G_1}$, then by Proposition \ref{cp1} we have $$N_{G}(H)=\hnn(N_{G_1}(H), H, \tau_{g_1} f) \mbox{ and } N_{\widehat{G}}(H)=\hnn(N_{\widehat{G_1}}(H), H, \tau_{g_1} f).$$ In this case $N_{G}(H)=\left \langle N_{G_1}(H), tg_1^{-1}\right\rangle$ and $N_{\widehat{G}}(H)=\overline{\left \langle N_{\widehat{G_1}}(H), tg_1^{-1}\right\rangle}$, and the result follows similarly as in the case of the non-conjugate associated subgroups. 

\medskip
%If $H$ is finite then 
Since $H$ is finite we have 
  $$\overline{N}_{\widehat{G}}(H)=\overline{N}_{G}(H),$$ since 
$N_{G}(H)$ is dense in $N_{\widehat{G}}(H)$. 


%$$\overline{N}_{\widehat{G}}(H)=\left\langle \overline{N}_{\widehat{G_1}}(H),\overline{N}_{\widehat{G_2}}(H)\right\rangle=\left\langle \overline{N}_{G_1}(H),\overline{N}_{G_2}(H)\right\rangle=\overline{N}_{G}(H).$$ 

\end{proof}

%We conclude this section observing that the proof of Proposition \ref{cp1} is valid for any free pro-$\C$ product with finite amalgamation, where $\C$ is a class of finite groups closed for subgroups, quotient and extensions. We state it here for future references.

%\begin{prop}\label{profinite normalizer} Let  $\C$ be a class of finite groups closed for subgroups, quotients and extensions and $G=G_1\amalg_H G_2$ be a free pro-$\C$ product with finite amalgamation. Then $N_G(H)=N_{G_1}(H)\amalg_H N_{G_2}(H)$.\end{prop}

\subsection{OE-groups} \label{OE groups}

%A finitely generated residually finite group $G$ with less than 2 ends, i.e. a finite or one ended group, will be called an $OE$-group in the paper.   It follows from the famous Stallings theorem that $G$ is an $OE$-group  if and only if whenever it acts on a tree with finite edge stabilizers
%it  has a global fixed point. A finitely generated profinite group will be called $OE$-group if it has the same property:   whenever it acts   on a  profinite tree  with finite edge stabilizers, then it fixes a vertex. 

 %A finitely generated residually finite group $G$  will be called $\widehat{OE}$-group if $\widehat G$ is $OE$-group. Note that an $\widehat{OE}$-group is automatically $OE$-group, since if a residually finite group $G$ splits as an amalgamated free product or HNN-extension over a finite group then so does $\widehat G$.  

%The next proposition gives a sufficient condition for $G$ to be an $\Oe$-group.

A finitely generated residually finite group $G$ with less than 2 ends, i.e. a finite or one-ended group, will be called an $\OE$-group in the paper.   It follows from the famous Stallings theorem \cite{Sta} that $G$ is an $\OE$-group  if and only if whenever it acts on a tree with finite edge stabilizers
it  has a global fixed point. A finitely generated profinite group will be called $\OE$-group if it has the same property:   whenever it acts on a  profinite tree  with finite edge stabilizers, then it fixes a vertex. 

 A finitely generated residually finite group $G$  will be called $\Oe$-group if $\widehat G$ is $\OE$-group. Note that an $\Oe$-group is automatically $\OE$-group, since if a finitely generated residually finite group $G$ splits as an amalgamated free product or HNN-extension over a finite group then so does $\widehat G$.  


Since we do not know whether 1-endedness is a profinite property, we do not know any single example of finitely generated residually finite 1-ended group which is not $\Oe$.  In this section we prove   
 several propositions that give  sufficient conditions for $G$ to be an $\Oe$-group  showing that the class of $\Oe$-groups is quite large and contains many important families of examples of $\Oe$-groups.


\begin{prop}\label{virtually cyclic}(\cite[Proposition 3.1]{BPZ}) Let $G$ be a finitely generated residually finite group such that $\widehat G$ does not have a non-abelian free pro-$p$ subgroup. If  $G$ is not virtually infinite cyclic, then $G$ is an $\Oe$-group.
\end{prop}


%The class of residually finite $\widehat{OE}$-groups is quite large. All arithmetic groups of $rank >1$ as well as indecomposable into a free product  Fuchsian groups, 3-manifold groups and Right Angled Artin groups and many others are $\Oe$-groups.
 
 
If $G$   satisfies an identity, then $\widehat G$ satisfies the same identity and so does not have non-abelian free  pro-$p$ subgroups, so it is an $\OE$ and an $\Oe$-group by Proposition \ref{virtually cyclic} unless it is virtually infinite cyclic. %All residually finite $Fab$ groups and in particular just-infinite groups are also $\widehat{OE}$-groups.  


The result below was proven in Propositions 3.2 and 3.3 of \cite{BPZ2} and gives several examples of $\Oe$-groups.

\begin{prop}\label{examples} The following groups are $\Oe$-groups. \begin{enumerate}
\item[(i)]  Virtually surface groups;
\item[(ii)] Indecomposable (into a free product) 3-manifold groups;
\item[(iii)] Indecomposable (into a free product) finitely generated Right-Angled Artin groups;
 \item[(iv)] Arithmetic groups of $rank >1$;
 \item[(v)] Just infinite and finitely generated residually finite Fab groups. 
 \end{enumerate}
\end{prop}  

The following result shows that  the class $\OE$-groups is  closed for extensions and commensurability.

\begin{prop} \label{extensao de OE group}  \begin{enumerate}
    \item[(i)] Let $H$ be a normal (closed) subgroup of a (profinite) group $G.$ If $H$ and $G/H$ are (profinite) $\OE$-groups then $G$ is a (profinite) $\OE$-group; 
      
    \item[(ii)]  Let $H$ be a (closed) subgroup of finite index in a (profinite) group $G.$ If $H$ is an (profinite) $\OE$-group then $G$ is an (profinite) $\OE$-group.     
\end{enumerate}
\end{prop} 

\begin{proof} Let $T$ be a (profinite) tree with finite edge stabilizers on which $G$ acts (continuously) and without inversions.  

1. Note that $G/H$ acts on the tree $T^{H} \neq \emptyset.$ Therefore $G/H$ fixes some vertex of $T^{ H}$ and therefore $ G$ fixes a vertex of $T.$ 

 2.  Let $H_{G}$ be the core of $H$ in $G.$ It is clear that $H_{G}$ is a (closed) normal subgroup of finite index in $G$ (hence an open subgroup in the profinite case). Therefore $T^{H_{G}} \neq \emptyset$ and $G/ H_{G}$ acts on the tree $T^{H_{G}}.$ As $G/ H_{G}$ is finite, by \cite[Theorem 2.10]{ZM} $T^{H_{G}}$ has a fixed point for this action, therefore  $T^{G} \neq \emptyset.$       
\end{proof} 


\subsection{ General results} \label{general results}

The following results will be important  for  next Sections. We shall prove the profinite version of the first statement only, as the abstract version follows mutadis mutandis.  

\begin{prop}\label{normal isomorfismo} Let $G = \hnn(G_{1}, H, f, t)$ and $B = \hnn(B_{1}, H_1, f_1, t_1)$ be abstract HNN-extensions with $H$ and $H_1$ finite associated subgroups of $G_1$ and $B_1$, respectively. 
Suppose $G_1$ is an $\OE$-group and $B$ is finitely generated  residually finite. If $G\cong B$ then there exists an isomorphism $\psi: G \longrightarrow B$ such that $\psi(G_1)= B_1$ and   $\psi(H)= H_1^{t_1^{\epsilon}}, \mbox{ where } \epsilon \in \{0,1\}.$  
\end{prop}



\begin{prop} \label{profinito normal iso} Let $G = \hnn(G_{1}, H,f,t)$ and $B = \hnn(B_{1}, H_1,f_1,t_1)$ be profinite HNN-extensions  with $H$ and $H_1$ finite associated subgroups of $G_1$ and $B_1$, respectively. Suppose $G_1$ is a profinite $\OE$-group and $B_1$ is finitely generated. If $G\cong B$ then there is a continuous isomorphism $\psi: G \longrightarrow B$ such that $\psi(G_1)= B_1$ and $\psi(H)= H_1^{t_1^{\epsilon}}$, where $ \epsilon \in \{0,1\}$.
\end{prop}

\begin{proof} Let $\psi: G\rightarrow  B$ be an isomorphism. Then $\psi(G_1)$ is conjugate into $B_1$, say $\psi(G_1)\leqslant {B_{1}}^{w}$ for some $w\in B$ (cf. \cite[Example 6.3.1]{R}). Thus replacing $\psi$ by $\tau_{w} \circ \psi$, if necessary (recall that $\tau_w$ is the inner  automorphism corresponding to $w$), we may assume that $\psi(G_1) \leqslant B_1$.  
Note that $H=G_1\cap G_1^{t^{-1}}$, so by Corollary 7.1.5 c) in \cite{R} we have $$\psi(H) = B_{1} \cap  B_{1}^{\psi(t^{-1})} \leqslant H_1^{t_1^{\epsilon}b}, \mbox{ where } b\in B_1 \mbox{ and } \epsilon \in \{0,1\}.$$ Thus replacing $\psi$ by $\tau_{b} \circ \psi$,  we get $\psi(H) \leq H_1^{t_1^{\epsilon}}.$ Suppose without loss of generality that $\psi(H) \leq H_1$. If the base groups of $G$ and $B$ are finite, applying the inverse homomorphism $\psi^{-1}$ in the previous calculations it can be shown that $\psi(G_1)=B_1$ and $\psi(H)=H_1$ because the groups will have the same orders, respectively. 

We now show that the containments $\psi(G_1) \leqslant B_1$ and $\psi(H) \leqslant H_1$ are in fact equalities even if the base subgroups are not finite. 
Since $\psi$ is an isomorphism $B=\hnn(\psi(G_{1}),  \psi(H), \ \psi \circ f\circ \psi^{-1}, \psi(t)  ).$ Let $U$ be an open normal subgroups of $B$ that intersect $H_1$  trivially. 
Set $U_1 := U \cap B_1$. Note that  $U_1$ is an open normal subgroup of $B_1$ satisfying $U_1 \cap H_1 =\{1\}=U_1 \cap \psi(H).$    
%Let $\widetilde U= \overline{ \langle\langle U_1\rangle\rangle}$ the (topological)  normal closure $\langle\langle U_1\rangle\rangle$ of subgroup $\langle U_1\rangle$ of $B.$ 
Let $\widetilde U= \overline{  U_{1}^{B}}$ the (topological) normal closure of subgroup $ U_1$ of $B.$ Consider the profinite HNN-extensions  $$\hnn( \psi(G_1)U_1/U_1,  \psi(H),  \psi \circ f\circ \psi^{-1}, \psi(t))\  {\rm and}\  \hnn( B_1/U_1,  H_1,  f_1, t_1)$$ of finite groups. We will show that such groups are isomorphic. Indeed, let $$\pi_1: B \rightarrow \hnn( B_1/U_1,  H_1,  f_1, t_1)$$ and $$\pi_2: B \rightarrow \hnn( \psi(G_1)U_1/U_1,  \psi(H),  \psi \circ f\circ \psi^{-1}, \psi(t))$$ be the natural epimorphisms. Note that in both cases the kernel is $\widetilde U$. Indeed to show that $Ker(\pi_2)=\widetilde U$ consider the following commutative diagram of epimorphisms
$$\xymatrix{B\ar[r]\ar[rd]& \hnn( \psi(G_1)U_1/U_1,  \psi(H),  \psi \circ f\circ \psi^{-1}, \psi(t))\ar[d]\\
& \hnn( B_1/U_1,  H_1,  f_1, t_1)}$$

Note that the restriction $$\hnn^{abs}( \psi(G_1)U_1/U_1,  \psi(H),  f^{\psi^{-1}}, \psi(t))\rightarrow \hnn^{abs}( B_1/U_1,  H_1, f_1, t_1)$$ of the vertical map is the natural map of abstract HNN-extensions and therefore is injective. Since these groups are virtually free, so is the vertical map of the diagram that coincides with the map of  their profinite completions. Thus the vertical map is an isomorphism. As the base groups are finite, it follows as previously discussed that $$(\psi(G_1)U_1)/U_1= B_1/U_1 \mbox{ and } H\cong H_1.$$ Since this holds for any $U$ such that $U\cap H_1=\{1\}$ we deduce that $\psi(H) = H_1$ and $\psi(G_1) =B_{1}.$ 
\end{proof}

The following result will be used later. It shows that  an HNN-extension cannot be isomorphic to an  amalgamated free product of $\OE$-groups.

\begin{prop} \label{HNN nao iso Amalg}
Let $G= \hnn(G_1, H_1, t_1, f_1)$ and $B=G_2 *_{H_2} G_3$ where $G_i \,(i=1,2,3)$ are $\OE$-groups and the $H_j \, (j=1,2)$ are finite. Then $G \not \cong B.$  
\end{prop}

\begin{proof} Suppose there exists an isomorphism $\psi: B \longrightarrow G.$ Since $G_i \, (i=2,3)$ are $\OE$-groups, we have that $\psi(G_i) \subset G_1^G \,(i=2,3),$ therefore $\psi$ is not surjective, an absurd. 
\end{proof} 

%\begin{proof} Suppose by absurd that $G \cong B$.  Then we have $\widehat G = \hnn(\widehat G_1, H_1, t_1) \cong \widehat B = \widehat G_{2} \coprod_{H_2} \widehat G_{3}.$ Denote by $G=\pi_1(\mathcal{G}, \Gamma)$ and $B= \pi_1(\mathcal{B}, \Delta)$ their respective JSJ-decompositions. By Proposition 4.2 in \cite{BPZ} we have that $G$ and $B$ belong to the class $\mathcal{A}$ (see this definition on p. 218 in \cite{BPZ}). Since $G= \hnn(G_1, H_1, t_1)$ and $B=G_2 *_{H_2} G_3$ are their respective JSJ-decompositions it follows from Theorem 1.1 in \cite{BPZ} that there are bijections between the respective sets of vertices and edges of $\Gamma$ and $\Delta$, which is a contradiction since $\Gamma$ is a loop and $\Delta$ is an edge with two distinct vertices.
%\end{proof}

The result below shows that an HNN-extension with conjugate associated finite subgroups  in the base group cannot be isomorphic to an HNN-extension with non-conjugate associated finite  subgroups in its base group. 
% We will give its proof below and in the profinite case, which we will see later, we will only state it because the proof is analogous to what we will do next (see Proposition \ref{normal non iso nao normal profinite}). 



\begin{prop} \label{normal non iso nao normal}
 Let $$G= \hnn(G_1, H, K, f, t)\  {\rm and}\  B=\hnn(B_1, H_1, K_1,f_1, t_1)$$ be HNN-extensions of an $\OE$-group $G_1$ with $H, K$ and $H_1, K_1$ finite subgroups of $G_1$ and $B_1$, respectively. Suppose $B$ is finitely generated residually finite. If $H$, $K$ are  conjugate in $G_1$ and $H_1$, $K_1$ are non-conjugated in $B_1$, then $G$ and $B$ are non-isomorphic. 
\end{prop} 

\begin{proof}
 Suppose on the contrary that there is an isomorphism $\psi: G \longrightarrow B.$ By Proposition \ref{normal isomorfismo} we have $\psi(G_1)=B_1$ and $\psi(H)=H_1 \mbox{ or } K_1.$  Let us assume without loss of generality that $\psi(H)=H_1.$ Then $$N_G(H) \cong \psi\left(N_{G}(H)\right)=N_{B}(H_1).$$ By Proposition \ref{cp1} \begin{equation} \label{cara dos normalizers} N_{G}(H)=\hnn(N_{G_1}(H), H, f') \cong N_{B}(H_1)=N_{B_1}(H_1)*_{H_1
 } N_{{B_1}^{t_{1}^{-1}}}(H_1).\end{equation}However, this is impossible since $N_G(H)\not \leq G_1^G$ and $N_B(H_1)\leq B_1^B$. 
\end{proof}

%\begin{proof}
 %Suppose by absurd that there is a isomorphism $\psi: G \longrightarrow B.$ By Proposition \ref{normal isomorfismo} we have $\psi(G_1)=B_1$ and $\psi(H)=H_1 \mbox{ or } K_1.$  Let us assume without loss of generality that $\psi(H)=H_1.$ Therefore $N_G(H) \cong \psi\left(N_{G}(H)\right)=N_{B}(H_1).$ By Proposition \ref{cp1} we have \begin{equation} \label{cara dos normalizers} N_{G}(H)=\hnn(N_{G_1}(H), H, f') \cong N_{B}(H_1)=N_{B_1}(H_1)*_{H_1
 %} N_{{B_1}^{t_{1}^{-1}}}(H_1),\end{equation} whence it follows that  $$N_{\widehat G}(H)=\hnn(N_{\widehat G_1}(H), H, f') \cong N_{\widehat B}(H_1)=N_{\widehat B_1}(H_1)\coprod_{H_1
 %} N_{{\widehat B_1}^{t_{1}^{-1}}}(H_1).$$ Since $G_1\cong B_1$ it follows that $B_1$ is an OE-group, whence it follows that $N_{G}(H)$ and $N_{B}(H_1)$ belong to the class $\mathcal{A}$ by Proposition 4.2 in \cite{BPZ} since the normalizers that are the vertex groups in $N_{G}(H)$ and $N_{B}(H_1)$ (see (\ref{cara dos normalizers})) are also OE-groups. As in Proposition \ref{HNN nao iso Amalg} we have a contradiction by Theorem 1.1 of \cite{BPZ} since on the one hand the associated graph will be a loop and on the other an edge with two distinct vertices.  
%\end{proof}

Using the same arguments as the previous proof and the Propositions \ref{cp1} %and \ref{profinito normal iso} 
we have the following profinite version: 

\begin{prop} \label{normal non iso nao normal profinite}
Let $G_1, B_1$ be profinite $\OE$-groups, $H, K$ and $H_1, K_1$ finite  subgroups of $G_1$ and $B_1$ respectively. Let $G= \hnn(G_1, H, K,  f,  t)$ and $B=\hnn(B_1, H_1, K_1, f_1, t_1)$ be profinite HNN-extensions such that in the first the associated subgroups are conjugate in $G_1$ and in the second they are non-conjugate in $B_1.$ Then $G$ and $B$ are non-isomorphic.  
\end{prop} 


%In this section we will deal with the isomorphism problem for profinite HNN-extensions. The proof of the following proposition is analogous to the abstract case (see Proposition \ref{fixed hnn abstract}), however we will give the proof for the reader's convenience. 


The following results will be important  for Sections \ref{iso abstract HNN} and \ref{secao 4}. We shall prove the profinite version of the  statement only, as the abstract version follows mutadis mutandis. 

\begin{prop}\label{fixed hnn abstract} Let $$G = \hnn(G_{1}, H,K, f, t)\  {\rm and}\  B = \hnn(B_{1}, H_1, K_1, f_1, t_1)$$ be abstract HNN-extensions, where $H, K$ and $H_1, K_1$ are finite associated subgroups of $G_1$ and $B_1,$ respectively.  Suppose that $G_1$ is an $\OE$-group and $B$ is residually finite. If $G\cong B$ then there is an isomorphism $\psi: G \longrightarrow B$ such that $\psi(G_1)= B_1$, $\psi(H)=H_1$ and $\psi(K)= K_{1}^{b_1}$, where $ b_1\in B_1$. Furthermore, if $f_1=\tau_{b_1}f^{\psi^{-1}}\in \Iso(H_1,K_1)$ then the converse also holds.
%then the converse of the Proposition is true.
\end{prop}

\begin{prop}\label{2a} Let $$G = \hnn(G_{1}, H,K, f, t)\  {\rm and}\  B = \hnn(B_{1}, H_1, K_1, f_1,\\ t_1)$$ be profinite HNN-extensions, where $H, K$ and $H_1, K_1$ are finite associated subgroups.  Suppose that $G_1$ be a profinite  $\OE$-group.  If $G\cong B$ then there is a continuous isomorphism $\psi: G \longrightarrow B$ such that $\psi(G_1)= B_1$, $\psi(H)=H_1$ and $\psi(K)= K_{1}^{b_1}$, where $ b_1\in B_1$. Furthermore, if $f_1=\tau_{b_1}f^{\psi^{-1}}\in \Iso(H_1,K_1)$ then the converse of also holds.
\end{prop}

\begin{proof} Let $\psi: G\rightarrow  B$ be a  continuous isomorphism. By Proposition \ref{profinito normal iso} we can assume that $\psi(H) = H_1$ and $\psi(G_1) =B_{1}.$ 

\textbf{Case 1)}.  $H$ and $K$ are conjugate subgroups of $G_1$, i.e.,  $H^{g_1}=K$ with $g_1\in G_1$. Then by  Proposition \ref{normal non iso nao normal profinite}  $H_1$ and $K_1$ are conjugate subgroups of $B_1$, i.e.,  $K_1^{s}=H_1$ with $s\in B_1$. Therefore $\psi(K)=\psi(H^{g_1})=H_1^{\psi(g_1)}=K_1^{s\psi(g_1)}$. Put $b_1=s\psi(g_1) \in B_1$. Then $\psi(K)=K_1^{b_1}$ as needed.

\textbf{Caso 2)}.  $H$ and $K$ are not conjugate subgroups of $G_1$. Then by Proposition \ref{normal non iso nao normal profinite} $H_1$ and $K_1$ are not conjugates in $B_1$. It follows that
 $\psi(K) =\psi(G_1\cap G_1^{t})= B_1 \cap  B_1^{\psi(t)} =K_1^{b_1}$ for some $b_1\in B_1$  (we cannot have $\psi(K) = H_{1}^{ b_1}$ because  $H$ and $K$ are not conjugate in $G_1$).

\bigskip
Conversly, if $f_1=\tau_{b_1}f^{\psi^{-1}}$ we define a map $\varphi$  from $G$ to $B$ as follows $$\varphi: \left\{
\begin{array}{ccl}
g_1 & \longmapsto  & \psi(g_1), \forall  g_1 \in  G_1,\\
t & \longmapsto  & t_1 b_1\\
\end{array}
\right.$$ and observe that $\varphi$ is a continuous isomorphism.
\end{proof}

The result below is valid for abstract and profinite HNN-extensions.  

\begin{lem} \label{conjugados da normal} Let $G=\hnn(G_1, H, K, f, t)$ be an  abstract (profinite) HNN-extension such that $H^{g_1}=K$ for some $g_1 \in G_1.$ Then $$G \cong B=\hnn(G_1, H, H, f_1, t_1),$$  where $f_1=\left( \tau_{g_1} \circ f \right)^{\pm 1} \in \aut(H).$
\end{lem}

\begin{proof} First suppose that $f_1=\tau_{g_1} \circ f.$ Define the following isomorphism of $G$ to $\hnn(G_1, H, H, \tau_{g_1} \circ f, t_1),$  
 $$\psi: \left\{
\begin{array}{ccl}
g_1 & \longmapsto  & g_1, \forall  g_1 \in  G_1,\\
t & \longmapsto  & t_1 g_1.\\
\end{array}
\right.$$ Looking at presentation one sees indeed that this defines uniquely an isomorphism. To finish the proof, define the  isomorphism  $$\hnn(G_1, H, H, \tau_{g_1} \circ f, t_1)\longrightarrow \hnn(G_1, H, H, (\tau_{g_1} \circ f)^{-1}, t_{1}^{})$$   setting
 $$\varphi: \left\{
\begin{array}{ccl}
g_1 & \longmapsto  & g_1, \forall  g_1 \in  G_1,\\
t_1 & \longmapsto  & t_1^{-1}.\\
\end{array}
\right.$$ 
\end{proof}


%\begin{cor}\label{conjugados da normal} Let $G=\hnn(G_1, H, K, f, t)$ be an (abstract or profinite) HNN-extension such that $H^{g_1}=K$ for some $g_1 \in G_1.$ Then $G \cong B=\hnn(G_1, H, f_1, t_1),$  where $f_1=(\tau_{g_1} \circ f)^{\pm 1} \in Aut(H).$ 
%\end{cor}

%\begin{proof} Suppose $f_1=\tau_{g_1} \circ f$. Then by Proposition \ref{2a} $G\cong \hnn(G_1,H,K,f_1,t_1)$. To finish the proof note that  $$\varphi: \left\{
%\begin{array}{ccl}
%g_1 & \longmapsto  & g_1, \forall  g_1 \in  G_1,\\
%t_1 & \longmapsto  & t_1^{-1}\\
%\end{array}
%\right.$$ 
%define an isomorphism between $\hnn(G_1,H,K,f_1,t_1)$ and $\hnn(G_1,H,K,f_1^{-1},t_1)$.
%\end{proof}

%\begin{cor}\label{class B} Let $G = \hnn(G_{1}, H, K, f, t)$ be an abstract HNN-extension with finite associated subgroups. Suppose $G_1$ is an  $\Oe$-group  and $B$ is a finitely generated residually finite not 1-ended group such that $\widehat G \cong \widehat B$. Then $B$ is an abstract HNN-extension $\hnn(B_{1},H_1,K_1, f_1, t_1)$ with $\widehat G_1\cong \widehat B_1$ and $H\cong H_1$.
%\end{cor}
%\begin{proof}
%Let $B$  be a finitely generated residually finite group and 
%$\psi: \widehat{G} \longrightarrow \widehat{B}$ is a continuous isomorphism. Then $B$ is infinite and by Stalling's Theorem (see \cite{Sta}), $B$ is either a non-fictitious free product with finite  amalgamation or an HNN-extension with finite associated subgroups. 

%Suppose $B$ decomposes as a non-fictitious free product with finite amalgamation, say $B=B_1*_{H_1} B_2$, but does not decompose as an HNN extension. Then $\widehat B = \widehat B_{1} \amalg_{H_1} \widehat B_{2}$ and since $G_1$ is an $\Oe$-group, $\psi(\widehat G_1)$ will be contained in a conjugate of $\widehat B_1$ or $\widehat B_2,$ say $\psi(\widehat G_1) \leq (\widehat B_1)^b$ for some $b \in \widehat{B}.$ Thus replacing $\psi$ by $\tau_{b} \circ \psi$, if necessary (recall that $\tau_b$ is the inner of automorphism corresponding to $b$), we may assume that $\psi(G_1) \leqslant \widehat B_{1}$. Note that $H=\widehat G_1 \cap \widehat G_1^{t^{-1}},$ and so by Proposition 2.1 (ii) in \cite{BPZ2} that $$\psi(H) \leq \widehat{B}_{1} \cap \widehat{B}_{1}^{\psi(t^{-1})} \leqslant H_1^{b'}, \mbox{ where } b'\in \widehat{B}_1.$$ 
%As done previously, we can once again replace $\psi$ by $\tau_{b'} \circ \psi$, if necessary, to assume that  $\psi(H) \leq H_{1}.$ Since $\psi$ is an isomorphism $\widehat B=\hnn(\psi(\widehat{G}_{1}),  \psi(H), \ \psi \circ f, \psi(t)  ).$ Let $U$ be an open normal subgroups of $\widehat B$ that intersect $H_1$  trivially and such that $U\psi(H)\neq U\psi(G_1)$.
%Set $U_i := U \cap \widehat{B}_i (i=1,2)$. Note that the $U_i$ is an open normal subgroup of $\widehat{B}_i \,(i=1,2)$ satisfying $U_1 \cap H_1 =\{1\}=U_1 \cap \psi(H)$ and $U_1\psi(H)\neq \psi(\widehat G_1)$. Let 
%$W=\overline{ \langle U_1, U_2\rangle^{\widehat{B}} }$ 
%$\widetilde U$ be the (topological)  normal closure of $U_1$ in $\widehat B$.
%subgroup $\langle U_1, U_2\rangle$ of $\widehat B$. %Let $W= \overline{ \langle\langle U_1, U_2\rangle\rangle}$ the (topological)  normal closure of subgroup $\langle U_1, U_2\rangle$ of $\widehat B$. 
%Since $\widetilde U \leq U$ we have that $\psi(H) \cap \widetilde=\{1\}$ and therefore $\psi(H)\widetilde U/\widetilde U \cong \psi(H);$ furthermore, since $U_1 \subset \widetilde U$ is an open normal subgroup of $\widehat B_{1}$ it follows that $\psi(\widehat G_1)\widetilde U/\widetilde U$ is a finite group. Now consider the profinite HNN-extension $\hnn( \psi(\widehat G_1) U_1/U_1,  \psi(H),  \psi \circ f, \psi(t))$ and the profinite amalgamation free product $\left( \widehat{B}_1/U_1 \right) \amalg_{H_1} \left(\widehat{B_2} \right)$. Let $\pi_1$ and $\pi_2$ be the respective epimorphisms: $\pi_1: \widehat{B} \twoheadrightarrow \left( \widehat{B}_1/U_1 \right) \amalg_{H_1} \left(\widehat{B_2} \right)$ and $\pi_2: \widehat{B} \twoheadrightarrow \hnn( \psi(\widehat G_1)U_1/U_1,  \psi(H),  \psi \circ f, \psi(t)).$ Note that $\widetilde U$ is the kernel of the two epimorphisms which implies that  \begin{equation} \label{eq1}     
%\widehat{B}/ \widetilde U  \cong \left( \widehat{B}_1/U_1 \right) \amalg_{H_1} \widehat{B_2} \cong \hnn( \psi(\widehat G_1)U_1/U_1,  \psi(H),  \psi \circ f, \psi(t)),\end{equation}. 




%a contradiction with Proposition \ref{HNN nao iso Amalg}.  



%Denote by $\varrho$ the second isomorphism in (\ref{eq1}). Note that $\varrho(\widehat{B}_i/U_i)$ is contained in some conjugate of $\psi(\widehat G_1)W/W$ for each $i=1,2$ because $\widehat{B}_i/U_i$ is finite (see Corollary 7.1.3 (c) in \cite{R}), a contradiction to the fact that $\varrho$ is surjective.  

%Denote by $\varrho$ the isomorphism from $\left( \widehat{B}_1/U_1 \right) \amalg_{H_1} \left(\widehat{B_2} /U_2\right)$ to $\hnn( \psi(\widehat G_1)W/W, \psi(H), \psi \circ f, \psi(t))$. Note that $\varrho(\widehat{B}_i/U_i) \leq \psi(\widehat G_1)W/W$ for each $i=1,2$ since $\widehat{B}_i/U_i \, (i=1,2)$ are finite, a contradiction to the fact that $\varrho$ is surjective.

%which is a contradiction since on the one hand $\widehat{B}/W$ is generated by finite order elements and on the other hand it is generated by the element $\psi(t)$ of infinite order.  


%Therefore $B$ is an abstract HNN-extension, say $B=\hnn(B_{1}, H_1, K_1, f_1, t_1)$ where $H_1, K_1$ are finite groups, whence it follows by Proposition \ref{2a} that there exists an isomorphism $\varphi: \widehat G\rightarrow \widehat B$ such that
%$\varphi(\widehat G_1) = \widehat B_1$ and $\varphi(H)=H_1.$ %Since $G_1$ is profinitely rigid, it follows that $G_1\cong B_1.$
%\end{proof}

%\begin{cor} \label{rigidez} 
 %Suppose in addition to the hypothesis of Corollary \ref{class B}  $G_1$ is profinitely rigid. Then in addition  one gets $G_1\cong B_1.$\end{cor}

%\begin{proof} In the proof of Corollary \ref{class B} we proved that there exists an isomorphism $\varphi: \widehat G\rightarrow \widehat B$ such that
%$\varphi(\widehat G_1) = \widehat B_1$ and $\varphi(H)=H_1.$ Since $G_1$ is profinitely rigid, it follows that $G_1\cong B_1.$\end{proof}


\begin{prop}\label{nonnormalhnnabs} 
Let  $$G= \hnn(G_1, H, K, f, t)\  {\rm and}\  B=\hnn(G_1, H, K, f_1, t_1)$$ be HNN-extensions  with  associated finite subgroups in an $\OE$-group $G_1$. Suppose that $K=H$ or $H$ and $K$ are not conjugate in $G_1$. Then $G\cong B$ if and only if \begin{enumerate}


\item [(i)] there exist $\psi\in \aut(G_1)$ with $\psi(H)=K $ and $ \psi(K)=H^{g_1}$ for some $g_1 \in G_1$ such that $f_1^{\psi^{-1}}(k)= \tau_{g_1^{-1}} f^{- 1}(k), \forall k \in K,$ % where $w=g_1^{-1}x$ with $x\in N_{G_1}(H)$ 
or 
\item [(ii)] there exist $\psi\in \aut(G_1)$ with $\psi(H)=H $ and $ \psi(K)=K^{g_1}$ for some $g_1 \in G_1$ such that $f_1^{\psi^{-1}}(h)= \tau_{g_1^{-1}} f(h), \forall h \in H.$ 
\end{enumerate}



\end{prop}

\begin{proof}
$(\Rightarrow)$ By Proposition \ref{fixed hnn abstract} there is an isomorphism $\psi: B \longrightarrow G$ such that $\psi(G_1) = G_1$,  $\psi(H)=H \mbox{ and }$ $\psi(K)=K^{w}$ or  $\psi(G_1) = G_1$, $\psi(H)=K \mbox{ and } \psi(K)=H^{w}$ for some $w\in G_1$.  

 If $H=K$ then $w\in N_{G_{1}}(H)$ and $\psi(t_1)=h_1t^{\pm 1}h_2\in N_{G}(H)$ with $h_1, h_2 \in N_{G_1}(H)$ by Lemma 2.11, p. 822 in \cite{BZ}.  As we have discussed throughout this article, we can replace the isomorphism $\psi$ by $\tau_{h_{1}^{-1}}\psi$, if necessary, to assume that $\psi(G_1)=G_1,$
  $\psi(H)=H$ and $\psi(t_1)= t^{\pm 1}h_2h_1.$ Put $g_1=h_2h_1 \in N_{G_1}(H)$. Then for  so for all $h \in H,$ we have $$\psi(f_1(h))=\psi(h^{t_1})=\psi(h)^{t^{\pm 1}g_1}=(\tau_{g_1^{-1}}f^{\pm 1}  \psi)(h).$$ 
 Therefore $f_1^{\psi^{-1}}(h) = \tau_{g_1^{-1}}f^{\pm 1}(h), \forall h \in H$ and (i) or (ii) holds. 

If $H$ and $K$ are not conjugate in $G_1$. We will divide the proof  into two parts.
  
 \textbf{ Case i)} Suppose that $\psi(G_1) = G_1$, $\psi(H)=K \mbox{ and } \psi(K)=H^{w}$ with $w\in G_1$.  Put $\psi(t_1) = g \in G$. Note that $$\psi(K) = H^{w}=\psi(H)^{\psi(t_1)}=K^g=H^{tg},$$ so that $w g^{-1} t^{-1} \in N_G(H)$. It follows from Lemma 2.11 in \cite{BZ} that $w g^{-1}=xty \in \left( N_{G_1}(H) \cdot t \cdot N_{G_1}(K) \right)$. Thus replacing $\psi$ by $\tau_y \circ \psi$, if necessary (where $\tau_y \in \overline{N}_{G_1}(K)$), we may assume that  $\psi(G_1)=G_1$, $\psi(H)=K$, $\psi(K)=H^{g_1}$ and $\psi(t_1)=t^{-1}g_1$, where $g_1=x^{-1}wy^{-1}.$ %for some $x\in N_{G_1}(H)$ and $y\in N_{G_1}(K)$. 
 For each $h \in H$, we have $$\psi(f_1(h))=\psi(h^{t_1})=\psi(h)^{\psi(t_1)}=\psi(h)^{t^{-1}g_1}=\tau_{g_1^{-1}} f^{-1}(\psi(h)).$$ 
 Therefore $f_1^{\psi^{-1}}(k)= \tau_{g_{1}^{-1}} f^{- 1}(k), \forall k \in K$ and (i) holds.
 
 
 \textbf{ Case ii)} Suppose $\psi(G_1) = G_1$,  $\psi(H)=H \mbox{ and }$ $\psi(K)=K^{w}$. Similar as in case (i), put $\psi(t_1) = g \in G$ and observe that $gw^{-1}t^{-1}\in N_{G}(H)$. It follows from Lemma 2.11 in \cite{BZ} that $g w^{-1}=xty \in \left( N_{G_1}(H) \cdot t \cdot N_{G_1}(K) \right)$. Thus replacing $\psi$ by $\tau_{x^{-1}} \circ \psi$, if necessary (where $\tau_{x^{-1}} \in \overline{N}_{G_1}(H)$), we may assume that $\psi(G_1)=G_1$, $\psi(H)=H$, $\psi(K)=K^{g_1}$ and $\psi(t_1)=tg_1$, where $g_1=ywx.$ %for some $x\in N_{G_1}(H)$ and $y\in N_{G_1}(K)$. 
 For each $h \in H$, we have$$\psi(f_1(h))=\psi(h^{t_1})=\psi(h)^{\psi(t_1)}=\psi(h)^{tg_1}=\tau_{g_1^{-1}} f(\psi(h)),$$
 Therefore $f_1^{\psi^{-1}}(h)= \tau_{g_1^{-1}} f(h), \forall h \in H$ and (ii) holds. 
 

$(\Leftarrow)$ Follows from the Proposition \ref{fixed hnn abstract}. %\ref{2a}.

\end{proof} 

%\begin{lem} Let $G=A\amalg_{H,\varphi} A$ be a free amalgamated product where $H$ is a subgroup of $A$ and $\varphi:H\longrightarrow  A$ be a monomorphism into a second factor. Then it embeds naturally into $HNN(A,H, \varphi, t)$. Moreover, if $G'=A\amalg_{H\varphi'} A$ is another such amalgam, then $G\cong G'$ if and only if  $HNN(A,H, \varphi', t')$. \end{lem}

%\begin{proof} Clearly $G=\langle A, A^{t}\rangle$ in $HNN(A,H, \varphi, t)$. Now if $G\cong G'$ then presentations $\langle A, A\mid rel (A), h=\varphi(h)\rangle$ and  $\langle A, A\mid rel (A), h=\varphi'(h)\rangle$ are equivalent. So presentations $\langle A, t\mid rel (A), h^t=\varphi(h)\rangle$ and $\langle A, A\mid rel (A), h^{t'}=\varphi(h)\rangle$ are equivalent and so $HNN(A,H, \varphi', t')$ and $HNN(A,H, \varphi, t)$ isomorphic.

%\end{proof} 

%\begin{prop} Let $G= \hnn(G_1, H, K, f, t)$ and $B=\hnn(G_1, H_1, K_1, f_1, t_1)$ be HNN-extensions (resp. profinite HNN-externsions) with  associated finite subgroups in an $\OE$-group $G_1$. Suppose that $K=H$ or $H$ and $K$ are not conjugate in $G_1$. Then $G\cong B$ if and only if $\langle G_1, G_1^t\rangle \cong \langle G_1, G_1^{t_1}$. \end{prop}

%\begin{proof} We prove it for profinite case, abstract being similar. 

%'íf'.   $\langle G_1, G_1^t\rangle \cong \langle G_1, G_1^{t_1}$. Note that $\langle G_1, G_1^t=G_1\amalg_H G_1^t$ and $\langle G_1, G_1^{t_1}=G_1\amalg_H G_1^{t_1}$. Then by \cite[Proposition 3.5]{BPZ} there exists an ismorphism $\psi:G\longrightarrow B$ such that $\psi(G_1)=G_1, \psi(H)=H$ and $\psi(G_1^t)=(G_1^{t_1})^b$ for some $b\in B$. Moreover, by \cite[Proposition 5.1]{BPZ} $b\in N_B(H)$. Form an HNN-extension $HNN(G_1, H, K_1^r , t_1r)$. 

%\end{proof}


\section{Isomorphism problem for abstract HNN-extensions} \label{iso abstract HNN}
%\section{ Isomorphism  Problem} \label{secao 5}


In this section we will deal with the isomorphism problem of abstract HNN-extensions. 

%\begin{prop}\label{nonnormalhnnabs} 
%Let $G_1$ be an OE-group and $H, K$ non-conjugated finite  subgroups in $G_1.$ Let $G= HNN(G_1, H, K, t_1, f_1)$ and $B=HNN(G_1, H, K, t_2, f_2)$ be HNN-extensions with $f_i(H)=K\,(i=1,2)$. Then $G\cong B$ if and only if \begin{enumerate}
%\item [(i)] there exist $\phi\in Aut(G_1)$ with $\phi(H)=K, \phi(K)=H^{\tilde{g_1}}$ for some $\tilde{g_1} \in G_1$ such that $f_2^{\phi^{-1}}=\phi f_2\phi^{-1}= \tau_w f_1^{- 1}$ for some $w \in G_1$, or
%\item [(ii)] there exist $\phi\in Aut(G_1)$ with $\phi(H)=H, \phi(K)=K^{\tilde{g_1}}$ for some $\tilde{g_1} \in G_1$ such that $f_2^{\phi^{-1}}=\phi f_2\phi^{-1}= \tau_w f_1$ for some $w \in G_1$. 
%\end{enumerate}
%\end{prop}



%\begin{prop}\label{nonnormalhnnabs} 
%Let $G_1$ be an OE-group and $H, K$ non-conjugated finite  subgroups in $G_1.$ Let $G= HNN(G_1, H, K, t_1, f_1)$ and $B=HNN(G_1, H, K,\\ t_2, f_2)$ be HNN-extensions with $f_i(H)=K\,(i=1,2)$. Then $G\cong B$ if and only if \begin{enumerate}
%\item [(i)] there exist $\psi\in Aut(G_1),\, x \in N_{G_1}(H),\, y \in N_{G_1}(K)$ with $\psi(H)=K, \psi(K)=H^{g_1}$ for some $g_1 \in G_1$, such that $f_2^{(\tau_y \psi)^{-1}}= \tau_w (f_1)^{- 1},$ where $w=y g_{1}^{-1}x$, or
%\item [(ii)] there exist $\psi\in Aut(G_1),\, x \in N_{G_1}(H),\, y \in N_{G_1}(K)$ with $\psi(H)=H, \psi(K)=K^{g_1}$ for some $g_1 \in G_1$, such that $f_2^{(\tau_{x^{-1}} \psi)^{-1}}= \tau_{w^{-1}} f_1,$  where $w=y g_{1}x.$
%\end{enumerate}
%\end{prop}
%and $H, K$ non-conjugated finite  subgroups in $G_1.$
 
%\subsection{Abstract HNN-extensions}
Let $G=\hnn(G_1, H, K, f_1, t_1) \mbox{ and } B=\hnn(G_1, H, K, f_2, t_2)$ be abstract HHN-extensions, where $G_1$ is an OE-group and $f_{i}: H\longrightarrow K\,(i=1,2)$ are  isomorphisms between finite subgroups $H, K$ of $G_1$, i.e $f_i \in \Iso(H, K)\,(i=1,2)$. We find necessary and 
sufficient conditions for $G\cong B$ in terms of the orbits of the actions that will be defined below. Note that we are not trying to find an algorithm that decides whether given HNN-extensions are isomorphic. Our goal is to describe in terms of HNN-extensions when they are isomorphic. 
%Let $G=HNN(G_1, H, K, f, t)$ be an HNN-extension where $G_1$ is a OE-group and $H, K$ are non-conjugate finite associated subgroups in $G_1.$
Some parts of the text in this section are adaptations of Sections 3 and 4 in \cite{BZ}.

Let $G=\hnn(G_1, H, K, f, t)$ be an abstract HNN-extension.  Consider the subgroup 
 \begin{equation} \label{gamma H} \Gamma_H:=G_1 \rtimes \aut_{G_1}(H)\end{equation} of the holomorph $G_1\rtimes \aut(G_1)$. 
%a multiplication by setting 
%$(\tau_{g_1},\alpha)(\tau_{g_2},\beta)=(\tau_{g_1}%\tau_{g_2}^{\alpha^{-1}},\alpha\beta)$. 
Sometimes we shall interpret the action of an element $g_1$ of $G_1$ in $\Gamma_H$ as the inner automorphism $\tau_{g_1},$ which should not be confused with the element of $\aut(G_1)$ (from the copy of $\inn(G_1)$ in $\aut(G_1)$).
%Sometimes we shall interpret an element $g_1$ of $G_1$ in $\Gamma_H$ as the element of the inner automorphism $\tau_{g_1}$ that should not be confused with the element of $\aut(G_1)$ (from the copy of $\inn(G_1)$ in $\aut(G_1)$).

Denote by $S$ the set of all finite subgroups of $G_1$. There is a natural action of $\Gamma_{H}=G_1\rtimes \aut_{G_1}(H)$ on the set $S$ namely \begin{equation} \label{acaoemX} (g_1,\alpha)\cdot X= g_1\alpha(X)g_1^{-1} \in S, \mbox{ where } (g_1,\alpha) \in \Gamma_H, X \in S. \end{equation}
 
 
 Let us  consider now
$\Upsilon :=\displaystyle\bigcup_{X\in S} \Iso(H,X)$.  Then  the action of $\Gamma_H$ on $S$ induces the natural action
of  $\Gamma_H$  on the left on $\Upsilon$ as the following lemma shows.
%$\mathcal{K}$.

\begin{lem}\label{elementary}
The semidirect product  $\Gamma_{H}=G_1 \rtimes \aut_{G_1}(H)$
acts from the left upon $\Upsilon$ by
\begin{equation} \label{acaosemidireta} (g_1,\alpha) \cdot f=\tau_{g_1} \alpha f \alpha^{-1} = \tau_{g_1} f^{\alpha^{-1}}, \mbox{ where } (g_1,\alpha) \in \Gamma_H \mbox{ 
 and } f \in \Upsilon.
 \end{equation}
\end{lem}
\begin{proof} For showing that the action is well-defined choose arbitrary $f\in \Upsilon$ and $\gamma=(g_1,\alpha)\in \Gamma_{H}$. Then $(g_1,\alpha) \cdot f \in \Iso(H, L)$ for $L:=\left[ \alpha(f(H))\right]^{g_1^{-1}} \in S$ and so $\gamma \cdot f \in \Upsilon$. Now certainly $(1,1)\cdot f=f$ holds for all $f\in \Upsilon$. Finally, if  $(g_1,\alpha), (g_2,\beta) \in \Gamma_H$ and $f \in \Upsilon$ then we have $$(g_1,\alpha) \cdot ((g_2,\beta) \cdot f)=\tau_{g_1}(\tau_{g_2} f^{\beta^{-1}})^{\alpha^{-1}}=\tau_{g_1}\alpha\tau_{g_2}\beta f \beta^{-1}\alpha^{-1}=$$ $$=\tau_{g_1}\tau_{g_2}^{\alpha^{-1}}f^{(\alpha\beta)^{-1}}=((g_1,\alpha)(g_2,\beta)) \cdot f.$$
\end{proof}


%%%%%%%%%%%%%%%%%

%\subsubsection{Isomorphism problem for abstract HNN-extensions of an OE-group}  


 Note that for $f\in \Iso(H, K)$, we have $(g_1,\alpha) \cdot f \in \Iso\left(H, [\alpha(K)]^{g_1^{-1}}\right) \subset \Upsilon$ holds true ( see (\ref{acaosemidireta})).  
%Let $\Gamma_{HK}$ denote the set stabilizer of $\Iso(H,K)$ according to action of $\Gamma_H$ on  $\Upsilon$ ( see (\ref{acaosemidireta})). Then $\Gamma_{HK}$ is nothing but the point stabilizer of $K$ for $\Gamma_{H}$ acting on $\Xi$ ( see (\ref{acaoemX})). 


Suppose there exist $\psi \in \aut(G_1)$ that swaps the conjugacy classes of $H$ and $K$, i.e. such that $\psi(K)=H^{g_1}$ and $\psi(H)=K$, for some $g_1\in G_1$.  Define $\iota: \Iso(H,K) \longrightarrow \Iso(H,K)$ by \begin{equation} \label{def iota abst} \iota(f)=\tau_{\psi^{-1}(g_1^{-1})}\left(f^{\psi}\right)^{-1}, f \in \Iso(H,K).\end{equation} 
%Suppose there exist $\psi \in Aut(G_1)$ such that $\psi(K)=H^{g_1}$ and $\psi(H)=K$, for some $g_1\in G_1$.  Define $\iota: ISO(H,K) \longrightarrow ISO(H,K)$ by $$\iota (f)=\tau_{\psi^{-1}(g_1^{-1})}(f^{-1})^{\psi}.$$


% When $H=K$ we can choose $\psi=id$ and $g_1=1,$ so $\iota$ applies $f$ to $f^{-1}.$ For this situation we have an isomorphism between $\hnn(G_1,H,f)$ and $\hnn(G_1,H,f^{-1})$ (see Remark \ref{f e f inverse}). 

When $H=K$ we can choose $\psi=id$ and $g_1=1,$ so $\iota$ sends $f$ to $f^{-1}.$ In this situation we have an isomorphism between $\hnn(G_1,H,f)$ and $\hnn(G_1,H,f^{-1})$ (see Lemma \ref{conjugados da normal}). Note that $\iota$ will be a bijection from $\Iso(H, H)=\aut(H)$ to $\Iso(H, H)=\aut(H)$ in this case. More generally, $\iota$ belongs to the symmetric group $Sym(\Iso(H, K)),$ of $\Iso(H, K)$. 



We denote by $\Gamma_{HK}$ be the stabilizer of $K\in S$ of this action of $\Gamma_H$ on $S$. Then $\Gamma_{HK}$ leaves $\Iso(H, K)$ invariant according to the action of $\Gamma_H$ on $\Upsilon$.  
Since $\Gamma_{HK}$ acts on $\Iso(H, K),$ we can consider the natural image of $\Gamma_{HK}$ in $Sym(\Iso(H, K))$ and denote it by $\widetilde\Gamma_{HK}$.

%When $H=K$ then, as has been  said earlier, $f\mapsto f^{-1}$ gives rise to an isomorphism between $\hnn(G_1,H,f)$ and $\hnn(G_1,H,f^{-1})$ (see Remark \ref{f e f inverse}). In this case, $\iota$ will be the map  that sends $f$ to $f^{-1}$ (since we can choose $\psi = id$ and $g_1=1$), so $\iota$ is a function from $\Iso(H, H)=\aut(H)$ to $\Iso(H, H)=\aut(H)$.

%Denote by $T(\Upsilon)$ the group of transformations acting to the left of $\Upsilon$.  




\begin{lem} \label{semidiretocomC2} Suppose there exist $\psi \in \aut(G_1)$ that swaps the conjugacy classes of $H$ and $K$.  Then $\iota $ normalizes $\widetilde\Gamma_{HK}$ and so $\langle \widetilde\Gamma_{HK}, \iota  \rangle=\widetilde\Gamma_{HK} \cdot  \left\langle  \iota \right\rangle$ is a subgroup of $Sym(\Iso(H, K))$. \end{lem}  

\begin{proof}  We need to show that if $\tilde\omega\in \widetilde\Gamma_{HK}$ then $\tilde\omega^{\iota}\in \widetilde\Gamma_{HK}$. Indeed,  pick $\omega=(x, \alpha) \in \Gamma_{HK} $ such that $\tilde w$ is its image in $\widetilde\Gamma_{HK}.$  Choose $f\in \Iso(H,K)$ and observe that

$\tilde\omega^\iota. f=\iota^{-1}\tilde\omega(\tau_{\psi^{-1}(g_1^{-1})}(f^{-1})^\psi)= \iota^{-1}(\tau_{x\alpha\psi^{-1}(g_1^{-1})}(f^{-1})^{\psi\alpha^{-1}})=\newline =(f)^{\psi\alpha^{-1}\psi^{-1}}\tau_{\psi(\alpha\psi^{-1}(g_1)x^{-1})g_1^{-1}}=\tau_{\psi(\alpha\psi^{-1}(g_1)x^{-1})g_1^{-1}} (f)^{\psi\alpha^{-1}\psi^{-1}\tau_{\psi(\alpha\psi^{-1}(g_1)x^{-1})g_1^{-1}}}= \newline =(\tau_{\psi(\alpha\psi^{-1}(g_1)x^{-1})g_1^{-1}}, \tau_{g_1\psi(x\alpha\psi^{-1}(g_1^{-1}))}\psi\alpha\psi^{-1}). f$. But $$(\psi(\alpha\psi^{-1}(g_1)x^{-1})g_1^{-1}, \tau_{g_1\psi(x\alpha\psi^{-1}(g_1^{-1}))}\psi\alpha\psi^{-1}) \in \Gamma_{HK}$$ and its image in $Sym(\Iso(H, K))$ is $\tilde{w}^{\iota} \in \tilde{\Gamma}_{HK},$ as desired. 
%But 
%$$(\tau_{\psi(\alpha\psi^{-1}(g_1)x^{-1})g_1^{-1}}, \tau_{g_1\psi(x\alpha\psi^{-1}(g_1^{-1}))}\psi\alpha\psi^{-1})\in \widetilde\Gamma_{HK}$$ as required.
\end{proof}


%\begin{lem} \label{semidiretocomC2} Suppose there exist $\psi \in Aut(G_1)$ such that $\psi(K)=H^{g_1}$ and $\psi(H)=K$, for some $g_1\in G_1$. Let $\left\langle \Gamma_{HK}, \iota \right \rangle$ be the subgroup of \textcolor{blue}{$\mathcal{T}(\Iso(H, K))$} generated by $\Gamma_{HK}$ and $\iota$. Then \textcolor{blue}{$\left\langle \Gamma_{HK}, \iota \right \rangle=\Gamma_{HK} \rtimes \left\langle  \iota \right\rangle$.}

%Let $\overline{\Gamma}_{HK}$ denote the subgroup of \textcolor{blue}{$\mathcal{T}(\Iso(H, K))$} generated by $\Gamma_{HK}$ and $\iota_{_{\psi, g_1}}$. Then \textcolor{blue}{$\overline{\Gamma}_{HK}=\Gamma_{HK} \rtimes \left\langle  \iota_{_{\psi, g_1}}\right\rangle$.}



%\end{lem}  \begin{proof} \textcolor{red}{ We need to show that if $\omega\in \Gamma_{HK}$ then $\omega^{\iota}\in \Gamma_{HK}$. In fact, first $\omega^{\iota} \in \Gamma_H$ and if $f \in \Iso(H, K)$ we have $\omega^{\iota} \cdot f = (\iota^{-1} \omega \iota) \cdot f \in \Iso(H, K)$, that implies $\omega^{\iota} \in \Gamma_{HK}.$ Finally note that $\iota^k$ does not belong to $\Gamma_H$ for any $k \geq 1,$ thus $\Gamma_{HK} \cap \left\langle  \iota \right\rangle =\{1\}.$}
%\end{proof}

\begin{lem}\label{carmo} The subgroup $\widetilde\Gamma_{HK} \cdot  \left\langle  \iota \right\rangle$ of $Sym(\Iso(H, K))$ does not depend upon the choice of $\psi$.
\end{lem}

%\begin{proof} Consider $\zeta \in \aut(G_1)$ such that $\zeta(K)=H^{r}$ and $\zeta(H)=K$, for some $r\in G_1$. %Let %$\Theta:\Iso(H,K)\rightarrow\Iso(H,K)$
%Let $\Theta \in Sym(\Iso(H,K))$ such that  $\Theta(f):=\tau_{_{\zeta^{-1}(r^{-1})}}\left(f^{-1}\right)^{\zeta}$, where $f \in \Iso(H, K)$. Take $\upsilon=(\tau_{_{\zeta^{-1}(r^{-1}  \cdot g_1)}}, \zeta^{-1}\psi) \in \widetilde\Gamma_{HK}$ and observe that $\Theta=\upsilon \iota$. Then  $\Theta=\upsilon \iota$, i.e., $\Theta \in \widetilde\Gamma_{HK} \cdot \left\langle \iota\right\rangle$.
%\end{proof}
\begin{proof} Consider $\zeta \in \aut(G_1)$ such that $\zeta(K)=H^{r}$ and $\zeta(H)=K$, for some $r\in G_1$. Let $\Theta \in Sym(\Iso(H,K))$ such that  $\Theta(f):=\tau_{_{\zeta^{-1}(r^{-1})}}\left(f^{-1}\right)^{\zeta}$, where $f \in \Iso(H, K)$. Take $\upsilon=(\zeta^{-1}(r^{-1}  \cdot g_1), \zeta^{-1}\psi) \in \Gamma_{HK}$ such that $\tilde{\upsilon}$ is its image in $\widetilde\Gamma_{HK},$ and observe that $\Theta=\tilde{\upsilon} \iota \in \widetilde\Gamma_{HK} \cdot \left\langle \iota\right\rangle.$
\end{proof}




%{When $K$ is in the $\Gamma_{H}$-orbit of $H$, i.e., there exists $\gamma:=(\tau_g, \varphi) \in \Gamma_H$ such that $H=\gamma \cdot K=(\tau_g, \varphi) \cdot K =\varphi(K)^{g^{-1}}$. Define $\iota^{\gamma} := \gamma^{-1} \circ \iota \circ \gamma$. Note that $(\iota^{\gamma})^2 = \gamma^{-1} \iota \gamma \gamma^{-1} \iota \gamma=\gamma^{-1} \iota^2 \gamma =id$, moreover,  for $f \in \Iso(H, K)$, we have $$i^{\gamma}(f)=\gamma^{-1} \cdot \iota (\tau_g \varphi f \varphi^{-1})= ( (\tau_{g^{-1}})^{\varphi}, \varphi^{-1} ) \cdot (\varphi f^{-1} \varphi^{-1} \tau_{g^{-1}})=\tau_{g^{-1}}^{\varphi} f^{-1} \tau_{g^{-1}}^{\varphi},$$ therefore we obtain \begin{equation} \label{igamma} i^{\gamma}(f)=\tau_{_{\varphi^{-1}(g^{-1})}} f^{-1} \tau_{_{\varphi^{-1}(g^{-1})}}.\end{equation}  

%Then, having the commutative diagram $$\xymatrix{
%\Iso(H,K) \ar^{\gamma}[r] \ar@{.>}^{\iota^{\gamma}}[d]  & \Iso(H,H) \ar^{\iota}[d]   \\ \Iso(H,K) & \Iso(H,H), \ar_{\gamma^{-1}}[l]}$$ 
%in mind we can extend $\Gamma_{HK}$ to include the involution $\iota^{\gamma}$.  

%Denote by $T(\Upsilon)$ the group of transformations acting to the left of $\Upsilon$.  

%\begin{lem} \label{semidiretocomC2} Suppose that $\gamma(K)=H$ for some $\gamma=(\tau_{g},\varphi)\in \Gamma_{H}$. Let $\overline{\Gamma}_{HK}$ denote the subgroup of $T(\Upsilon)$ generated by $\Gamma_{HK}$ and $\iota^{\gamma}$. Then $\overline{\Gamma}_{HK}=\Gamma_{HK}\rtimes \left\langle \iota^{\gamma}\right\rangle$. %Also, for $\delta\in \Gamma_{A}$ with $\delta(B)=A$ we have $\left\langle \Gamma_{AB},\iota^{\delta}\right\rangle=\tilde{\Gamma}_{AB}$.
%\end{lem}

%\begin{proof} We need to show that if $\omega\in \Gamma_{HK}$ then $\omega^{\iota^{\gamma}}\in \Gamma_{HK}$. In fact, if $f \in \Iso(H, K)$ we have $\omega^{\iota^{\gamma}} \cdot f = (\iota^{\gamma} \omega \iota^{\gamma}) \cdot f \in \Iso(H, K)$, that implies $\omega^{\iota^{\gamma}} \in \Gamma_{HK}.$  



%The first fact follows easily by observing that if $\omega\in \Gamma_{AB}$ then $\omega^{\iota^{\gamma}}\in \Gamma_{AB}$. For the second statement it is only needed to show that $\iota^{\delta}\in \tilde{\Gamma}_{AB}$. Now $\delta^{-1}\gamma:\Iso(A,B)\rightarrow\Iso(A,B)$ implies that $\delta^{-1}\gamma \in \Gamma_{AB}$. Therefore $\iota^{\delta}=\delta^{-1}\gamma\iota^{\gamma}\gamma^{-1}\delta\in \tilde{\Gamma}_{AB}$, as claimed.

%\end{proof}

%With the notations in this subsection we have the following result.



%\begin{lem}\label{carmo} Suppose that $\gamma(K)=H$ for some $\gamma=(\tau_{g},\varphi)\in \Gamma_{H}$. Consider  $\psi\in \aut(G_1)$ such that $\psi(H)=K$ and $\psi(K)^{r}=H$ for some $r\in G_1$. Let $\Theta:\Iso(H,K)\rightarrow\Iso(H,K)$ such that  $\Theta(f):=\tau_{\psi^{-1}(r)}\left(f^{-1}\right)^{\psi}$, where $f \in \Iso(H, K)$. Then there is $\omega\in \Gamma_{HK}$ such that $\Theta=\omega \iota^{\gamma}$, i.e., $\Theta \in \overline{\Gamma}_{HK}=\Gamma_{HK} \rtimes \left\langle \iota^{\gamma}\right\rangle$.
%\end{lem}

%\begin{proof} First rewrite Equation (\ref{igamma}) as  $\iota^{\gamma}(f)=\tau_{\varphi^{-1}(g^{-2})}\left(f^{-1}\right)^{\tau_{\varphi^{-1}(g^{-1})}}$. Take $\omega\in \Gamma_{HK}$ where
%$$\omega=(\tau_{\psi^{-1}(r \varphi^{-1}(g^{2} ))}, \psi^{-1}\tau_{\varphi^{-1}(g^{-1})})$$ and observe that $\Theta=\omega \iota^{\gamma}$.
%\end{proof}

 

\begin{deff}
\label{novaacao} Given finite subgroups $H$ and $K$ of $G_1$ we set $$\overline{\Gamma}_{HK}:=\widetilde\Gamma_{HK}\cdot \left\langle \iota\right\rangle$$ if there exist $\psi \in \aut(G_1)$ such that $\psi(K)=H^{g_1}$ and $\psi(H)=K$, for some $g_1\in G_1$ and  put $$\overline{\Gamma}_{HK}:=\widetilde\Gamma_{HK}$$ otherwise. Further let $$\overline\Gamma_{HK}\backslash \Iso(H,K)$$ denote the set of orbits of $\Iso(H, K)$ modulo the action  of $\overline{\Gamma}_{HK}$.
\end{deff}

%\begin{rem}
%The Lemma \ref{carmo} shows that the previously defined orbits are independent of the automorphism $\psi$ chosen in the definition of $\iota$ (see formula (\ref{def iota abst})).    
%\end{rem}


%\begin{rem} \label{novaacao2} When $H=K$ then $\Iso(H,H)=\aut(H)$, $\overline{\inn}_{G_1}(H)=\bar{N}_{G_1}(H)$. In this case one can define  $\Gamma_{H}=\overline{\inn}_{G_1}(H)\rtimes\aut_{G_1}(H)= \bar{N}_{G_1}(H) \rtimes\aut_{G_1}(H)$, $\Gamma_{HH}=\Gamma_{H}$ as being the stabilizer of $H$ of  the action of $\Gamma_H$ on  $S$ (see (\ref{acaoemX})) and 
%$\overline{\Gamma}_{HH}=\Gamma_{H}\cdot C_2$, where $C_2= \left\langle \iota\right\rangle$ acts by inversion on $\aut(H)$. Note that if $\iota \in \Gamma_H$ then $\overline{\Gamma}_{HH}=\Gamma_{H},$ otherwise  $\overline{\Gamma}_{HH}=\Gamma_{H}\rtimes C_2.$ 
%In particular, $\left[HH\right]= \overline{\Gamma}_{HH} \backslash \aut(H)$ denote the set of orbits of $\aut(H)$ modulo the action  of $\overline{\Gamma}_{HH}$.
%\end{rem}

%%%%%%%%%%%%%%%%%%%%%%%%%%%%%%%%%%%%%%%%%%%%%%%%%%%%%%%%%%%%%%%%%%%%%%

\begin{rem} \label{novaacao2} When $H=K$ then $\Iso(H,H)=\aut(H)$. In this case one can define   $\Gamma_{HH}=\Gamma_{H}=N_{G_1}(H) \rtimes\aut_{G_1}(H)$ as being the stabilizer of $H$ of  the action of $\Gamma_H$ on  $S$ (see (\ref{acaoemX})) and 
$\overline{\Gamma}_{HH}=\widetilde\Gamma_{H}\cdot C_2$, where $C_2= \left\langle \iota\right\rangle$ acts by inversion on $\aut(H)$. Note that if $\iota \in \widetilde\Gamma_H$ then $\overline{\Gamma}_{HH}=\widetilde\Gamma_{H},$ otherwise  $\overline{\Gamma}_{HH}=\widetilde\Gamma_{H}\rtimes C_2.$ 
In particular, $\overline{\Gamma}_{HH} \backslash \aut(H)$ denote the set of orbits of $\aut(H)$ modulo the action  of $\overline{\Gamma}_{HH}$.
\end{rem}




%%%%%%%%%%%%%%%%%%%%%%%%%%%%%%%%%%%%%%%%%%%%%%%%%%%%%%%%%%%%%%%%%%%%%%






%$\overline{\Gamma}_{HH}=\Gamma_{H}\rtimes C_2$, where $C_2= \left\langle \iota\right\rangle$ acts by inversion on $\aut(H)$. %and $\Gamma_{HH}$ is the stabilizer of $H$ given by an action similar to (\ref{acaoemX}).


The result below is a consequence of  Propositions   \ref{nonnormalhnnabs} and the results of this subsection. They generalize Theorem 4.5 in \cite{BZ} to $\OE$-groups. %We divided it into two parts to avoid the density of information in a single result. 

%\begin{thm} \label{classnonnormalabs} Let $G=\hnn(G_1, H, K, f_1, t_1)$ and $B=\hnn(G_1, H, K, f_2, t_2)$ be abstract  HNN-extensions with  non-conjugated finite associated  subgroups in $G_1$ such that $G_1$ be an OE-group. For $f_i \in \Iso(H, K),$ $(i=1, 2),$ the corresponding HNN-extensions $G$ and $B$ are isomorphic iff both belong to the same orbit in $[HK]$.   
%\end{thm}

\begin{thm} \label{classnonnormalabs} Let $G_1$ be an $\OE$-group and $$G=\hnn(G_1, H, K, f, t),\  B=\hnn(G_1, H, K, f_1, t_1)$$ be  abstract HNN-extensions with finite associated  subgroups. Then $G$ and $B$ are isomorphic iff $f$ and $f_1$ belong to the same $\overline\Gamma_{HK}$-orbit of $\Iso(H, K)$.    
\end{thm}

\begin{proof} 

$(\Rightarrow)$ We have two possibilities: either $H$ and $K$ are conjugate in the base group $G_1$ or  not.  Moreover, in the first case we may assume that $H=K$ by Lemma  \ref{conjugados da normal}.

\medskip
Thus suppose first that $H=K$. By Proposition \ref{nonnormalhnnabs} there  exist $\varphi\in \aut_{G_1}(H)$ and $ g_1 \in N_{G_1}(H)$ such that

$$( f_1)^{\varphi^{-1}}= \tau_{g_1^{-1}} ( f)^{\epsilon} \iff \tau_{g_1}f_1^{\varphi^{-1}}=( f)^{\epsilon},$$
where $\epsilon \in \{\pm 1\}.$ %Note that $\psi^{-1} \in \aut_{G_1}(H),\,\,  \psi^{-1}(r) \in N_{G_1}(H)$ and $G \cong \hnn(G_1, H, f_1^{-1}, t_1^{-1}) $ by Remark \ref{f e f inverse}.
If $\epsilon=1$ take $(g_1,\varphi)\in \Gamma_{H}$ and  observe that $$(g_1,\varphi)\cdot f_1=\tau_{g_1}f_1^{\varphi^{-1}}= f.$$

If $\epsilon=-1$ take $\beta=(\varphi^{-1}(g_1^{-1}), \varphi^{-1})\in \Gamma_{H}$ and observe that  $\iota(f)=f^{-1}$ (just consider $\psi=id$ in the Definition of $\iota$ in (\ref{def iota abst})). In this way, we have 
$$(\beta, \iota) \cdot f = \tau_{\varphi^{-1}(g_{1}^{-1})} \left( f^{-1}\right)^{\varphi}   =f_1.$$ %Therefore $f_1$ and $f$ belong to the same orbit.
Therefore for any $\epsilon=\pm 1$ we have that $f_1$ and $f$ are in the same $\overline\Gamma_{HH}$-orbit.



\medskip

Now suppose
that $H$ and $K$ are not conjugate in $G_1$. By Proposition \ref{nonnormalhnnabs} $G \cong B$ if, and only if, the subcases i) or ii) of Proposition \ref{nonnormalhnnabs} are satisfied. 

 \textbf{Subcase i)}: There exist $\psi\in \aut(G_1)$ with $\psi(H)=K $ and $ \psi(K)=H^{g_1}$ for some $g_1 \in G_1$ such that $f_1^{\psi^{-1}}(k)= \tau_{g_1^{-1}} f^{- 1}(k), \forall k \in K$. By definition of  $\iota$ in (\ref{def iota abst}) we have

$$\iota(f)=\tau_{\psi^{-1}(g_1^{-1})}(f^{-1})^{\psi}=f_1.$$
Therefore $f_1$ and $f$ belong to the same  $\overline{\Gamma}_{HK}$-orbit of $\Iso(H, K)$.

 \textbf{Subcase ii)}: There exist $\psi\in \aut(G_1)$ with $\psi(H)=H $ and $ \psi(K)=K^{g_1}$ for some $g_1 \in G_1$ such that $f_1^{\psi^{-1}}(h)= \tau_{g_1^{-1}} f(h), \forall h \in H$. Observe that $(g_1, \psi)\in \Gamma_{HK}$ and $$(g_1, \psi)\cdot f_1= \tau_{g_1}f_1^{\psi^{-1}}=f.$$
Therefore $f_1$ and $f$ belong to the same orbit.


 
\bigskip

\hspace{-0,8cm} $(\Leftarrow)$ It follows from Proposition   \ref{nonnormalhnnabs}. 
\end{proof}

\begin{cor}\label{number of isomorphisms} Let $G_1$ be an $\OE$-group and $H\cong K$ be fixed finite subgroups of $G_1$.
The number of isomorphism classes of abstract HNN-extensions of the form $\hnn(G_1, H, K, f, t)$  equals to the number of $\overline\Gamma_{HK}$-orbits 

\begin{equation}\label{formula numero hnn extensao}
|\overline{\Gamma}_{HK} \backslash \Iso(H, K)|
\end{equation}
in $\Iso(H, K)$.
\end{cor}


%\begin{rem} \label{Delta}
%By Theorem \ref{classnonnormalabs}  there is a bijection $\Delta$ from the set of isomorphism classes of the groups of the form $\hnn(G_1, H, K, f, t)$ and the set $[HK]$. We will use $\Delta$ in some parts of this text. %We will use $\Delta$ in the result below and elsewhere in this text. 
%\end{rem}

\subsection{The number of isomorphism classes of abstract HNN-extensions %with non-conjugate associated subgroups
} \label{cota bonita non conjugados}

 Our first goal is to count the isomorphism classes of HNN-extensions  $$\hnn(G_1, H, K, f, t)$$ with finite associated subgroups. 

%$$\textcolor{red}{G_1 \cdot L}=\inn(G_1)\cdot L=\{L^{g_1}| g_1 \in G_1\}.$$ 

%Let $S_{_{H}} \subset S$ the set of all finite subgroups of $G_1$ isomorphic to $H$. We have an action of $\inn(G_1)$ in $S_{_{H}}$  
%defined by $\tau_{g_1} \cdot L= L^{g_1^{-1}}$ for all $\tau_{g_1}\in \inn(G_1)$ and $L \in S_{_{H}}.$ 
Let $ S$ be the set of all finite subgroups of $G_1$. We have an action of $G_1$ on $S$ by conjugation. So, given $L\in S$ the orbit of $L$ with respect to this action (i.e., the conjugacy class of $L$ in $G_1$) is  %$\inn(G_1) \cdot L= \{L^{g_1}| g_1 \in G_1\}.$ 
$$\langle\langle L\rangle\rangle^{G_1}= \{L^{g_1}| g_1 \in G_1\}.$$ 
Denote by $\overline{S}$ the space of orbits of this action (the set of conjugacy classes of finite subgroups of $G_1$ ), i.e., $$\overline{S}= G_1\backslash S:=\{ \langle\langle L\rangle\rangle^{G_1} \,|\, L \in S\}.$$    

Let $[\overline{S}]^2$ be the collection of all unordered pairs of the form $\{X, Y\}$ such that $X, Y \in \overline{S}.$  Define an action of $\aut(G_1)$ on $[\overline{S}]^2$ by \begin{equation*} \varepsilon: \aut(G_1) \times [\overline{S}]^2 \longrightarrow [\overline{S}]^2, \end{equation*}
\begin{equation} \label{acao epsilon}
( \varphi, \{ \langle\langle H_1\rangle\rangle^{G_1} , \langle\langle K_1\rangle\rangle^{G_1} \} ) \mapsto \{ \langle\langle \varphi(H_1)\rangle\rangle^{G_1} , \langle\langle  \varphi(K_1)\rangle\rangle^{G_1}  \}.
\end{equation}

%\langle\langle L\rangle\rangle^{G_1}

%Given $H_1$ and $K_1$ finite subgroups of $G_1$ isomorphic to $H$ we represent the $\aut(G_1)$-orbit of $\{H_1, K_1\}$ by $$\aut(G_1)\{H_1, K_1\}=\{\{\varphi(H_1)^{g_1}, \varphi(K_1)^{g_2}\}|\varphi \in \aut(G_1), g_i \in G_1 (i=1,2)\}.$$
Note that $\{ \langle\langle H\rangle\rangle^{G_1}
, \langle\langle K\rangle\rangle^{G_1}
 \}$ and $\{ \langle\langle H_1\rangle\rangle^{G_1}
, \langle\langle K_1\rangle\rangle^{G_1}
 \}$ belong to the same $\aut(G_1)$-orbit of action (\ref{acao epsilon}) if and only if  there exist $\psi\in \aut(G_1)$ such that $\psi(H)=H_1$ and $\psi(K)=K_1^{g_1}$ for some $g_1\in G_1.$ 

Let $H_1, K_1 \in S,$ by Proposition  \ref{fixed hnn abstract},   for any $f\in \Iso(H,K)$ one has  $$\hnn(G_1, H, K, f) \cong \hnn(G_1, H_1, K_1, \tau_{g_1}f^{\psi^{-1}}).$$ 

%\begin{deff}
%We say that $\{H_1, K_1\}\in S_H\times S_H$ is a representative of the  $\aut(G_1)$-orbit of $\{\inn(G_1) \cdot H, \inn(G_1) \cdot K\}$  if there exist $\psi\in \aut(G_1)$ such that  $\psi(H)=H_1$ and $\psi(K)=K_1^{g_1}$ for some $g_1\in G_1$, i.e., $\{\inn(G_1) \cdot H, \inn(G_1) \cdot K\}$ and $\{\inn(G_1) \cdot H_1, \inn(G_1) \cdot K_1\}$ belong to the same $\aut(G_1)$-orbit according to action (\ref{acao epsilon}). 
%\end{deff}


%The map below is an action of $\aut(G_1)$ on the set of unordered  pairs of elements of $S_H$, say $S_H \times S_H,$ as follows  \begin{equation*} \varepsilon: \aut(G_1) \times \left( S_H \times S_H\right) \longrightarrow S_H \times S_H, \end{equation*}
%\begin{equation} \label{acao epsilon}
%( \varphi, \{H_1, K_1\} ) \mapsto \{\varphi(H_1), \varphi(K_1)\}.
%\end{equation} Since $\inn(G_1) \triangleleft \aut(G_1)$ we have that $\{\rho(H)\,|\, \rho \in \aut(G_1)\}= \{ \varphi(H) ^{g} |  g \in G_1 \mbox{ and } \varphi \in \aut(G_1)\}.$ 

%Given $H_1$ and $K_1$ finite subgroups of $G_1$ isomorphic to $H$ we represent the $\aut(G_1)$-orbit of $\{H_1, K_1\}$ according to  action (\ref{acao epsilon}) as  being the following \begin{equation} \label{orbita abstrata} \aut(G_1)\{H_1, K_1\}=\{\{\varphi(H_1)^{g_1}, \varphi(K_1)^{g_2}\}|\varphi \in \aut(G_1), g_i \in G_1 (i=1,2)\}.\end{equation}

%Note that $\{H, K\}$ and $\{H_1, K_1\}$ represent the same  $\aut(G_1)$-orbit if and only if $\{ \langle\langle H\rangle\rangle^{G_1}
%, \langle\langle K\rangle\rangle^{G_1}
% \}$ and $\{ \langle\langle H_1\rangle\rangle^{G_1}
%, \langle\langle K_1\rangle\rangle^{G_1}
 %\}$ are in the same $\aut(G_1 )$-orbit of action %(\ref{acao epsilon}). 

If $H$ and $K$ are finite subgroups of  $G_1,$ we saw in Corollary \ref{number of isomorphisms} that the number of isomorphism classes of abstract HNN-extensions of the form $\hnn(G_1, H, K, f, t),$ with $f \in \Iso(H, K)$ is 
$|\overline{\Gamma}_{HK} \backslash \Iso(H, K)|$. 

\medskip
Denote the set of orbits of action (\ref{acao epsilon}) as $\aut(G_1)\backslash [\overline{S}]^2$. From the preceding discussion and Theorem \ref{classnonnormalabs} we have  

\begin{thm} \label{general finite associated} Let $G_1$ be  an $\OE$-group. Then the number  of isomorphism classes of abstract HNN-extensions $$\hnn(G_1, H,K, f,t)$$ with finite associated subgroups  is $$\sum_{\{H, K\}} |\overline{\Gamma}_{HK} \backslash \Iso(H, K)|,$$ where $\{H, K\}$ ranges over representatives of the orbits $\aut(G_1)\backslash [\overline{S}]^2.$ \end{thm}



Let now $H$ be a finite subgroup of an OE-group $G_1.$ Let $S_{_{H}} \subset S$ be the set of all finite subgroups of $G_1$ isomorphic to $H$.
By Proposition \ref{fixed hnn abstract}   isomorphic HNN-extensions with finite associated subgroups must have isomorphic associated subgroups and by Proposition \ref{2a} isomorphic profinite HNN-extensions with finite associated subgroups must have isomorphic associated subgroups as well. 
Therefore for the calculation of genus it is important to count isomorphism classes of HNN-extensions with fixed associated subgroup $H$ up to isomorphism.

To do this we restrict  the action given by (\ref{acao epsilon}) of $\aut(G_1)$ on $[\overline{S}]^2,$ to the set $[\overline{S_H}]^2$ of all unordered pairs of conjugacy classes of finite  subgroups of $G_1$ isomorphic to $H$. Denote by $\aut(G_1) \backslash [\overline{S_H}]^2$ the space of orbits of this action. Thus we have the following

\begin{thm} \label{classe geral familia F} Let $H$ and $K$ be finite subgroups of an $\OE$-group $G_1$. Then the number $\eta$ of isomorphism classes of abstract HNN-extensions $$\hnn(G_1, H_1,\\
K_1, f_1,t_1)$$ such that %the subgroups $H_1$ and $K_1$ are non-conjugated in $G_1,$ 
$H_1 \cong H \cong K \cong  K_1$ and $f_1 \in \Iso(H_1, K_1)$ is $$\eta  = \sum_{\{H, K\}} |\overline{\Gamma}_{HK} \backslash \Iso(H, K)|,$$ where $\{H, K\}$ ranges over representatives of the orbits $\aut(G_1)\backslash [\overline{S}_H]^2.$
\end{thm} 

%Let $G_1$ be  an $\OE$-group. Then the number  of isomorphism classes of abstract HNN-extensions $$\hnn(G_1, H_,
%K, f,t)$$ with finite associated subgroups  is $$\sum_{\{H, K\}} |\overline{\Gamma}_{HK} \backslash \Iso(H, K)|,$$ where $\{H, K\}$ ranges over representatives of the orbits $\aut(G_1)\backslash [\overline{S}]^2,$  where $[\overline{S}]^2$ denotes the set of unordered pairs of conjugacy classes of isomorphic finite subgroups of $G_1$.


\subsection{Normal HNN-extension} \label{cota bonita conjugados}

 




%\begin{rem} \label{remark da contagem} 
%The formula that calculates the number of isomorphism classes of normal HNN-extensions of the form $\hnn(G_1, H, f,t),$ with $f\in \aut(H)$, is the same, that calculates the number of isomorphism classes of HNN-extensions  of the form $\hnn(G_1, H, K, f,t)$ with $H$ and $K$ conjugated in $G_1$ and $f\in \Iso(H,K)$.
%\end{rem}

\subsubsection{Fixed associated subgroups}

The formula (\ref{formula numero hnn extensao}) is a consequence of  Theorem \ref{classnonnormalabs} and calculates the number of isomorphism classes of HNN-extensions of the form $\hnn(G_1,H,K,f,t)$ with $H,K$ fixed finite subgroups of $G_1$, where $ f\in \Iso(H,K)$.
The objective of this subsection is to make more explicit  formula (\ref{formula numero hnn extensao})  by showing a more elegant version in the case where $H$ and $K$ are conjugate finite subgroups in $G_1$.
 %$G=\hnn(G_1,H,K,f,t)$ such that $H^{g_1=K$, for some $ g_1\in G_1$, where $f\in\Iso(H,K).$ 
 For this, we observe that if $G=\hnn(G_1,H,K,f,t)$ where $H^{g_1}=K$, for some $ g_1\in G_1$ then $G\cong \hnn(G_1, H, (\tau_{g_1}f)^{\pm 1}, t), \mbox{ where } (\tau_{g_1}f)^{\pm 1} \in \aut(H) \mbox{ and } g_1 \in G_1$ (see Lemma \ref{conjugados da normal}).  

\begin{deff}
An (profinite) abstract HNN-extension in which the stable letter normalizers associated subgroup will be called \textbf{normal}. 
%If the associated subgroups are not conjugate in the base group, we will say that the HNN-extension is \textbf{not normal}. 
\end{deff}

For normal HNN-extensions  it is possible to write the set of orbits of the Theorem \ref{classnonnormalabs} in a more natural way and this will be extremely useful to deduce several corollaries of  of this theorem.

From Proposition \ref{nonnormalhnnabs} we have that  $\overline{N}_{G_1}(H)$ acts on $\aut(H)$ by composition on the left: we will denote by $\overline{N}_{G_1}(H) \cdot f =\{\tau_g \circ f \,|\, g \in N_{G_1}(H), f \in \aut(H) \}$  the orbit of $f$ with respect to this action and the space of right cosets $\overline{N}_{G_1}(H) \backslash \aut(H)$ is the space of these orbits. Note that $\overline{N}_{G_1}(H) \triangleleft \overline{\aut}_{G_1}(H)$; furthermore, $ \overline{\aut}_{G_1}(H) \times C_2$ acts on the right on $\aut(H)$, where $C_2=\left\langle \iota\right\rangle$ acts by inversion, as follows 
%%%%%%%%%%%%%%%%%%%%%%%%%%
\begin{equation}\label{acao1} f \cdot (\psi, \iota) = \psi^{-1} f^{-1}\psi = (f^{-1})^{\psi}=\left(f^{\psi}\right)^{-1},\end{equation}
%%%%%%%%%%%%%%%%%%%%%%%%%%%%%%%%%%%%
where  $( \psi, \iota) \in  \overline{\aut}_{G_1}(H) \times C_2, f \in \aut(H).$ This induces a natural action on the right of the group $\overline{\aut}_{G_1}(H) \times C_2$ on the space of orbits $\overline{N}_{G_1}(H) \backslash \aut(H)$, given by
%%%%%%%%%%%%%%
\begin{equation} \label{acao2} \left( \overline{N}_{G_1}(H) \cdot f \right) \cdot (\psi, \iota) = \overline{N}_{G_1}(H) \cdot  (f^{-1})^{\psi}=\overline{N}_{G_1}(H) \cdot \left(f^{\psi}\right)^{-1},\end{equation}
%%%%%%%%%%%%%%%%%%%%%%
where $( \psi, \iota) \in  \overline{\aut}_{G_1}(H) \times C_2, f \in \aut(H).$ We will denote the space of orbits of the action given in (\ref{acao2}) as $\overline{N}_{G_1}(H) \backslash \aut(H)/ ( \overline{\aut}_{G_1}(H) \times C_2)$; this is not the space of double cosets as $C_2$ is not a subgroup of $\aut(H)$ and the action of  $\overline{\aut}_{G_1}(H)$ is by conjugation.

From  Formulas (\ref{acao1}) and (\ref{acao2}) it can be seen that the set of orbits of the action first on the left by $\overline{N}_{G_1}(H)$ in $\aut(H)$ followed by the action of $ \overline{\aut}_{G_1}(H) \times C_2$ on $\aut(H)$ on the right  is in bijection with $\overline{\Gamma}_{HH} \backslash \aut(H)$, so from now on we will identify $\overline{\Gamma}_{HH} \backslash \aut(H) = \overline{N}_{G_1}(H) \backslash \aut(H)/ ( \overline{\aut}_{G_1}(H) \times C_2)$. 

Thus by %Theorem \ref{classe geral familia F} 
Corollary \ref{number of isomorphisms} we have the following

\begin{thm} Let $G_1$ be an $\OE$-group and $H$ be a fixed finite subgroup of $G_1$.
Then there is a bijection between   the set of orbits \begin{equation} \label{a bijecao das orbitas} \overline{N}_{G_1}(H) \backslash \aut(H)/ ( \overline{\aut}_{G_1}(H) \times C_2) \end{equation} and the set of isomorphism classes of HNN-extensions of the form $\hnn(G_1, H, f, t),$ where $f\in \aut(H)$.\end{thm}

 Note that $\inn(H)\leq \overline{N}_{G_1}(H)$; so we can pass to $\out(H)$ using tilde instead of overline. 
Therefore
we can state the following
 



 \begin{cor}\label{corhnn} Let $H$ and $K$ be fixed conjugate finite subgroups of an $\OE$-group $G_1$. Then the number $n_{_{H}}$  
 of isomorphism classes of HNN-extensions $\hnn(G_1, H, K, f,t)$,  $f \in \Iso(H, K)$  %(with $H$ and $K$  fixed) 
 is $$n_{_{H}}=|\widetilde{N}_{G_1}(H) \backslash \out(H)/  ( \widetilde{\aut}_{G_1}(H) \times C_2) |.$$ 
 \end{cor}  

%%%%%%%%%%%%%%%REMARK DO DELTA QUE VAGNER COMENTOU DE POR DEPOIS DE 4.12 ABAIXO

%\begin{rem} \label{Delta}
%By Theorem \ref{classnonnormalabs}  there is a bijection $\Delta$ from the set of isomorphism classes of the groups of the form $\hnn(G_1, H, K, f, t)$ and the set $[HK]$. We will use $\Delta$ in the result below and elsewhere in this text. 
%\end{rem}


\begin{cor}\label{g1} %Suppose  $\widetilde{N}_{G_1}(H)$ is normal in $\out(H)$. 
Suppose that $\widetilde{\aut}_{G_1}(H)$ is central in $\out(H)$. Then $\widetilde{N}_{G_1}(H) \\ \unlhd \out(H)$ and the number of isomorphism classes of groups of the form $G=\hnn(G_1,H,f, t)$, with $f\in \aut(H)$ is equal to $\frac{n+d}{2}$, where $d$ is the number of elements of order $\leq 2$ in the quotient  $\out(H)/\widetilde{N}_{G_1}(H)$ and $n=|\out(H)/\widetilde{N}_{G_1}(H)|$.
\end{cor}
%\begin{cor}\label{g1} Suppose  $\widetilde{N}_{G_1}(H)$ is normal in $\out(H)$. Then the number of isomorphism classes of groups of the form $G=\hnn(G_1,H,f, t)$, with $f\in \aut(H)$ is equal to $\frac{n-d}{2}+d=\frac{n+d}{2}$, where $d$ is the number of elements of order $\leq 2$ in $\widetilde{N}_{G_1}(H)\backslash \out(H)$ and $n=| \widetilde{N}_{G_1}(H) \backslash \out(H)|$.
%\end{cor}

\begin{proof} Since $\widetilde{N}_{G_1}(H)\leq \widetilde{\aut}_{G_1}(H)$ and $\widetilde{\aut}_{G_1}(H)$ is central in $\out(H)$, we have $\widetilde{N}_{G_1}(H)^{\gamma}=\widetilde{N}_{G_1}(H)$ for all $\gamma \in \out(H)$, i.e $\widetilde{N}_{G_1}(H)\unlhd \out(H)$.   Denote by $\tilde{f_1}$ the image of $f_1 \in \aut(H)$ in $\out(H).$ 
By formula (\ref{a bijecao das orbitas}) combined with Remark \ref{novaacao2} and Corollary \ref{corhnn}, the isomorphism class $[B]$ of $B=\hnn(G_1, H, f_1), \mbox{ with } f_1\in \aut(H)$ corresponds to $$\left\{ \tilde{\tau}_{k}  \left(\tilde{f}_1\right)^{\pm 1} \ | \ k
\in  N_{G_1}(H) \right\}=[f_1],$$ where $[f_1]$ denotes the orbit of $f_1.$ Note that we can partition $\widetilde{N}_{G_1}(H) \backslash \out(H)$ as $\tilde{N}_{1} \dot\cup %\uplus 
\tilde{N}_{2},$ where $\tilde{N}_{1}$ is the set containing the orbits $\widetilde{N}_{G_1}(H) \cdot \tilde{f}_1$ such that $|\tilde{f}_1|> 2$ and $\tilde{N}_{2}$ is the set containing the orbits $\widetilde{N}_{G_1}(H) \cdot \tilde{f}_1$ such that $|\tilde{f}_1| \leq 2.$ Thus $$ (\out(H)/\widetilde{N}_{G_1}(H)) /  ( \widetilde{\aut}_{G_1}(H) \times C_2) =  \tilde{N}_1 /  ( \widetilde{\aut}_{G_1}(H) \times C_2) \dot \cup \tilde{N}_2 /  ( \widetilde{\aut}_{G_1}(H) \times C_2).$$ By (\ref{acao2}) and (\ref{a bijecao das orbitas}) we have that $|\tilde{N}_2 / ( \widetilde{\aut}_{G_1}(H) \times C_2)|=d$ and $|\tilde{N}_{1} / ( \widetilde{\aut}_{G_1}(H) \times C_2)|= \frac{n - d}{2}$ and so all together the number of orbits is $\frac{n - d}{2}+d=\frac{n + d}{2}$. 
\end{proof}



%\begin{rem} \label{so normal}
%If $\widetilde{N}_{G_1}(H) \unlhd \out(H),$ we have the same result as in the preceding corollary.    
%\end{rem}

\begin{rem}
We note that Corollary \ref{g1} applies when $H$ is a finite cyclic subgroup in an $\OE$-group $G_1$. 
\end{rem}

Recall that a subgroup $H$ of $G$ is {\it malnormal} if $H\cap H^g=\{1\}$ for every $g \in G \backslash H$.

%\begin{cor}\label{malnormal} If $H$ is malnormal in $G_1$ then $n=|\out(H)|$.\end{cor}

\begin{cor} \label{malnormal} 
If $\widetilde{\aut}_{G_1}(H)$ is central in $\out(H)$ and $\widetilde{N}_{G_1}(H)=1$ (in particular if $H$ is malnormal in $G_1$), then the number  of isomorphism classes of groups of the form $G = \hnn(G_1, H, f, t),$ with
$f \in \aut(H)$ is equal to $\frac{|\out(H)| +d}{2},$ where $d$ is the number of elements of order $\leq 2$ in $\out(H).$ 
\end{cor}




\subsubsection{Not fixed associated subgroups}

Let $L\leq G_1$ be a finite subgroup of $G_1$ isomorphic to $H$, i.e. $L \in S_H.$ Then by  Corollary \ref{corhnn} the number of isomorphism classes of abstract normal HNN-extensions  $\hnn(G_1, L, f, t)$ is  $$n_{_{L}}:=| \widetilde{N}_{G_1}(L) \backslash \out(L)/ ( \widetilde{\aut}_{G_1}(L) \times C_2) |.$$ Note that in our case $K$ belongs always to the same orbit as $H$ and so the action that we consider is $\aut(G_1)$  on $\overline S_H$ or simply on $S_H$. From Theorem \ref{classe geral familia F} then, we deduce then the following  

%We observe that this number does not depend on the choice of representative $W$ in $Aut(G_1)$-orbit  $Aut(G_1)L$, because if $\varphi\in Aut(G_1)$ is such that $\varphi(L)=W$ then we have $$HNN(G_1, L, f) \cong HNN(G_1, \varphi(L), \varphi f\varphi^{-1}) = HNN(G_1, W,  f^{\varphi^{-1}})$$ for all $f \in Aut(L)$ (see Proposition \ref{hnnabs}).

%Let $G=HNN(G_1, H, f, t)$ be an HNN-extension. We will say that the subgroup $H$ is not fixed if $f\in Inj(H, G_1)$, that is, when $f(H) \in S$. Then we deduce the following

\begin{thm} Let $H$  be  a finite subgroup of the $\OE$-group $G_1$. Then the number of isomorphism classes of abstract HNN-extensions $\hnn(G_1, L, K, f, t)$ with conjugate associated subgroups isomorphic to $H$ equals $$\sum_{L} n_{_{L}},$$ where $L$ ranges over representatives of the orbits  $\aut(G_1)\backslash S_H$.
\end{thm}




\begin{cor} The number of isomorphism classes of abstract HNN-extensions $\hnn(G_1, L, K, f, t)$ with conjugate $L$ and $K$   in $G_1$ and isomorphic to $H$ is less or  equal to $$|\aut(G_1)\backslash {S}_H| \cdot |\out(H)/C_2|.$$
\end{cor}

 \begin{rem} If $|\out(H)|$ is odd then  $$| \out(H)/ C_2|=\frac{|\out(H)|-1}{2} + 1.$$ Therefore, in this case the number of isomorphism classes in the previous Corollary is less than or equal to $$|\aut(G_1)\backslash S_H| \cdot \left( \frac{|\out(H)|-1}{2} + 1\right).$$
\end{rem}

%\begin{cor} Let $G= G_1*_H (H\times \Z)$ be a free product with finite amalgamation, where $G_1$ is an $\OE$-group. Then the number of isomorphism classes of such free products with amalgamation $G_1*_H (H\times \Z)$ with $H$ fixed is 
%$$n_{_{H}}:=| \widetilde{N}_{G_1}(H) \backslash \out(H)/ ( \widetilde{\aut}_{G_1}(H) \times C_2) |.$$
%Moreover, if $H$ is malnormal, then 
%$$n_{_{H}}:=| \out(H)/ ( \widetilde{\aut}_{G_1}(H) \times C_2) |.$$
%\end{cor}


%%%%%%%%%%%%%%%%%%%%%%%%%%%%%%%%%%%%%%%%%%%%%%%%%%%%%%%%%%%%%%%%%%%%%%

Now we deduce the formula for the number of isomorphism classes of all HNN-extensions with finite conjugate associated subgroups.

\begin{thm} Let $G_1$ be an $\OE$-group. Then the number of isomorphism classes of abstract HNN-extensions $\hnn(G_1, L, K, f, t)$ with finite conjugate associated subgroups  equals $$\sum_{L} n_{_{L}},$$ where $L$ ranges over representatives of the orbits  $\aut(G_1)\backslash S$.
\end{thm}

\begin{cor} The number of isomorphism classes of abstract HNN-extensions $\hnn(G_1, L, K, f, t)$ with $L$ and $K$ finite and conjugate  in $G_1$  is less or  equal to $$|\aut(G_1)\backslash {S}| \cdot |\out(H)/C_2|.$$
\end{cor}

\subsubsection{Examples}


\begin{exa} Let $G= G_1*_{H, \mu} (H\rtimes_{\alpha} \Z)$ be a free product with finite amalgamation, where $G_1$ is an $\OE$-group and  $\mu$ is a monomorphism from $H$ to $H\rtimes_{\alpha} \Z= H\rtimes_{\alpha} \left\langle t \right\rangle.$ Note that $G\cong \hnn(G_1, H, \gamma, t),$ where $\gamma = \alpha(t)\circ \mu \in \aut(H).$ This case was missing in the study of free products with finite amalgamation in \cite{BPZ2}. The results of this section fill this gap, in fact, the number of  isomorphism classes of such free products with amalgamation $ G_1*_{H, \mu} (H\rtimes_{\alpha} \Z)$  with $H$ fixed is 
$$n=| \widetilde{N}_{G_1}(H) \backslash \out(H)/ ( \widetilde{\aut}_{G_1}(H) \times C_2) |.$$
Moreover, if $H$ is malnormal, then 
$$n=| \out(H)/ ( \widetilde{\aut}_{G_1}(H) \times C_2) |.$$
\end{exa}





The following result is an adaptation of \cite[Example 4.8]{ BZ}.


\begin{exa}\label{bb}
Let $G_1$ be an $\OE$-group, $H\leq G_1$ such that $H$ is central cyclic of order $n$ in $G_1$. As $H$ is central in $G_1$ it follows that $\overline{N}_{G_1}(H)=1=\widetilde{N}_{G_1}(H)$. Note that 
$\aut(H)=\prod_{p| n} \aut(C_{p^m})=\prod_{p| n} \out(C_{p^m})=\out(H),$   where $p^m$ is the order of the $p$-Sylow subgroup of $H$. 
Note that
$$\aut(C_{p^m})=\left\{\begin{array}{l}
1, \mbox{  if } m=1 \mbox{ and } p=2,\\
C_{2}, \mbox{  if } m=2 \mbox{ and } p=2,\\
C_{2}\times C_{2^{m-2}}, \mbox{  if } m\geq3\mbox{ and } p=2,\\
C_{p-1}\times C_{p^{m-1}}, \mbox{  if } p \mbox{ odd.}
\end{array} \right.$$
Thus $\aut(H)$ is abelian, and moreover  $\aut(C_{p^m})$ contains $2$ elements of order $\leq 2$ if $p$ is odd and $4$ elements if  $p=2$ and $m\geq3$. Let $r \geq 0$ be the number of primes $p \neq 2$ dividing $n.$ Then by Corollary \ref{g1}  the number of isomorphism classes of the groups $G=\hnn(G_1, H, f, t)$, where $f \in \aut(H)$ equals to:
$$\left\{\begin{array}{l} 

\frac{\prod_{2\neq p|n} \left( (p-1)p^{m-1}\right) -2^r}{2} +2^r, \mbox{ if the 2-Sylow subgroup of } H \mbox { has order } \leq 2;\\
\\

\frac{2 \cdot  \prod_{2\neq p|n} \left( (p-1)p^{m-1}\right) -2^{r+1}}{2} +2^{r+1}, \mbox{ if the 2-Sylow subgroup of } H \mbox{ has  order } 4;\\
\\
\frac{2^{m-1} \cdot  \prod_{2\neq p|n} \left( (p-1)p^{m-1}\right) -2^{r+2}}{2} +2^{r+2}, \mbox{ if the 2-Sylow subgroup of } H \mbox{ has order } 2^m \\ \mbox{ with } m\geq 3.
\end{array} \right.$$

In particular, if $H\cong C_{p^m}$, then the number of isomorphism classes is equal to 
$$\left\{\begin{array}{l}
1, \mbox{  if } m=1 \mbox{ and } p=2,\\
2, \mbox{  if } m=2 \mbox{ and } p=2,\\
(2^{m-1}-4)/2 +4, \mbox{  if  }m\geq3 \mbox{ and } p=2, \\
\frac{(p-1)p^{m-1}-2}{2} +2,\mbox{ if } p \mbox{ odd}.
\end{array} \right.$$
\end{exa}

%\begin{exa}
%Let $G_1$ be an $\OE$-group and $H, K$ isomorphic finite  subgroups such that $H$ is normal subgroup in $G_1$ and $K$ is not. Consider the following class of groups $\{G=\hnn(G_1, H, K, f), \mbox{ where } f\in \Iso(H, K)\,\}$. Then by Theorem  \ref{classnonnormalabs} the number of isomorphism classes of these groups equals to the cardinality of the orbit set $\overline{\Gamma}_{HK} \backslash \Iso(H, K).$
%\end{exa}

\begin{rem} \label{finite-by-cyclic} 
 
Let $G=\hnn(G_1, H, K, f, t)$ be an HNN-extension with finite associated subgroups such that $G_1=M_1\rtimes \Z,$ where $M_1$ is a finite subgroup, distinct from $H$ and $K.$ % where $M_1$  finite subgroup $M_1$ and distinct from $H$ and $K.$
%Let $G=\hnn(G_1, H, K, f, t)$ be an HNN-extension with finite associated subgroups such that $G_1=M_1\rtimes \Z$ with $ H \neq M_1 \neq K$ finite.  %such that $M_1\neq H$ %and $B=\hnn(B_1, Y, W, g, z)$ a  finitely generated residually finite group with finite associated subgroups.  
Then $G_1$ is 2-ended and hence is not an $\OE$-group. %However, Propositions \ref{fixed hnn abstract} and \ref{nonnormalhnnabs} also holds in this case. 
However, the results obtained in the previous sections for this $G_1$ remain valid.\end{rem}

\section{Profinite HNN-extensions} \label{secao 4}
\label{profinitehnnextensions}


In this section we will solve the isomorphism problem for profinite HNN-extensions.  % and we will give some quotas, especially in the case where the associated subgroups are conjugate in the base group (in particular for normal HNN-extensions). 

%\subsection{Isomorphism between profinite HNN-extensions} \label{normal profinite HNN}
%The result below generalizes Corollary 2.5 in Bessas-Zalesskii \cite{BZ} and will be very important in the calculations of the  genus  in the following sections. 
%\begin{cor}
%Let $G_1, B_1$ be OE-groups, $H$ and $H_1$ finite  subgroups of $G_1$ and $B_1$ respectively. Let $G= HNN(G_1, H,  f)$ and $B=HNN(B_1, H_1, f_1)$ be HNN-extensions such that in the first the associated subgroups are conjugated in $G_1$ and in the second they are non-conjugated in $B_1.$ Then $\widehat G \not \cong \widehat B.$    
%\end{cor}
Let $G= \hnn(G_1, H, K, f, t)$ be a profinite HNN-extension. We define the set \begin{equation} \label{N mais}  N^{+}:=\{ gt^{- 1} \in N_{G}(H) \hspace{0.1cm} | \hspace{0.1cm}  \overline{\left\langle g, G_1 \right\rangle}=G\}.\end{equation} 
Consider $g_1\in N_{G_1}(H)$ and $g=g_1t$. Then we have $g_1=gt^{-1}\in N^{+}.$ Therefore \begin{equation} \label{NG1H em N+} N_{G_1}(H) \subset N^{+} \subset N_{G}(H) .\end{equation}


 %%%%%%%%%%%%%%%%%%%%%%%%%%%%%%%%%%
 Further denote by $\overline{N}^{+}$ and $\widetilde{N}^{+}$  their natural images in $\aut(H)$ and $\out(H),$ 
respectively. Note that it is $N^+$ that marks the difference between the abstract and the profinite case, as we will see later, it will be crucial for the calculation of the genus.  

The following proposition will be fundamental to find a bound for the number of isomorphism classes of profinite HNN-extensions. % The following proposition will be fundamental to find the number of isomorphism classes of profinite HNN-extensions. 
 
 %\begin{prop}\label{refhnn}
%Let $H$ be a finite subgroup of a profinite $OE$-group $G_1$. Let  $G= HNN(G_1, H, t, f)$ and $B=HNN(G_1, H, t_1, f_1)$ be profinite HNN-extensions  with $ f, f_1\in Aut(H)$. Then $G\cong B$ if and only if there exists $\psi_1\in Aut_{G_1}(H)$ and $k\in N^{+}$ such that  $\psi_1f_1\psi_1^{-1}= f^{\epsilon}\tau_k$ where $k=gt^{-\epsilon}$, $g\in N_{G}(H)$ and $\epsilon= \pm 1$.
%\end{prop}






\begin{prop}\label{nonnormalhnnprof} Let $$G= \hnn(G_1, H, K, f_1, t_1)\  {\rm and}\  B=\hnn(G_1, H, K, f_2, t_2)$$ be profinite HNN-extensions with associated finite subgroups in a profinite $\OE$-group $G_1$. Suppose that $K=H$ or $H$ and $K$ are not conjugate in $G_1$. Then $G\cong B$ if and only if 

\begin{enumerate}
\item [(i)] there exist $\varphi\in \aut(G_1)$ with $\varphi(H)=K, \varphi(K)=H^{g_1}$ for some $g_1 \in G_1$ such that $ f_2^{\varphi^{-1}} = \tau_{g_1^{-1}} \cdot  (f_1\tau_{n^{-1}})^{-1},$ with  $n\in N^{+}$;
\item [(ii)] there exist $\varphi\in \aut(G_1)$ with $\varphi(H)=H, \varphi(K)=K^{g_1}$ for some $g_1 \in G_1$ such that $f_2^{\varphi^{-1}} = \tau_{g_1^{-1}} \cdot f_1 \cdot \tau_{n^{-1}} ,$ with $n\in N^{+}$.
\end{enumerate}
\end{prop}

\begin{proof}
%$(\Leftarrow)$ The proof is analogous to what was done in the abstract case (see Proposition \ref{nonnormalhnnabs}). 

$(\Longrightarrow)$ By Proposition \ref{2a} there exists a continuous  isomorphism \break 
$\psi:B\longrightarrow G$ such that $\psi(G_1) = G_1$, $\psi(H)=K$ and $\psi(K)=H^{g_1}$ or  $\psi(G_1) = G_1$, $\psi(H)=H$ and $\psi(K)=K^{g_1}$ for some $g_1\in G_1$.  


First suppose that $H=K$ and observe that then $g_1\in N_{G_1}(H)$. Put $g:=\psi(t_2) \in G$. Since $t_2 \in N_{B}(H)$ we conclude that $gg_1^{-1} \in N_{G}(H)$ and $ 
 \overline{\left\langle  G_1, (gg_1^{-1})^{\epsilon} \right\rangle}=G$, where  $\epsilon \in \{\pm 1\}.$  %, moreover, $g \in N^{+}$. 
Put $n:= (gg_1^{-1})^{\epsilon}t_1^{-1} \in N^{+}$ and observe that for all $h \in H$ one has $$\psi(f_2(h)) = \psi(h)^{\psi(t_2)} = \psi(h)^{(nt_1)^{\epsilon}g_1} =\tau_{g_1^{-1}}(f_1\tau_{n^{-1}})^{\epsilon} \psi(h), \mbox{ where } \, \epsilon = \pm 1.$$ 
%$$\psi(f_1(h)) = \psi(h)^{\psi(t_1)} = \left(\left(\psi(h)\right)^{k}\right)^{t^{\epsilon}} =f^{\epsilon}\left(  \tau_{k} \psi(h)\right).$$
Take $\varphi:=\psi|_{_{G_1}} \in \aut_{G_1}(H)$. Then $\varphi f_2\varphi^{-1}=  \tau_{g_1^{-1}}(f_1^{-1}\tau_{n^{-1}})^{\epsilon}$ as required. 

\bigskip
If $H$ and $K$ not conjugate in $G_1$ let's divide it into two cases. 

\textbf{Case 1:} Suppose  $\psi(G_1) = G_1$, $\psi(H)=K$ and $\psi(K)=H^{g_1}$ . 
Put $\psi(t_2) = g \in G$. Note that $ H^{g_1}=\psi(K) =\psi(H^{t_2})=\psi(H)^{\psi(t_2)}=K^g=H^{t_1g}$, from where it follows that $n:=g_1 g^{-1} t_{1}^{-1} \in N_G(H)$, furthermore $n\in N^{+}$. % Set $\rho^{-1}:=n^{-1}g_1 \in G,$ thus $\varphi(t_2)=g=t_{1}^{-1}n^{-1}g_1=t_1^{-1}\rho^{-1}$. 
Then for all $h \in H,$ we have$$\psi(f_2(h))=\psi(h^{t_2})=\psi(h)^{\psi(t_2)}=\psi(h)^{t_{1}^{- 1}n^{-1}g_1}=(\tau_{g_1^{-1}}\tau_{n} f_1^{- 1}  \psi)(h).$$ Therefore taking $\varphi = \psi|_{_{G_1}} \in \aut(G_1)$ we have  $f_2^{\psi^{-1}} = f_2^{\varphi^{-1}} = \tau_{g_1^{-1}} \cdot \tau_{n} \cdot f_1^{-1}= \tau_{g_1^{-1}} \cdot   (f_1\tau_{n^{-1}})^{-1},$ with $g_1 \in G_1$ and $n\in N^{+}$. 

\textbf{Case 2:} Suppose $\psi(G_1) = G_1$, $\psi(H)=H$ and $\psi(K)=K^{g_1}$. %It follows analogously to Case 1 that  %$\psi(K) \leq \psi(G_1)^{\psi(t_2)}\cap \psi(G_1)=G_1^{\psi(t_2)} \cap G_1 \leq H^{g_1} \mbox{ or } K^{g_1}$ for some $g_1 \in G_1$  (since $\psi$ is
%an isomorphism and $H$ and $K$ are finite non-conjugate subgroups of $G_1$) we have 
%$\psi(K)=K^{g_1}$ for some $g_1 \in G_1$. 
Take $\psi(t_2)=g\in G$ and $n:=gg_1^{-1}t_1^{-1} \in N^{+}$. 
Then for all $h\in H$, we have $$\psi(f_2(h))=\psi(h^{t_2})=\psi(h)^{\psi(t_2)}=\psi(h)^{ n t_1g_1}=(\tau_{g_1^{-1}} f_1 \tau_{n^{-1}} \psi)(h).$$ 
So  taking $\varphi=\psi|_{_{G_1}}$ the result follows. 

\bigskip
$(\Longleftarrow)$ Define a map $\psi$  from $B$ to $G$ as follows $$\psi: \left\{
\begin{array}{ccl}
g_1 & \longmapsto  & \varphi(g_1), \forall  g_1 \in  G_1,\\
t_2 & \longmapsto  & (nt_1)^{\epsilon}g_1, \epsilon=\pm 1. \\
\end{array}
\right.$$
If (i) holds, consider $\epsilon=-1$ and if  (ii) holds, consider $\epsilon=1$;  in both cases $\psi$ is an isomorphism.

\end{proof} 


\subsection{Isomorphism problem for profinite HNN-extensions} \label{iso profinite HNN}


In this section we will give a bound for the number of isomorphism classes of profinite HNN-extensions with base group $G_1$ being a profinite $\OE$-group and fixed finite associated  subgroups. In the  case when the HNN-extensions have the conjugate associated subgroups  in the base group, some corollaries will also be shown.

Let $G=\hnn(G_1, H, K, f, t)$ be a profinite HNN-extension, where $H, K$ are finite associated subgroups of a profinite OE-group $G_1$. 

  Consider the closed subgroup 
 \begin{equation} \label{gamma H2} 
 \Gamma_{\widehat H}:=G_1 \rtimes \aut_{G_1}(H)\end{equation} of the holomorph $G_1\rtimes \aut(G_1)$ and as in the abstract case sometimes we shall write $\tau_{g_1}$ for the inner automorphism corresponding to an element $g_1$ of $G_1$. 

%Define on \begin{equation} \label{gammachapeuH} \Gamma_{\widehat{H}}:=G_1 \times \aut_{G_1}(H),\end{equation} 
 %of the holomorph $G_1\rtimes \aut(G_1)$. 
Denote by $P=\mbox{Subgr}(G_1)$ the set of all closed subgroups of $G_1.$ Note that $\mbox{Subgr}(G_1)$ is a profinite space because if $G_1=\varprojlim_{U\triangleleft_o G}  G_1/U$ then $\mbox{Subgr}(G_1) = \varprojlim_{U\triangleleft_o G}  \mbox{Subgr}(G_1/U)$. 
There is a continuous action of $$\Gamma_{\widehat{H}}=G_1\rtimes \aut_{G_1}(H)$$ on the set $P$ namely \begin{equation} \label{acaoemXp} (g_1,\alpha)\cdot X= \tau_{g_1}(\alpha(X))=\left[\alpha(X)\right]^{g_1^{-1}} \in P, \mbox{ where } (g_1,\alpha) \in \Gamma_{\widehat{H}}, X \in P. \end{equation}
 Let $\Gamma_{\widehat{HK}}$ be the stabilizer of $K$ in $\Gamma_{\widehat{H}}$ for this action on $P$ (see (\ref{acaoemXp})). Then $\Gamma_{\widehat{HK}}$ is a profinite group (see 5.6.3 in \cite{RZ}).
 
 
 
 Now let's fix a finite subgroup $H\leq G_1$ and consider 
$\widehat{\Upsilon} :=\displaystyle\bigcup_{X \leq G_1, |X|<\infty} \Iso(H,X)$.  Similarly to the abstract case, $\Gamma_{\widehat{H}}$ acts on $\widehat{\Upsilon}$.

%$\overline{\Gamma}_{\widehat{HK}}$

\begin{lem}\label{elementaryp}
The semidirect product  $\Gamma_{\widehat{H}}=G_1 \rtimes \aut_{G_1}(H)$
acts on the left upon $\widehat{\Upsilon}$ by
\begin{equation} \label{acaosemidiretap} (g_1,\alpha) \cdot f=\tau_{g_1} \alpha f \alpha^{-1} = \tau_{g_1} f^{\alpha^{-1}}, \mbox{ where } (g_1,\alpha) \in \Gamma_{\widehat{H}} \mbox{ 
 and } f \in \widehat\Upsilon.
 \end{equation}
\end{lem}
\begin{proof} It is word by word the same as of Lemma \ref{elementary}. 
\end{proof}
 
 
  For $f\in \Iso(H, K)$, we have $(g_1,\alpha) \cdot f \in \Iso\left(H, [\alpha(K)]^{g_1^{-1}}\right) \subset \widehat \Upsilon$ and
the group $\Gamma_{\widehat{HK}}$ leaves $\Iso(H, K)$ invariant with respect  to the action of $\Gamma_{\widehat{H}}$ on $\widehat\Upsilon$ (see (\ref{acaosemidiretap})).  

From now on, let $$G=\hnn(G_1, H, K, f_1, t_1) \mbox{ and } B=\hnn(G_1, H, K, f_2, t_2),$$ be profinite HHN-extensions, where $G_1$ is a profinite OE-group and \break $f_i \in \Iso(H, K)\,(i=1,2)$. Our goal is to describe  when $G$ and $B$ are isomorphic. 


Suppose there exist $\psi \in \aut(G_1)$ such that $\psi(K)=H^{g_1}$ and $\psi(H)=K$, for some $g_1\in G_1$.  Define $\iota: \Iso(H,K) \longrightarrow \Iso(H,K)$ by \begin{equation} \label{def iota prof} \iota(f)=\tau_{\psi^{-1}(g_1^{-1})}(f^{-1})^{\psi}=\tau_{\psi^{-1}(g_1^{-1})} \left(f^{\psi}\right)^{-1}, f \in \Iso(H,K).\end{equation} 



If $H=K$ we can choose $\psi=id$ and $g_1=1,$ so $\iota$ sends $f$ to $f^{-1}.$ In this situation we have an isomorphism between $\hnn(G_1,H,f)$ and $\hnn(G_1,H,f^{-1})$ (see Lemma \ref{conjugados da normal}). Note that $\iota$ will be a bijection from $\Iso(H, H)=\aut(H)$ to $\aut(H)$. More generally, $\iota$ belongs to the group of symmetries $Sym(\Iso(H, K))$. Since $\Gamma_{\widehat{HK}}$ acts on $\Iso(H, K),$ we  denote by  $\widetilde\Gamma_{\widehat{HK}}$ the natural image of $\Gamma_{\widehat{HK}}$ in $Sym(\Iso(H, K))$ induced by the natural homomorphism $\Gamma_{\widehat{HK}} \longrightarrow Sym(\Iso(H, K)).$ The proof of the following lemma is similar to the proof of Lemma 
\ref{semidiretocomC2}. 

\begin{lem} \label{semidiretocomC2p}  Suppose there exist $\psi \in \aut(G_1)$ that swaps the conjugacy classes of $H$ and $K$.  Then $\iota $ normalizes $\widetilde\Gamma_{\widehat{HK}}$ and so $\langle \widetilde\Gamma_{\widehat{HK}}, \iota  \rangle=\widetilde\Gamma_{\widehat{HK}} \cdot  \left\langle  \iota \right\rangle$ is a subgroup of $Sym(\Iso(H, K))$.   

\end{lem}  


The next Lemma is proved similarly as Lemma \ref{carmo}

\begin{lem}\label{carmoprof} The subgroup $\widetilde\Gamma_{\widehat{HK}} \cdot  \left\langle  \iota \right\rangle$ of $Sym(\Iso(H, K))$ does not depend upon the choice of $\psi$.
\end{lem}

%As the orbits of the action of $\Gamma_{\widehat{HK}}$ on $\Iso(H, K)$ are the same orbits of the induced action of $\widetilde\Gamma_{\widehat{HK}}$ on $\Iso(H, K)$, the definition below follows. 


\begin{deff}
\label{novaacaopp} Given finite subgroups $H$ and $K$ of $G_1$ we set $$\overline{\Gamma}_{\widehat{HK}}:=\widetilde{\Gamma}_{\widehat{HK}}\cdot \left\langle \iota\right\rangle$$ if there exist $\psi \in \aut(G_1)$ such that $\psi(K)=H^{g_1}$ and $\psi(H)=K$, for some $g_1\in G_1$, and put $$\overline{\Gamma}_{\widehat{HK}}:=\widetilde{\Gamma}_{\widehat{HK}},$$ otherwise. Further let $$\overline{\Gamma}_{\widehat{HK}} \backslash \Iso(H, K)$$ denote the set of orbits of $\Iso(H, K)$ modulo the action  of $\overline{\Gamma}_{\widehat{HK}}$.
\end{deff}

\begin{rem} \label{novaacao2pp} When $H=K$ then $\Iso(H,H)=\aut(H)$. In this case one can define  $\Gamma_{\widehat{H}}=N_{G_1}(H) \rtimes\aut_{G_1}(H)$, $\Gamma_{\widehat{HH}}=\Gamma_{\widehat{H}}$ as being the stabilizer of $H$ of the action of $\Gamma_{\widehat{H}}$ on  $P$ (see (\ref{acaoemXp})) and 
$\overline{\Gamma}_{\widehat{HH}}=\widetilde\Gamma_{\widehat{H}}\cdot C_2$, where $C_2= \left\langle \iota\right\rangle$ acts by inversion on $\aut(H)$. Note that if $\iota \in \widetilde\Gamma_{\widehat{H}}$ then $\overline{\Gamma}_{\widehat{HH}}=\widetilde{\Gamma}_{\widehat{H}},$ otherwise $\overline{\Gamma}_{\widehat{HH}}=\widetilde{\Gamma}_{\widehat{H}} \rtimes C_2.$ 
In particular, $ \overline{\Gamma}_{\widehat{HH}} \backslash \aut(H)$ will denote the set of $ \overline{\Gamma}_{\widehat{HH}}$-orbits of $\aut(H)$.
\end{rem}

The result below is a consequence of Propositions   \ref{nonnormalhnnprof} and the discussion of this subsection. In the profinite case we can only give an estimate for the number of isomorphism classes of profinite HNN-extensions with fixed associated subgroups, as follows.  \begin{thm}\label{profinite isomorphism number} Let $H, K$ be fixed finite subgroups of a profinite $\OE$-group $G_1$. Then the number of isomorphism classes of profinite HNN-extensions \\$\hnn(G_1, H, K, f, t)$ is less than or equal to $$|\overline{\Gamma}_{\widehat{HK}} \backslash \Iso(H, K)|.$$\end{thm}

\subsection{Bounds for the number of isomorphism classes of profinite HNN-extensions with not fixed associated subgroups } \label{quotasprofnonconjugates}

Let $H$ be a finite subgroup of a profinite OE-group $G_1.$ Our goal is to find a bound for the number of isomorphism classes of profinite HNN-extensions of  $G_1$ with associated subgroups isomorphic to $H$. This will be important for the calculation of the genus since by Proposition \ref{profinito normal iso} if $H\not\cong H_1$ then $\hnn(G_1, H, K, f, t)\not\cong \hnn(G_1, H_1, K_1, f_1, t_1)$. 


Our considerations in this subsection  will be give analogous to what was done in the abstract case (see Subsection \ref{cota bonita non conjugados}). %Let $H$ and $K$ be finite isomorphic subgroups of a profinite $\OE$-group  $G_1$ 
Let $P_H \subset P$ the set of all finite subgroups of $G_1$ isomorphic to $H$. We have an action of $G_1$ on $P_H$  by conjugation. So, given $L\in P_H$ the orbit of $L$ with respect to this action (i.e., the conjugacy class of $L$ in $G_1$) is $$\langle\langle L\rangle\rangle^{G_1}= \{L^{g_1}| g_1 \in G_1\}.$$ Put $$\overline{P}_H= G_1\backslash P_H:=\{ \langle\langle L\rangle\rangle^{G_1} \,|\, L \in P_H\}.$$   

Let $[\overline{P}_H]^2$ be the collection of all unordered pairs of the form $\{X, Y\}$ such that $X, Y \in \overline{P}_H.$ Define similarly to what was done in the abstract case (see Subsection \ref{cota bonita non conjugados}) an action of $\aut(G_1)$ on $[\overline{P}_H]^2$ as follows  

\begin{equation*} \rho: \aut(G_1) \times [\overline{P}_H]^2 \longrightarrow [\overline{P}_H]^2, \end{equation*}  %\begin{equation} \label{acao p}
%( \varphi, \{\inn(G_1) \cdot H_1, \inn(G_1) \cdot K_1\} ) \mapsto \{\inn(G_1) \cdot \varphi(H_1),  \inn(G_1) \cdot \varphi(K_1)\}.
%\end{equation}
\begin{equation} \label{acao p}
( \varphi, \{ \langle\langle H_1\rangle\rangle^{G_1} , \langle\langle K_1\rangle\rangle^{G_1} \} ) \mapsto \{ \langle\langle \varphi(H_1)\rangle\rangle^{G_1} , \langle\langle  \varphi(K_1)\rangle\rangle^{G_1}  \}.
\end{equation}
%The map below is an action of $\aut(G_1)$ on the set of unordered  pairs of elements of $P_H$, say $P_H \times P_H,$ as follows  \begin{equation*}   \rho: \aut(G_1) \times \left( P_H \times P_H\right) \longrightarrow P_H \times P_H, \end{equation*} \begin{equation} \label{rho}
%( \varphi, \{H_1, K_1\} ) \mapsto \{\varphi(H_1), \varphi(K_1)\}.
%\end{equation}

Consider  finite subgroups $H_1$ and $K_1$ of $G_1$ isomorphic to $H$ and $K,$ respectively.  Note that $\{ \langle\langle H\rangle\rangle^{G_1}
, \langle\langle K\rangle\rangle^{G_1}
 \}$ and $\{ \langle\langle H_1\rangle\rangle^{G_1}
, \langle\langle K_1\rangle\rangle^{G_1}
 \}$ belong to the same $\aut(G_1)$-orbit according to action (\ref{acao p}) if and only if  there exist $\psi\in \aut(G_1)$ such that $\psi(H)=H_1$ and $\psi(K)=K_1^{g_1}$ for some $g_1\in G_1.$ 

%\begin{deff}
%We say that $\{H_1, K_1\}\in P_H\times P_H$ is a representative of the  $\aut(G_1)$-orbit of $\{\inn(G_1) \cdot H, \inn(G_1) \cdot K\}$  if there exist $\psi\in \aut(G_1)$ such that  $\psi(H)=H_1$ and $\psi(K)=K_1^{g_1}$ for some $g_1\in G_1$, i.e., $\{\inn(G_1) \cdot H, \inn(G_1) \cdot K\}$ and $\{\inn(G_1) \cdot H_1, \inn(G_1) \cdot K_1\}$ belong to the same $\aut(G_1)$-orbit according to action (\ref{acao p}). 
%\end{deff}

 By Proposition  \ref{2a}, for each $f\in \Iso(H,K)$ we have $$\hnn(G_1, H, K, f,t) \cong \hnn(G_1, H_1, K_1, \tau_{g_1}f^{\psi^{-1}},t_1).$$


When $H$ and $K$ are fixed finite subgroups of a profinite OE-group $G_1,$ we saw from Theorem \ref{profinite isomorphism number} that the number $n_{_{\widehat{HK}}}$ of isomorphism classes of profinite HNN-extensions of the form $\hnn(G_1, H, K, f, t),$ with $f \in \Iso(H, K)$ satisfies the following $$n_{_{\widehat{HK}}} \leq | \overline{\Gamma}_{\widehat{HK}} \backslash \Iso(H, K)|.$$ 

Denote the set of orbits of action (\ref{acao p}) by $\aut(G_1)\backslash [\overline{P}_H]^2.$  % From the preceding discussion and Proposition \ref{normal non iso nao normal} 


\begin{thm}\label{profinite problem} Let $H$ and $K$ be finite subgroups of a profinite  $\OE$-group $G_1$. Then the number $\hat \eta$ of isomorphism classes of profinite HNN-extensions $\hnn(G_1, H_1, K_1, f_1)$ such that %the subgroups $H_1$ and $K_1$ are non-conjugated in $G_1,$ 
$H_1 \cong H \cong K \cong  K_1$ and $f_1 \in \Iso(H_1, K_1)$ satisfies the following inequalities $$\hat \eta\leq \sum_{\{H, K\}} | \overline{\Gamma}_{\widehat{HK}} \backslash \Iso(H, K)|,$$ where $\{H, K\}$ ranges over  representatives of the $\aut(G_1)$-orbits of $[\overline{P}_H]^2.$
\end{thm} 


Note that the set of all finite sugroups of a profinite group is not closed in general in the space of subgroups of it. However, the subset of subgroups of order bounded by a natural number $n$ is closed. This allows to extend the action $\rho$ given by (\ref{acao p}) of $\aut(G_1)$ on $[\overline{P}_H]^2,$ to the set $[\overline{P_n}]^2$ of all unordered pairs of conjugacy classes of  subgroups of $G_1$ of order $\leq n$. Denote by $\aut(G_1) \backslash [\overline{P_n}]^2$ the space of orbits of this action. Then we can extend the preceding theorem by establishing a bound on the number of isomorphism classes of all profinite HNN-extensions of $G_1$ with finite associated subgroups.

\begin{thm}\label{profinite HNN-extensions with finite associated} Let $G_1$ be a profinite  $\OE$-group. Then the number  of isomorphism classes of profinite HNN-extensions $\hnn(G_1, H, K, f, t)$ with finite associated subgroups does not exceed  $$ \sum_{\{H, K\}} | \overline{\Gamma}_{\widehat{HK}} \backslash \Iso(H, K)|,$$ where $\{H, K\}$ ranges over representatives of the orbits $\aut(G_1)\backslash [\overline{P_n}]^2.$ 
%where $\{H, K\}$ ranges over  representatives of the $\aut(G_1)$-orbits of unordered pair of finite subgroups of $G_1$.
\end{thm} 

\begin{rem} We can extend the action of $Aut(G_1)$ to the the set of all finite sugroups of $G_1$ and get the statement as written in Theorem \ref{int general finite associated}, but it differs from the theorem above only when the number of isomorphisms is infinite.\end{rem}


\subsubsection{ Fixed conjugate associated subgroups } \label{quotasprofconjugates}

The objective of this subsection is to find  bounds for  the number of isomorphism classes of profinite HNN-extensions of the form $G=\hnn(G_1, H, K, f)$ with $H,K$ conjugate in $G_1$. For this, we observe that $G$ is isomorphic to a normal profinite HNN-extension, namely, $G\cong \hnn(G_1, H,   \tau_{g_1}f)$ (see Lemma \ref{conjugados da normal}). 

As discussed in the abstract case (see Section \ref{cota bonita conjugados}), we have an action on the right of $\overline{\aut}_{G_1}(H) \times C_2$ on the space of orbits $\overline{N}_{G_1}(H) \backslash \aut(H)$, given by  \begin{equation} \label{acao bunitaap} \left( \overline{N}_{G_1}(H) \cdot f \right) \cdot (\psi, \iota) = \overline{N}_{G_1}(H) \cdot  (f^{-1})^{\psi}=\overline{N}_{G_1}(H) \cdot \left(f^{\psi}\right)^{-1},\end{equation}
where $( \psi, \iota) \in  \overline{\aut}_{G_1}(H) \times C_2$ and $f \in \aut(H).$ We will denote the space of orbits of this action by $\overline{N}_{G_1}(H) \backslash \aut(H)/ ( \overline{\aut}_{G_1}(H) \times C_2).$ 


%Let $H$ be a finite subgroup of $G_1$. Then $\overline{N}_{G}(H)$ acts on $Aut(H)$ by composition on the left, furthermore $\overline{N}_{G_1}(H)$ also acts on $Aut(H)$  by restricting the action of $\overline{N}_{G}(H)$ to
%$\overline{N}_{G_1}(H)$.  We also can define an action of $(\overline{Aut}_{G_1}(H) \rtimes C_2)$ on $Aut(H)$ on the right as follows:
%$$f \cdot (\psi, c)= \psi f^{-1}\psi^{-1}=(f^{-1})^{\psi^{-1}}, \mbox{ with } (\psi, c) \in  \overline{Aut}_{G_1}(H) \rtimes C_2, f \in Aut(H),$$ where $c$ is the generator of $C_2$.



Note that $\inn(H)\leq \overline{N}_{G}(H)$ so we can pass to $\out(H)$ using tilde instead of overline. Then we can deduce the bound  for the number  of isomorphism classes of the profinite HNN-extensions $\hnn(G_1, H, K, f, t)$, where $H$ and $K$ are fixed conjugate finite subgroups. 
%with fixed conjugate $H$ and $K$. 

\begin{thm}\label{profinite isomorphisms} Let $H$ and $K$ be fixed finite conjugate subgroups of a profinite $\OE$-group $G_1$. Then the number $\hat n_H$ of the isomorphism classes of profinite HNN-extensions $G=\hnn(G_1, H, K, f, t)$ is less or equal than $$\left| \widetilde N_{G_1}(H) \backslash \out(H)/ \left( \widetilde{\aut}_{G_1}(H) \times C_2\right) \right|.$$
\end{thm}
\begin{proof}
Follows from Proposition \ref{nonnormalhnnprof}, Theorem \ref{profinite isomorphism number} and the fact that $N_{G_1}(H)\\ \subset N^{+} \subset N_{G}(H)$ (see (\ref{NG1H em N+})).  
\end{proof}


\begin{exa} Let $G= G_1\amalg_{H, \mu} (H\rtimes_{\alpha} \widehat\Z)$ be a  free profinite product with finite amalgamation, where $G_1$ is a profinite $\OE$-group and $\mu$ is the natural monomorphism from $H$ to $H\rtimes_{\alpha} \widehat\Z= H\rtimes_{\alpha} \overline{\left\langle t \right\rangle}.$ This case was missing in the study of free profinite products with finite amalgamation in \cite{BPZ2} .Note that $G\cong \hnn(G_1, H, \gamma, t),$ where $\gamma = \alpha(t)\circ \mu \in \aut(H).$ Then the number of isomorphism classes of such profinite free products with finite amalgamation $G_1\amalg_{H, \mu} (H\rtimes_{\alpha} \widehat\Z)$  with $H$ fixed is less or equal than  
$$n_{_{H}}:=| \widetilde{N}_{G_1}(H) \backslash \out(H)/ ( \widetilde{\aut}_{G_1}(H) \times C_2) |.$$
\end{exa}


%Let $\hat{S}$ be a set of finite subgroups of the profinite group $G_1$ isomorphic to a subgroup $H$ of $G_1$. The $Aut(G_1)$ acts on $\hat{S}$ in a natural way. Consider the set of orbits of this action given by $Aut(G_1)\backslash \hat{S}$. 

%If $L \in \hat{S}$, then by Corollary \ref{profinite isomorphisms} the number of isomorphism classes of  normal profinite HNN-extensions  $G=HNN(G_1, L, f, t)$  satisfies   $$\hat n_{L}\leq \left| \widetilde{N}_{G_1}(L) \backslash Out(L)/ ( \widetilde{Aut}_{G_1}(L) \times C_2)\right|.$$  

%Then we deduce the following from the Corollary \ref{profinite isomorphisms}.
%\textcolor{red}{Esse teorema abaixo ao meu ver não faz mais sentido. Ele é igual ao teorema \ref{profinite problem} com $H=K$. Eu sugiro tirar.}
%\begin{thm}
%Let $H$ and $K$ be conjugate finite subgroups in the profinite $\OE$-group $G_1$. Then the number of isomorphism classes of profinite HNN-extensions $\hnn(G_1, H, K, f, t)$ is less or equal than  $$\sum_{H} \hat n_{H},$$
%where $H$ ranges over representatives of the orbits  $Aut(G_1)\backslash \hat{S}$.
%\end{thm}



%\textcolor{red}{By Proposition \ref{refhnn},  the normal profinite HNN-extensions $HNN(G_1,H,f,t)$ and $HNN(G_1,H,f',t')$ are isomorphic whenever $\langle  f\rangle=\langle f'\rangle \leq Aut(H)$. Let $c_H$ be a number of cyclic subgroups in $Out(H)$. Then we have the following estimate.}

%Note that $f\in \overline N_G(H)$ and so the profinite normal HNN-extensions $HNN(G,H,f,t)$ and $HNN(G,H,f',t')$ are isomorphic whenever $\langle  f\rangle=\langle f'\rangle$. Let $c_H$ be a number of cyclic subgroups in $Out(H)$. Then we have the following estimate.
%\textcolor{red}{
%\begin{cor} Let $H$ and $K$ be   conjugate finite subgroups in the profinite OE-group $G_1$. The number of isomorphism classes of profinite HNN-extensions $\hnn(G_1, H, K, f, t)$  ($H$ not fixed) is less or  equal to $$c_H \cdot |Aut(G_1)\backslash \hat{S}|.$$
%\end{cor}
%}

\begin{rem} \label{finite-by-cyclic profinite} 
  Let $G=\hnn(G_1, H, K, f, t)$ %, $B=\hnn(B_1, Y, W, g, z)$ 
be a profinite HNN-extension with finite associated subgroups such that $G_1=M_1\rtimes \widehat\Z$, where $M_1$ is a finite subgroup, distinct from $H$ and $K.$ Then $G_1$ is not a profinite $\OE$-group. However, the results obtained in this section for this $G_1$  remain valid.


%Similarly, let $G=\hnn(G_1, H, K, f, t)$, $B=\hnn(B_1, Y, W, g, z)$ be  profinite HNN-extensions with finite associated subgroups, whit $G_1=M_1\rtimes \widehat\Z$, where $M_1$ finite such that $M_1\neq H$. Then $G_1$ is not a profinite $\OE$-group. \textcolor{red}{However, the results obtained in this section for when $G_1$ is a profinite $\OE$-group remain valid.}
%, but  Propositions \ref{2a} also holds in this case.

\end{rem} 

%\begin{rem} \label{finite-by-cyclic ii} 
%The results of this and previous sections are also valid for HNN-extensions which are of the form $G=\hnn(G_1, H, K, f, t)$ where $G_1$ is as in Remark \ref{finite-by-cyclic}.

%The results from the previous Sections are also valid for HNN-extensions $G$ are of the form $G=\hnn(G_1, H, K, f, t)$ where $G_1$ is as in the Remark \ref{finite-by-cyclic}.
%\end{rem}

\subsection{Profinite vs Abstract}

In this subsection we consider an abstract HNN-extension $\hnn(G_1,H,K,f,t)$ of an $\Oe$-group and its profinite completion $\hnn(\widehat G_1,H,K,f,t)$. 



\begin{rem} \label{rem dos normalizadores e densidade}
%Suppose $[G_1:N_{G_1}(H)]=[\widehat G_1:N_{\widehat G_1}(H)]<\infty$. Then we have $\overline{N}_{\widehat G_1}(H)=\overline{N}_{ G_1}(H)$ because $N_{G_1}(H)$ is dense in $N_{\widehat G_1}(H).$ 
Since $\inn(H)\leq \overline{N}_{G_1}(H)$ ( respectively $\inn(H) \leq \overline{\aut}_{G_1}(H)$), the statement that $\widetilde{N}_{\widehat G_1}(H)=\widetilde{N}_{ G_1}(H) $ ( respectively $\widetilde{\aut}_{G_1}(H) = \widetilde{\aut}_{\widehat G_1}(H)$) immediately implies $\overline{N}_{\widehat G_1}(H)=\overline{N}_{ G_1}(H)$ ( respectively $\overline{\aut}_{G_1}(H) = \overline{\aut}_{\widehat G_1}(H)$).
\end{rem}

\begin{rem} \label{tilde implica overline} Recall that
if the associated subgroups $H$ and $K$ are conjugate in  $G_1$, say $H^{g_1}=K\,(g_1 \in G_1),$ then $\hnn(G_1,H,K,f,t)$ is isomorphic to a normal  HNN-extension $G=\hnn(G_1, H, f_1) \, \mbox{ with } f_1 \in \aut(H)$  (see Lemma \ref{conjugados da normal}), i.e.,  one can assume $G= \hnn(G_1, H, f),$ where $f \in \aut(H)$ is normal.  If $\overline{N}_{\widehat G_1}(H)=\overline{N}_{G_1}(H)$ (this is the case if for example $N_{G_1}(H)$ is dense in $N_{\widehat G_1}(H)$), then by Proposition \ref{cp1} 

$$\overline{N}_{\widehat G}(H) = \langle \overline{N}_{\widehat G_1}(H),  f\rangle= \langle \overline{N}_{G_1}(H), f\rangle=\overline{N}_{G}(H),$$ $$\widetilde{N}_{\widehat G}(H) = \langle \widetilde{N}_{\widehat G_1}(H), \tilde f\rangle= \langle \widetilde{N}_{G_1}(H), \tilde f\rangle=\widetilde{N}_{G}(H),$$     
where $\tilde f$ is the image of $f$ in $\out(H)$.
\end{rem}

The situations described in the remark above occur in particular when $G_1$ is finite. 


Thus having abstract and profinite HNN-extensions $\hnn(G_1,H,K,f,t)$  and  $\hnn(\widehat G_1,H,K,f,t)$ of an $\Oe$-group $G_1$ we have two groups $\overline {\Gamma}_{HK}$ and $\overline{\Gamma}_{\widehat{HK}}$ acting on $\Iso(H,K)$ (the first in fact is a subgroup of the second) and the cardinality of genus is the difference between the the orbits of these action: more precisely, it is the number of $\overline {\Gamma}_{HK}$-orbits lying in $\overline{\Gamma}_{\widehat{HK}}$-orbit of $f$.  Of course the $\overline {\Gamma}_{HK}$-orbits will coincide with $\overline{\Gamma}_{\widehat{HK}}$-orbits if the groups $\overline{\Gamma}_{HK}$ and $\overline{\Gamma}_{\widehat{HK}}$ are equal.   


Recall that a pair $(H,K)$ of finite subgroups of $G_1$ (resp. $\widehat G_1$) is a swapping pair if there exists an  automorphism $\psi\in \aut(G_1)$ (resp. $\in \aut(\widehat G_1)$) that swaps the conjugacy classes of $H$ and $K$ in $G_1$ (resp. in $\widehat G_1)$. Of course, to have  $\overline{\Gamma}_{HK}=\overline{\Gamma}_{\widehat{HK}}$ it is necessary for $(H,K)$ to be either a swapping pair or not a swapping pair in both $G_1$ and $\widehat G_1$.  Having this, it suffices to have the density of $\Gamma_{HK}$ in $\Gamma_{\widehat{HK}}$, since $\Iso(H,K)$ is finite. However, this would imply the density of $\aut_{G_1}(H)$  in $\aut_{\widehat G_1}(H)$ and for infinite $G_1$ it happens not in too many cases despite the fact that it is quite difficult to check. The next lemma gives more transparent sufficient conditions that can be quite often checked when $\out(H)$ is small; it  will be used in  Subsection \ref{section 5 1}. 

%\begin{lem} \label{equal orbits}
%With the assumptions in this subsection, suppose  %$Aut_{G_1}(H)$ is dense in $Aut_{\widehat G_1}(H)$ and  
%     $(H,K)$ is either a swapping pair in  $G_1$ or not a swapping pair in $\widehat G_1$. 
 %Then $\overline{\Gamma}_{HK}=\overline{\Gamma}_{\widehat{HK}}$.   
%\end{lem}

\begin{lem} \label{orbits non conjugated}
With the assumptions in this subsection, suppose the following conditions are satisfied 

\begin{enumerate}
    \item [i)]  $\widetilde{\aut}_{\widehat{G}_1}(H)=\widetilde{\aut}_{G_1}(H)$ and $[G_1:N_{G_1}(K)]=[\widehat G_1:N_{\widehat G_1}(K)]<\infty$; 
    %and 
    %$N_{G_1}(K)C_{\widehat G_1}(K)=N_{\widehat G_1}(K)$;
    
    \item [ii)] $(H,K)$ is either a swapping pair in  $G_1$ or not a swapping pair in $\widehat G_1$. 
\end{enumerate}
Then $\overline{\Gamma}_{HK}=\overline{\Gamma}_{\widehat{HK}}$. 
\end{lem}
\begin{proof}
First assume that there exists $\psi \in \aut(G_1)$ such that $\psi(K)=H^{g_1}$ and $\psi(H)=K$, for some $g_1\in G_1$ and observe that $\psi$ extends to an isomorphism $\hat\psi\in \aut(\widehat G_1)$.  Let $f \in \Iso(H, K)$. It follows from the definitions of the respective $\iota$ in (\ref{def iota abst}) and (\ref{def iota prof})  that 
$$\overline{\Gamma}_{\widehat{HK}}\cdot f\ni \iota\cdot f =\tau_{\hat\psi^{-1}(g_1^{-1})} \left(f^{\hat\psi}\right)^{-1}=\tau_{\psi^{-1}(g_1^{-1})} \left(f^{\psi}\right)^{-1}\in \overline{\Gamma}_{HK}\cdot f$$

Take $(\hat{g}, \hat{\alpha}) \in \Gamma_{\widehat{HK}}$ such that $\tilde{w}$ is its image in $\widetilde{\Gamma}_{\widehat{HK}}.$ 
According to item i) %$\tau_{\hat{g}} \in \inn(\widehat G_1)$ and $\hat \alpha \in \aut_{\widehat G_1}(H)$ 
there exists $\tau_g \in \inn(G_1)$ and $\alpha \in \aut(G_1)$ such that $\tau_{\hat{g}}(k)=\tau_{g}(k), \, \forall k \in K$ and $\hat \alpha|_{_{H}} = \alpha|_{_{H}}$. Let $w$ be the image of $(g, \alpha)$ in $\widetilde{\Gamma}_{HK}.$ In this way, we have  $$\overline{\Gamma}_{\widehat{HK}}\cdot f\ni \tilde{w} \cdot f = \tau_{\hat g} f^{\hat{\alpha}^{-1}} = \tau_{ g} f^{\alpha^{-1}} =  w \cdot f\in \overline{\Gamma}_{HK}\cdot f,\, \forall f \in \Iso(H, K).$$

If the second possibility of item ii) occurs, the proof is similar but simpler since we do not have the action of $\iota$ since $\overline{\Gamma}_{HK}=\widetilde{\Gamma}_{HK}$ and $\overline{\Gamma}_{\widehat{HK}}=\widetilde{\Gamma}_{\widehat{HK}}.$
    
\end{proof}

\begin{rem} \label{altenative conditions}
Note that the condition $[G_1:N_{G_1}(K)]=[\widehat G_1:N_{\widehat G_1}(K)]<\infty$ in item i) of the Lemma \ref{orbits non conjugated} can be replaced by the condition $$N_{G_1}(K)C_{\widehat G_1}(K)=N_{\widehat G_1}(K)$$ together with the hypothesis that $K$ is a conjugacy distinguished subgroup in $G_1$ (that is, if $K$ and $L \leq G_1$ are conjugate in $\widehat G_1$, then they are conjugate in $G_1$), since in both cases it follows that given $\tau_{\widehat g} \in \inn(\widehat G_1)$ there exists $\tau_{g} \in \inn(G_1)$ such that $\tau_{\hat{g}}(k)=\tau_{g}(k), \, \forall k \in K.$
\end{rem}


\section{Genus of HNN-extensions}\label{secao 5}

Let $H$ and $K$ be finite subgroups of an $\Oe$-group $G_1$ and $f\in \Iso(H,K)$. 
Note here that the abstract HNN-extension $\hnn(G_1, H, K, f,t)$  is residually finite and its profinite completion is the proper profinite HNN-extension $\hnn(\widehat{G}_1, H, K, f,t)$ (see \cite{RZ}, Proposition 9.4.3)   and so we can apply the results of the previous sections.

\subsection{Fixed associated subgroups} \label{section 5 1}

Let  $$\mathcal{F}=\{ \hnn(G_1, H, K, f, t) | f \in \Iso(H,K) \mbox{ and } H, K \mbox{ fixed  }\}$$ be the class of all  
abstract HNN-extensions $\hnn(G_1, H,K, f, t)$, where $H$ and $K$ are fixed finite subgroups of $G_1$ and $f \in \Iso(H,K)$. 
Recall that
\begin{equation*}% \label{N mais com completamento}
N^+=\left\{ gt^{- 1} \in N_{\widehat G}(H) \hspace{0.1cm} \left| \hspace{0.1cm}  \overline{\left\langle g, \widehat G_1 \right\rangle}=\widehat G \right. \right\},   
\end{equation*} and  \begin{equation} \label{cadeia de N mais etc} N_{ G_1}(H) \subset N_{\widehat G_1}(H) \subset N^{+} \subset N_{\widehat G}(H). \end{equation} 
%Denote by  $\overline{N}^{+}$ and $\widetilde{N}^{+}$ the natural images of $N^{+}$ in $\aut(H)$ and $\out(H)$, respectively.  

%$N^{+}=\left\{g \cdot t^{- 1} \in N_{\widehat G}(H) \hspace{0.1cm} \left| \right. \hspace{0.1cm}  \overline{\left\langle g, \widehat G_1 \right\rangle}=\widehat G\right\},$



Given $G\in \mathcal{F}$, our goal now is to find the cardinality of the set $g(G, \mathcal{F})$.   Define \begin{equation} \label{Omega como uniao} \Omega:=\bigcup_{n\in N^{+}}\overline{\Gamma}_{\widehat{HK}} \cdot (f\tau_{n^{-1}}) \subset \Iso(H, K)\end{equation} as the union of orbits of  $f\tau_{n^{-1}}\,(n \in N^{+}) \in \Iso(H, K)$ with respect to action given in Definition \ref{novaacaopp}. Consider $B=\hnn(G_1,H,K,f_1,t_1)\in \mathcal{F}$. By Proposition \ref{nonnormalhnnprof} $\widehat G \cong \widehat B$ if and only if $f_1$ belongs to $\overline{\Gamma}_{\widehat{HK}}$-orbit of $f\tau_{n^{-1}}$ for some $n\in N^{+}$ if and only if $f_1 \in \Omega.$ 

Thus, we have the following 

\begin{thm}\label{fixed H} Let $H$ and $K$ be finite subgroups of an $\Oe$-group $G_1.$ Then the cardinality of  genus $g(G, \mathcal{F})$ of HNN-extension $G=\hnn(G_1,H,K,f,t)$ is 
 $$|g(G, \mathcal{F})| = \left|\overline{\Gamma}_{HK}  \backslash \Omega\right|\,\,(\mbox{ see formula  } (\ref{Omega como uniao})).$$ 
\end{thm}

%\begin{cor}\label{covered out(H)} If    $\widetilde{N}_{G_1}(H)=\out(H)$, then $|g(G, \mathcal{F})| =1.$ \end{cor}



Next  we find the cardinality of the genus of $G$ when the subgroups  $\overline{\Gamma}_{\widehat{HK}}$ and $\overline{\Gamma}_{HK}$ of $\Iso(H,K)$ coincide (see Lemma \ref{orbits non conjugated}).   Note that this is always the case if $G_1$ is finite ($G_1 = \widehat{G}_1$) (see Lemma \ref{orbits non conjugated}). %(see Lemmas \ref{orbits non conjugated} and \ref{orbits conjugated}), so we have the following




Let $\kappa: \Iso(H,K)\longrightarrow  \overline\Gamma_{HK}\backslash \Iso(H,K)$ be the natural map and \begin{equation} \label{O conjunto S} R:=\{f\tau_{n^{-1}}\,|\, n \in N^{+}\}.
\end{equation} 

The following theorem gives the precise value for the genus of HNN-extensions; in particular it strengthens  Theorem 5.1 in \cite{BZ}.

\begin{thm} \label{nova conta non conjugadosss} Let $H$ and $K$ be finite subgroups in the $\Oe$-group $G_1$ and $G=\hnn(G_1,H,K,f,t)$ be an HNN-extension. If  $\overline\Gamma_{HK}=\overline\Gamma_{\widehat{HK}}$, then $|g(G, \mathcal{F})| = |\kappa(R)|.$ This is the case if \begin{enumerate}
    \item [i)]  $\widetilde{\aut}_{\widehat{G}_1}(H)=\widetilde{\aut}_{G_1}(H),$  $[G_1:N_{G_1}(K)]=[\widehat G_1:N_{\widehat G_1}(K)]<\infty$ and; 
    %and 
    %$N_{G_1}(K)C_{\widehat G_1}(K)=N_{\widehat G_1}(K)$;
    
    \item [ii)] $(H,K)$ is either a swapping pair in  $G_1$ or not a swapping pair in $\widehat G_1$. 
\end{enumerate} 
\end{thm}

\begin{proof} Follows directly from Lemma \ref{orbits non conjugated}. \end{proof}



Since $\widetilde{N}^{+} \subset \widetilde{N}_{\widehat G}(H),$ %=\widetilde{N}_{G_1}(H)=\widetilde{N}_{\widetilde G_1}(H),$ 
the corollary below follows directly from the previous result together with Theorem \ref{classnonnormalabs} and the definition of $R$.  
\begin{cor} \label{normal iguais genero 1 non conjugated} Suppose, in addition to the
hypotheses of Theorem \ref{nova conta non conjugadosss} that $\widetilde{N}_{\widehat G}(H)=\widetilde{N}_{G_1}(H),$
then $|g(G, \mathcal{F})|=1$. In particular, this holds if $\widetilde{N}_{G_1}(H)=\out(H)$.
\end{cor}

\begin{proof} It follows from Lemma \ref{orbits non conjugated} that the actions of $\overline{\Gamma}_{HK}$ and  $\overline{\Gamma}_{\widehat{HK}}$ on $\Iso(H, K)$ correspond to the same set of orbits. The result then follows from Theorem \ref{fixed H}.
\end{proof}

 
%where \begin{equation} \label{N mais com completamento} N^+=\{ gt^{- 1} \in N_{\widehat G}(H) \hspace{0.1cm} | \hspace{0.1cm}  \overline{\left\langle g, \widehat G_1 \right\rangle}=\widehat G\}\,\, (\mbox{compare with formula } (\ref{N mais})).
%\end{equation}  

\subsubsection{Conjugate subgroups $H$ and $K$}

If the associated subgroups $H$ and $K$ are conjugate in the base group $G_1$,  one can assume w.l.o.g. that the HNN-extension $G=\hnn(G_1, H, K, f,t)$ is normal (see Lemma \ref{conjugados da normal}). Therefore,  the set $R$ from the preceeding subsection can be seen as a subset of $\aut(H).$  Furthermore,  
if $gt^{-1} \in N^{+}$ then $g \in N_{\widehat G}(H)$. %(see formula (\ref{N mais com completamento})). 
Note that the image of $gt^{-1}$ in $\aut(H)$ is given by $\tau_{_{(gt^{-1})^{-1}}}=f^{-1} \tau_{g^{-1 }}.$ Therefore the set $R$ (see formula (\ref{O conjunto S})) can be rewritten as \begin{equation} \label{conjunto S com BZ} R=\{f\tau_{n^{-1}}\,\mid\, n \in N^{+}\}= \{ \tau_{g}^{-1} \mid g \in N_{\widehat G}(H) \hspace{0.1cm} \mbox{and} \hspace{0.1cm}   \overline{ \langle g, \widehat G_1 \rangle}  =\widehat G \}.\end{equation} 

%\begin{equation} \label{conjunto S com BZ} R=\{f\tau_{n^{-1}}\,|\, n \in N^{+}\}= \left\{ \tau_{g^{-1}} \in \overline{N}_{\widehat G}(H) \hspace{0.1cm} \left|  \hspace{0.1cm}   \overline{ \left\langle g^{\pm 1}, N_{\widehat G_1}(H) \right \rangle} \right. =N_{\widehat G}(H) \right\}.\end{equation} 


%On the other hand, if we assume that $\widetilde{\aut}_{\widehat{G}_1}(H)=\widetilde{\aut}_{G_1}(H),$ then by Remarks \ref{rem dos normalizadores e densidade} and \ref{tilde implica overline} we have $\overline{N}_{\widehat G_1}(H) =\overline{N}_{G_1}(H)$ and $\overline{N}_{\widehat G}(H) =\overline{N}_{G}(H),$ and therefore we obtain that $S$ can again be rewritten in the following way \begin{equation} \label{conjunto S com BZZ} S=\left\{ \tau_{g^{-1}} \in \overline{N}_{G}(H) \hspace{0.1cm} \left| \hspace{0.1cm} \overline{ \left\langle g^{\pm 1}, N_{\widehat G_1}(H) \right \rangle} \right. =N_{\widehat G}(H) \right\}.\end{equation} %With this additional hypothesis, we note that the set given in (\ref{conjunto S com BZZ}) coincides with the set $S$ defined in Theorem 5.1 (1) of \cite{BZ}. 


%Now, (see Corollary \ref{corhnn}) $\kappa$ induces \begin{equation} \label{kappa function} \tilde\kappa: \out(H)\longrightarrow \widetilde{N}_{G_1}(H) \backslash \out(H)/  ( \widetilde{\aut}_{G_1}(H) \times C_2).\end{equation}

Suppose $\mu\in \aut(H)$ and define $\tilde{\mu}$ to be the image of $\mu$ in $\out(H).$ Define the following subset of $\out(H)$ \begin{equation}\label{cS} \cR= \tilde{f}\widetilde{N}^+ := 
\{\tilde f \tilde \tau_{n}^{-1}\,|\, n \in N^{+}\},\end{equation} where $\cR$ is the natural image of the subset $R$ (see (\ref{O conjunto S})) in $\out(H).$

It follows that $\cR=\tilde{f}\widetilde{N}^+$ is a subset of $\tilde f \widetilde{N}_{\widehat G}(H)$. As discussed in Subsubsection \ref{quotasprofconjugates}, $\widetilde{\aut}_{\widehat G_1}(H)$ acts on $\out(H)$ on the right by conjugation.  % and $C_2$ acts by inversion. 
The group $\widetilde{N}_{\widehat G_1}(H)$ acts on $\out(H)$ on the left by multiplication. It follows that $\widetilde{N}_{\widehat G_1}(H) \cdot \tilde{f} \widetilde{N}^+ \cdot \widetilde{\aut}_{\widehat G_1}(H)$ is invariant with respect to the action of $\widetilde{N}_{G_1}(H)$ on the left and with respect to the action of $\widetilde{\aut}_{G_1}(H)$ on the right. In the case of conjugated associated subgroups  in the base group, we have the following formula for the genus.

\begin{thm} \label{1.8 da introducao}  Let $H$ and $K$ be conjugate finite subgroups   in an $\Oe$-group $G_1$ and $G=\hnn(G_1,H,K,f,t)$ be an HNN-extension.  Then $$|g(G, \mathcal{F})| = | \widetilde{N}_{G_1}(H) \backslash \widetilde{N}_{\widehat G_1}(H) \cdot \cR  \cdot \widetilde{\aut}_{\widehat G_1}(H)/( \widetilde{\aut}_{G_1}(H) \times C_2)|,$$ where $C_2$ acts on $\out(H)$ by inversion. 
\end{thm}
%\tilde{f}\left( \widetilde{N}^+\right)^{-1}

\begin{proof} It can be assumed that $G$ is normal and $f \in \aut(H)$ (see Lemma \ref{conjugados da normal}). Let $B=\hnn(G_1,H,f_1,t_1)$ with $f_1 \in \aut(H).$ By Proposition \ref{nonnormalhnnprof}, $\widehat B \cong \widehat G$ if and only if $f_1  \in \overline{N}_{\widehat G_1}(H) \cdot R \cdot \overline{\aut}_{\widehat G_1}(H)$ if and only if  $\widetilde{f_1} \in \widetilde{N}_{\widehat G_1}(H) \cdot \cR \cdot \widetilde{\aut}_{\widehat G_1}(H)$ (since $\inn(H)\leq \overline{N}_{G_1}(H) \cap \overline{\aut}_{G_1}(H)$) and the result follows from Corollary \ref{corhnn}. 
\end{proof}


Now, let 
\begin{equation} \label{kappa function} \tilde\kappa: \out(H)\longrightarrow \widetilde{N}_{G_1}(H) \backslash \out(H)/  ( \widetilde{\aut}_{G_1}(H) \times C_2)\end{equation}
be the natural map (see Corollary \ref{corhnn}).
Then we have the following 

%Recall that $\overline{N}_{\widehat G_1}(H)=\overline{N}_{ G_1}(H)$ for example when  $[G_1:N_{G_1}(H)]=[\widehat G_1:N_{\widehat G_1}(H)]<\infty$ (see Remark \ref{rem dos normalizadores e densidade}) or when $G_1$ is finite. We have the following result.  

\begin{cor} \label{nova conta conjugadoss} Let $H$ and $K$ be conjugate finite subgroups of an $\Oe$-group $G_1$ and $G=\hnn(G_1,H,K,f,t)$ be an HNN-extension. Suppose that $\widetilde{\aut}_{\widehat G_1}(H)\\=\widetilde{\aut}_{G_1}(H)$ and $\widetilde{N}_{\widehat G_1}(H)=\widetilde{N}_{ G_1}(H).$  %(for example if $G_1$ is finite or when \textcolor{blue}{ $[G_1:N_{G_1}(H)]=[\widehat G_1:N_{\widehat G_1}(H)]<\infty$}). 
Then $$|g(G, \mathcal{F})| = |\tilde \kappa(\cR)|=| \widetilde{N}_{G_1}(H) \backslash \cR/  ( \widetilde{\aut}_{G_1}(H) \times C_2)|.$$ 
\end{cor}

\begin{proof} The result follows from Theorem \ref{1.8 da introducao}. 
\end{proof}


\begin{cor} \label{cotas n} Let $H$ be a finite subgroup of an $\Oe$-group $G_1$ and $G=\hnn(G_1,H,f,t)$  be a normal HNN-extension.  Suppose $\widetilde{N}_{\widehat G_1}(H)=\{1\}$ (for example if $H$ is malnormal in $\widehat G_1$) and $\widetilde{\aut}_{ G_1}(H)=\widetilde{\aut}_{\widehat{G}_1}(H)$ (this is the case if $|G_1|<\infty$). Then 

$$|g(G, \mathcal{F})|\leq \left\{\begin{array}{cc}
  1,&\mbox{ if } |\tilde f| \leq 2,\\
\displaystyle\frac{\phi(|\tilde f|)}{2} , &\mbox{ if } |\tilde f| \geq 3,
\end{array}\right.$$
where $\phi$ is Euler's function.  
\end{cor}

\begin{proof}
By Proposition \ref{cp1} $\widetilde{N}_{\widehat{G}}(H)=\left \langle \widetilde{N}_{\widehat{G }_1}(H), \tilde{f} \right\rangle=\left \langle \tilde{f} \right\rangle.$ Note that  $\mathcal{R}=\left\{ \tilde{\tau}_{g^{-1}}\,|\, \tilde{\tau}_{g^{-1}} \mbox{ is a generator of } \left \langle \tilde{f}\right \rangle \right\}.$ Therefore, the bound can be obtained by considering only the action of $C_2$ on $\mathcal{R}$ as in Corollary \ref{nova conta conjugadoss}. 
 %$\mathcal{R}= \tilde f \langle \tilde{f}^*\rangle,$ where $*$ denotes the set of generators.}
\end{proof} 

\begin{rem}\label{H subgroup cyclic}
In the preceding result, if $\aut(H)$ is abelian (for example, when $H$ is cyclic) then $|g(G, \F)|=\frac{\phi(|\tilde f|)}{2}$ if $|\tilde f|\geq 3.$      
\end{rem} 

%%%%%%%%%%
Denote by $tor(G)$ the subset of elements of finite order in a given group $G.$ 
%When $G$ is finitely generated nilpotent, $tor(G)$ is a finite subgroup of $G$.

\begin{exa} Suppose that $G_1$ is a finitely generated abelian group which is not virtually infinite cyclic and let $H$ be one of its finite subgroups. By Proposition \ref{virtually cyclic} $G_1$ is an $\Oe$-group. Let $G=\hnn(G_1,H,f,t)$ be a normal HNN-extension. Since $G_1$ is finitely generated abelian $\widetilde{N}_{G_1}(H)= \widetilde{N}_{\widehat{G }_1}(H)=\{1\}$ and furthermore $\widetilde{\aut}_{ G_1}(H)=\out(H)=\widetilde{\aut}_{\widehat{G}_1}(H).$ It follows from Corollary \ref{cotas n} that the genus satisfies %$$|g(G, \mathcal{F})| \leq \dfrac{\phi(|\tilde f|)+1}{2}.$$ 
$$|g(G, \mathcal{F})|\leq \left\{\begin{array}{cc}
  1,&\mbox{ if } |\tilde f| \leq 2,\\
\displaystyle\frac{\phi(|\tilde f|)}{2}, &\mbox{ if } |\tilde f| \geq 3.
\end{array}\right.$$ 

If $G_1$ is virtually infinite cyclic and $H \neq tor(G_1)$ the same result is valid by Remarks \ref{finite-by-cyclic} and \ref{finite-by-cyclic profinite}.  
For $H=tor(G_1)$ we have that the HNN-extension is of the form $G=H \rtimes F,$ where $F$ is a free group and the genus in this case was calculated in Theorem 3.26 in \cite{GZ}. 
\end{exa}

%\begin{exa} Suppose that $G_1$ is a finitely generated abelian group which is not virtually infinite cyclic and let $H$ be one of its finite subgroups. By Proposition \ref{virtually cyclic} $G_1$ is an $\Oe$-group. Let $G=\hnn(G_1,H,f,t)$ be a normal HNN-extension. $\widetilde{N}_{\widehat{G}}(H)=\left \langle \widetilde{N}_{\widehat{G }_1}(H), \tilde{f} \right\rangle=\left \langle \tilde{f} \right\rangle.$ It follows from Corollary \ref{nova conta conjugadoss} that  
 %$$|g(G, \mathcal{F})|= |\mathcal{R}/( \widetilde{\aut}_{G_1}(H) \times C_2)|\leq |\widetilde{N}_{G}(H)/( \widetilde{\aut}_{G_1}(H) \times C_2)|.$$ In particular, if $H$ is cyclic then $\widetilde{N}_{G}(H)=\left\langle f \right\rangle$ and therefore $$|g(G, \mathcal{F})|=|\mathcal{R}/C_2|\leq \left\{\begin{array}{cc}
 %\displaystyle\frac{n}{2}+1,&\mbox{ if } n=|f| \mbox{ is even, }\\
%\displaystyle\frac{n-1}{2}+1 , &\mbox{ if } n=|f| \mbox{ is odd.}
%\end{array}\right.$$ If $G_1$ is virtually infinite cyclic and $H \neq tor(G_1)$ the same result is valid for Remarks \ref{finite-by-cyclic} and \ref{finite-by-cyclic ii}. For $H=tor(G_1)$ we have that the HNN-extension is of the form $G=H \rtimes F,$ where $F$ is a free group and the genus in this case was calculated in Theorem 3.26 in \cite{GZ}. 
%\end{exa}


%%%%%%%%%%%%%%%%%%%%%%%%%%%%%%%%%%%%%%%%%%%
%%%%%  3-COROLARIO REPETITIVO FOI RETIRADO 
%%%%%%%%%%%%%%%%%

The second part of the result below, which deals with genus equal to 1, is similar to what was done in Theorem 3.1 in \cite{BZ} when $G$ is restricted to the class $\mathcal{F}$ and the finite associated subgroups are conjugate in $G_1,$ that is, when the HNN-extension can be considered normal.


\begin{cor} \label{autG classe} 
Let $H$ and $K$ be conjugate finite subgroups of an $\Oe$-group $G_1$  and $G=\hnn(G_1, H, K, f, t)$ be an HNN-extension with $f \in \Iso(H, K).$ Suppose $\widetilde{\aut}_{\widehat G_1}(H)=\widetilde{\aut}_{G_1}(H)$ and $\widetilde{N}_{\widehat G_1}(H)=\widetilde{N}_{ G_1}(H)$ (for example $|G_1|<\infty$).  Then $$|g(G, \mathcal{F})| \leq | \tilde\kappa(\widetilde{N}_{G}(H))|.$$ In particular, if $\widetilde{N}_{\widehat G}(H)=\widetilde{N}_{G_1}(H)$ (this is the case if $\tilde f\in \widetilde{N}_{ G_1}(H),$ where $\tilde f$ is the image of $f$ in $\out(H)$) or $|\out(H)|\leq 2,$ then $|g(G, \mathcal{F})|=1$.
\end{cor}
\begin{proof} Note that $\widetilde N_{\widehat G}(H)=\langle \widetilde N_{\widehat G_1}(H), \tilde f\rangle=\widetilde N_G(H)$ in this case. So  the first part follows directly from Corollary \ref{nova conta conjugadoss} and the fact that $N^{+} \subset N_{\widehat G}(H),$ because the genus is less than or equal to $$| \tilde\kappa(\tilde f \cdot \widetilde{N}_{G}(H))|=| \tilde\kappa( \widetilde{N}_{G}(H))|.$$ The second statement follows directly from the first. 
\end{proof}

\begin{rem} \label{case polycyclic}
Let $G_1$ be a polycyclic-by-finite group and $H$ a finite subgroup of $G_1$ such that $\widetilde{\aut}_{\widehat G_1}(H)= \widetilde{\aut}_{G_1}(H)$. Consider  a normal HNN-extension $G=\hnn(G_1, H, f, t)$. Then the cardinality of the genus of $G$ in $\mathcal{F}$  satisfies $$|g(G, \F)| \leq  | \tilde\kappa( \widetilde{N}_{G}(H))|,$$ by Proposition \ref{cp1},  Corollaries  \ref{8} and \ref{autG classe} since $\widetilde N_{\widehat G}(H)=\widetilde N_{G}(H)$. In particular, the result holds if $G_1$ is finite. 
\end{rem}

The  result, which provides more effective conditions for calculating the genus of normal HNN-extensions, will be given in the following corollary. Note that it is a generalization of Corollary 5.5 in \cite{BZ}.  

%and since $\widetilde{\aut}_{\widehat{G}_1}(H)=\widetilde{\aut}_{G_1}(H)$ one deduces from Remark \ref{rem dos normalizadores e densidade} that  $\overline{N}_{G_1}(H) = \overline{N}_{\widehat G_1}(H) \triangleleft \overline{N}_{\widehat G}(H)$ and therefore $\widetilde{N}_{G_1}(H)  \triangleleft  \widetilde{N}_{\widehat G}(H)$; moreover by Proposition \ref{cp1} we have $$\overline{N}_{\widehat G}(H) = \langle \overline{N}_{\widehat G_1}(H),  f\rangle= \langle \overline{N}_{G_1}(H),  f\rangle=\overline{N}_{G}(H).$$  


\begin{cor} \label{genus cotas faceis}
Let $H$ be a finite subgroup of an $\Oe$-group $G_1$ and $G=\hnn(G_1,H,f,t)$ be an HNN-extension, where $f \in \aut(H)$. Suppose that $\widetilde{\aut}_{\widehat{G}_1}(H)=\widetilde{\aut}_{G_1}(H)$, $\widetilde{N}_{\widehat G_1}(H)=\widetilde{N}_{ G_1}(H)$   and moreover,  one of the following conditions is satisfied \begin{enumerate}
    \item [i)] $f \in \overline{\aut}_{G_1}(H);$

    \item [ii)] $\overline{N}_{ G_1}(H)=\inn(H);$

    \item [iii)] $H \leq Z(G_1).$
\end{enumerate}  
Then \begin{enumerate}    \item [a)] $|g(G, \F)|=1,$ if $|\widetilde{N}_{ G}(H)/\widetilde{N}_{ G_1}(H)|\leq 2$;
    \item [b)] 
 $|g(G, \F)|\leq \frac{\phi\left(|\widetilde{N}_{G}(H)/\widetilde{N}_{G_1}(H)|\right)}{2},$ if $\left| \frac{\widetilde{N}_{ G}(H)}{\widetilde{N}_{ G_1}(H)}\right| \geq 3,$ where $\phi$ is the Euler function.
\end{enumerate} 

%Then \begin{enumerate}     \item [a)] $|g(G, \F)|=1,$ if $|\overline{N}_{ G}(H)/\overline{N}_{ G_1}(H)|\leq 2$;
 %   \item [b)] $|g(G, \F)|\leq \frac{1}{2} \phi\left(|\overline{N}_{G}(H)/\overline{N}_{G_1}(H)|\right),$ if $\left| \frac{\overline{N}_{ G}(H)}{\overline{N}_{ G_1}(H)}\right| \geq 3,$ where $\phi$ is the Euler function. 
%\end{enumerate}  
\end{cor}



\begin{proof} By Remarks \ref{rem dos normalizadores e densidade} and 
\ref{tilde implica overline} we have  $$\overline{N}_{\widehat G}(H) = \langle \overline{N}_{\widehat G_1}(H),  f\rangle= \langle \overline{N}_{G_1}(H),  f\rangle=\overline{N}_{G}(H)$$ and $\widetilde{N}_{\widehat G}(H) = \widetilde{N}_{G}(H)=\langle  \widetilde{N}_{G_1}(H),  \tilde{f}\rangle.$ Since $\overline{N}_{G_1}(H) \triangleleft \overline{\aut}_{G_1}(H),$ then for any of the hypotheses i)-iii) of this corollary we have that $\overline{N}_{G_1}(H) \triangleleft \overline{N}_{\widehat G}(H)=\overline{N}_{G}(H),$  whence it follows that $\widetilde{N}_{G_1}(H) \triangleleft \widetilde{N}_{\widehat G}(H)=\widetilde{N}_{G}(H).$ Therefore, the quotients $\overline{N}_{ G}(H)/\overline{N}_{ G_1}(H) \cong \widetilde{N}_{ G}(H)/\widetilde{N}_{ G_1}(H)$ are finite cyclic. Let $\tilde{\zeta}: \widetilde{N}_{G}(H) \longrightarrow \widetilde{N}_{ G}(H)/\widetilde{N}_{ G_1}(H)$ be the canonical epimorphism. By the definition of $R$ and $\mathcal{R}$ in (\ref{conjunto S com BZ}) and (\ref{cS}), we have that  $\tilde{\zeta}(\mathcal{R})$ is the set of all generators of the cyclic group $\widetilde{N}_{ G}(H)/\widetilde{N}_{ G_1}(H).$ Since $\widetilde{N}_{G_1}(H) \backslash \cR \subset \widetilde{N}_{G_1}(H) \backslash \widetilde{N}_{G}(H)$, we obtain the result by Corollary \ref{nova conta conjugadoss} by looking only at the orbits obtained by the action of $C_2$ on the generator set $\widetilde{N}_{G_1}(H) \backslash \cR.$
\end{proof}


%If hypotheses ii) or iii) are satisfied then $\widetilde{N}_{G_1}(H)=\widetilde{N}_{\widehat G_1}(H)=\{1\},$ on the other hand, if condition i) is satisfied, then since $\overline{N}_{G_1}(H) \triangleleft \overline{\aut}_{G_1}(H)$ it follows that $\widetilde{N}_{G_1}(H) \triangleleft \widetilde{\aut}_{G_1}(H)$ and therefore $\overline{N}_{G_1}(H) \triangleleft \overline{N}_{\widehat G}(H)=\overline{N}_{G}(H),$ whence it follows that $\widetilde{N}_{G_1}(H) \triangleleft \widetilde{N}_{\widehat G}(H)=\widetilde{N}_{G}(H).$

%\begin{proof}  
%The proof of it follows from the Corollary \ref{autG classe}, \textcolor{red}{Remarks \ref{rem dos normalizadores e densidade}, \ref{tilde implica overline}} and the versions of Corollary 5.2 and Theorem 5.4 of \cite{BZ} for $G_1$ an $\Oe$-group, %according to Remark \ref{super remark example}. 
%since for any of the hypotheses $i), ii) \mbox{ and  } iii)$ of Corollary \ref{genus cotas faceis}  we have $\overline{N}_{G_1}(H) \triangleleft \overline{N}_{G}(H),$  $\overline{N}_{\widehat G_1}(H) \triangleleft \overline{N}_{\widehat G}(H),$ and since $\overline{N}_{\widehat G_1}(H)=\overline{N}_{ G_1}(H)$ one deduces that $\overline{N}_{G_1}(H) \triangleleft \overline{N}_{\widehat G}(H),$ and therefore $\widetilde{N}_{G_1}(H)  \triangleleft  \widetilde{N}_{\widehat G}(H)$; \textcolor{red}{moreover by Proposition \ref{cp1} and Remark \ref{tilde implica overline}} we have $$\overline{N}_{\widehat G}(H) = \langle \overline{N}_{\widehat G_1}(H),  f\rangle= \langle \overline{N}_{G_1}(H),  f\rangle=\overline{N}_{G}(H) \mbox{ and } \widetilde{N}_{\widehat G}(H) = \widetilde{N}_{G}(H).$$  
%\end{proof}

\subsection{Genus 1}

In this subsection we will give some sufficient conditions for  the genus of HNN-extensions with  an $\Oe$ base group and finite associated subgroups to have only one element. Of course, if  $\widetilde{N}_{G_1}(H)=\out(H)$, then $|g(G, \mathcal{F})| =1
$ by Corollary %\ref{covered out(H)} 
\ref{normal iguais genero 1 non conjugated} and so we shall omit this condition from the list  below to avoid repeatition. The following results give a generalization of Theorem 1.1 in \cite{BZ}. 

%In any of the items below one has $\widetilde{N}_{\widehat G}(H)=\widetilde{N}_{G_1}(H)$ and so the result is an immediate consequence of the Corollary \ref{normal iguais genero 1 non conjugated}.

We start the case of normal HNN-extension (equivalently $H$ and $K$ are conjugate in the base group).

\begin{thm}\label{genus 1 normal}
%Let $H$ and $K$ be a finite subgroups of an $\Oe$-group $G_1$ 
Let $H$ be a finite subgroup of an $\Oe$-group $G_1$ and $$G=\hnn(G_1, H, f, t)$$ be a normal HNN-extension, i.e. $f \in \aut(H).$ 
Suppose that   $\widetilde{\aut}_{\widehat{G}_1}(H)=\widetilde{\aut}_{G_1}(H)$ and $\widetilde N_{G_1}(H)=\widetilde N_{\widehat G_1}(H)$. If $\widetilde{N}_{\widehat G}(H)=\widetilde{N}_{G_1}(H)$ (this is the case if $f\in \overline N_{G_1}(H)$), then $|g(G,\F)|=1$. 
\end{thm}
\begin{proof} Follows from Theorem \ref{1.8 da introducao}.
\end{proof}

%By Remarks \ref{rem dos normalizadores e densidade} and 
%\ref{tilde implica overline} we have  $$\overline{N}_{\widehat G}(H) = \langle \overline{N}_{\widehat G_1}(H),  f\rangle= \langle \overline{N}_{G_1}(H),  f\rangle=\overline{N}_{G}(H),$$ and so   $\widetilde{N}_{\widehat G}(H) = \widetilde{N}_{G}(H)=\langle  \widetilde{N}_{G_1}(H),  \tilde{f}\rangle.$

%%%%%%%%%%%%%%%%%%%%%%%%%%%%%%%%%%%%%%%%%%%%%%%%%%%%%%%%%%%%%%%%
%Now we state theorem for not normal HNN-extensions. Note that item (iii) of Theorem \ref{genus 1 intr} automatically enters in the next theorem, since if $H$ or $K$ is central in $G_1,$ then $H\neq K$ implies non-conjugacy of $H$ and $K$.

Now we state theorem for not 
normal HNN-extensions. Note that the particular case in item iv) of Theorem \ref{genus 1 intr} automatically enters in the next theorem,  since if $H$ or $K$ is central in $G_1,$ then $H\neq K$ implies non-conjugacy of $H$ and $K$.



%%%%%%%%%%%%%%%%%%%%%%%%%%%%%%%%%%%%%%%%%%%%%%%%%%%%%%%%%%%%%%%%
\begin{thm}\label{genus 1}
Let $H$ and $K$ be  finite non-conjugate subgroups of an $\Oe$-group $G_1$ and $G=\hnn(G_1, H, K, f, t)$ be an HNN-extension.
Suppose that $\widetilde{\aut}_{\widehat{G}_1}(H)=\widetilde{\aut}_{G_1}(H)$,   $[G_1:N_{G_1}(K)]=[\widehat G_1:N_{\widehat G_1}(K)]<\infty$ and 
     either $(H,K)$ is a swapping pair in $G_1$ or  not a swapping pair in $\widehat G_1$. 
 Then there is only one element in the genus of $G$ in $\F,$ if one of the following conditions is satisfied:
\begin{itemize} 
%\item[i)] $H$ and $K$ are conjugate in $G_1$;
\item[i)]   
     $f$ extends to an automorphism of $G_1;$    
    %% $f \in \aut(G_1);$
    \item [ii)] $N_{G_1}(K)=K \cdot C_{G_1}(K)$ or $N_{G_1}(H)=H \cdot C_{G_1}(H)$ (in particular,  $H$ or $K$ is central in $G_1$).  
   % \item [ii)] $H$ or $K$ is central in $G_1;$
\end{itemize}
\end{thm}

%\begin{proof} Items $i)$ and $iii)$ follow from Corollary \ref{normal iguais genero 1 non conjugated} since   $\widetilde{N}_{\widehat G}(H)=\widetilde{N}_{G_1}(H)$ .

%To see that  $\widetilde{N}_{\widehat G}(H)=\widetilde{N}_{G_1}(H)$ in (ii) one uses Remark  \ref{tilde implica overline} and Proposition \ref{cp1}.
%\end{proof}

\begin{proof}
First note that the hypothesis $[G_1:N_{G_1}(K)]=[\widehat G_1:N_{\widehat G_1}(K)]<\infty$ implies that $N_{G_1}(K)$ is dense in $N_{\widehat G_1}(K),$ and as $$N_{G_1}(H)=N_{G_1}(K^{t^{-1}})=\left(N_{G_1}(K)\right)^{t^{-1}}$$ it follows that $N_{G_1}(H)$ is dense in $N_{\widehat G_1}(H)$. Therefore  $\overline N_{G_1}(K) = \overline N_{\widehat G_1}(K)$ and $\overline N_{G_1}(H) = \overline N_{\widehat G_1}(H).$ Since $f(h)=h^t=k$ if and only if $f(h)^{t^{-1}}=h=k^{t^{-1}}, \forall h \in H, k \in K,$ one has $\inn(H)=f^{-1} \cdot \inn(K) \cdot f.$  By Proposition \ref{cp1} \begin{equation} \label{genus 1 prova}  \overline{N}_{\widehat G}(H) = \langle \overline{N}_{\widehat G_1}(H), \overline{N}_{\widehat G_1}(K)^{f} \rangle=\langle \overline{N}_{G_1}(H), \overline{N}_{G_1}(K)^{f} \rangle=\overline{N}_{G}(H)\end{equation} and so   $\widetilde{N}_{\widehat G}(H) = \widetilde{N}_{G}(H)=\langle  \widetilde{N}_{G_1}(H),  \widetilde{N}_{G_1}(K)^{\tilde f}\rangle.$ We need to find the image of $\overline{N}_{\widehat G_1}(H)$ and $\overline{N}_{\widehat G_1}(K)^{f}$ in $\aut(H).$ We will divide the rest of the proof into the following cases.

\textbf{Case 1:} Suppose that $f$ extends to an automorphism of $G_1.$ In this case, if $\hat{g}_1 \in N_{\widehat G_1}(K)$,  then $\tau_{\left( \hat{g}_1^{t^{-1}}\right)^{-1}}= f^{-1} \tau_{g_1^{-1}}f=\tau_{f^{-1}(g_1^{-1})}$ for some $g_1 \in N_{G_1}(K)$, since $\overline N_{G_1}(K) = \overline N_{\widehat G_1}(K)$ and since $f$ extends to $G_1$. It follows that $\tau_{f^{-1}(g_1^{-1})} \in \overline{N}_{G_1}(H).$ Therefore $\overline{N}_{\widehat G_1}(K)^{f} \subset \overline{N}_{G_1}(H),$ whence it follows by (\ref{genus 1 prova}) that $\widetilde{N}_{\widehat G}(H)=\widetilde{N}_{G_1}(H),$ thus the result follows from the Corollary \ref{normal iguais genero 1 non conjugated}. 

\textbf{Case 2:} Suppose $N_{G_1}(K)=K \cdot C_{G_1}(K).$ In this case, we have that $$\overline{N}_{G_1}(K)=\inn(K)=f\cdot \inn(H) \cdot f^{-1}=\overline{N}_{\widehat G_1}(K),$$  thus $\overline{N}_{\widehat G_1}(K)^{f} \subset \inn(H) \subset \overline{N}_{G_1}(H),$ and the result follows from (\ref{genus 1 prova}) and the Corollary \ref{normal iguais genero 1 non conjugated} since $\widetilde{N}_{\widehat G}(H)=\widetilde{N}_{G_1}(H).$ 

\textbf{Case 3:} Suppose $N_{G_1}(H)=H \cdot C_{G_1}(H).$ In this case $$\overline{N}_{G_1}(H)=\inn(H)=f^{-1}\cdot \inn(K) \cdot f\subset f^{-1} \cdot \overline{N}_{G_1}(K) \cdot f.$$ It follows from (\ref{genus 1 prova}) that $\overline{N}_{\widehat G}(H)=\overline{N}_{G_1}(K)^{ f}$ and therefore $\widetilde{N}_{\widehat G}(H)=\widetilde{N}_{G_1}(K)^{\tilde f}.$ Note that $R=f\overline{N}^+\subset \overline{N}_{ G_1}(K)f$ since $\overline{N}^+\subset \overline{N}_{\widehat G}(H)$ (see equations (\ref{cadeia de N mais etc}) and  (\ref{O conjunto S})). The result follows from Theorem \ref{nova conta non conjugadosss} since $|g(G, \F)|=|\kappa(R)| \leq |\kappa(\overline{N}_{ G_1}(K)f)|=1$ because any HNN-extension of the form $B=\hnn(G_1, H,\\ K, \tau_{g_1}f, t')$ with $g_1 \in N_{G_1}(K)$ are isomorphic by Theorem \ref{classnonnormalabs}. 
\end{proof}


%%%%%%%%%%%%%%%%%%%%%%%%%%%%%%%%%%%%%%%%%%%%%%%%%%%%%%%%%%%%%%%

The result below is an immediate consequence of item ii) of the previous result because if $G_1$ is finitely generated abelian then $\overline{\aut}_{G_1}(H)=\overline{\aut}_{\widehat{G_1}}(H)$ and the other conditions of the previous theorem are also satisfied. Note that if $G_1$ is virtually infinite cyclic and $H \neq tor(G_1)$ the result still follows from Theorem \ref{genus 1} ii) and Remark  \ref{finite-by-cyclic}. For $H=tor(G_1)$ we have that the HNN-extension is of the form $G=H \rtimes F,$ where $F$ is a free group and the genus in this case was calculated in Theorem 3.26 in \cite{GZ}. This result generalizes Corollary 3.3 in \cite{BZ}.

\begin{cor} \label{fin ger abelian} Suppose that $G_1$ is a finitely generated abelian group and $H\neq K$ its finite subgroups. Then
 $|g(\hnn(G_1, H, K, f, t), \mathcal{F})|=1$.  
\end{cor}

Recall that an indecomposable (into a free product) finitely generated Fuchsian group $\Gamma$ has the presentation of
the form: 
$$
\Gamma=\langle\,a_1,b_1,\ldots,a_n, b_n,\, c_1,\ldots, c_t\, \mid$$
$$ c_1^{e_1}=\ldots=
c_t^{e_t}=1,c_1^{-1}\cdot\ldots \cdot
c_t^{-1}\cdot[a_1,b_1]\cdot\ldots\cdot [a_n,b_n]=1 \, \rangle,$$  where $n,\,
t\geq 0$ and $e_j>1$ for $j=1,\ldots, t$ (see \cite{newman}).

\begin{exa}\label{grupo fuchsian}  Let $\Gamma$ be indecomposable into a free product finitely generated Fuchsian group and $H \neq \{1\}$ is one of its finite subgroup. Let $G_1=K\times \Gamma$ such that $H \cong K$. We first observe that indecomposable into a free product Fuchsian groups are virtually compact surface groups and so they are $\Oe$-groups by Proposition \ref{examples} (i). Therefore by Proposition \ref{extensao de OE group} $G_1$ is an $\Oe$-group. Consider $G=\hnn(G_1,H, K, f, t)$ where $f\in \Iso(H, K).$ Note that $H \neq K$ and therefore $G$ is not normal HNN-extension. By Corollary 7.4 of \cite{BPZ2}  $\overline{\aut}_{\Gamma}(H)= \aut(H)=\overline{\aut}_{\widehat \Gamma}(H)$ therefore we have $$\widetilde{\aut}_{\Gamma}(H)=\widetilde{\aut}_{G_1}(H)= \out(H)=\widetilde{\aut}_{\widehat \Gamma}(H)=\widetilde{\aut}_{\widehat G_1}(H).$$ Now note that $H$ is not a normal subgroup of $\widehat \Gamma$ by Corollary 5.2 in \cite{BCR16}, thus there is no $\hat \psi \in \aut(\widehat G_1)$ such that $\hat \psi(H)=K$ because $K$ is normal in $\widehat G_1.$ Furthermore, since $K \leq Z(\widehat G_1)$, it follows that $N_{G_1}(K)=G_1$ and $N_{\widehat G_1}(K)=\widehat G_1.$ Thus all the hypotheses of Theorem \ref{genus 1} are satisfied, and it follows from item ii) that $|g(G, \mathcal{F})|=1.$ 

%Therefore all hypotheses of Theorem \ref{genus 1} are satisfied since $K \leq Z(\widehat G_1)$ so $\tau_{\hat g_1}|_{_{K}}=id,\, \forall \hat g_1 \in \widehat G_1,$ \textcolor{red}{trocar essa última condição para: so $N_{\widehat{G_1}}(K)=\widehat G_1$ fica mais claro que satifaz a condição i)}  where it follows from item ii) that $|g(G, \mathcal{F})|=1.$

\end{exa}

%\begin{exa} Let $G_1=M\times H$ such that $M$ is an $\Oe$-group and $H$ is finite group. Consider $G=HNN(G_1,H,f,t)$ where $f\in Aut(H)$. Observe that $\widetilde{\aut}_{\widehat{G}_1}(H)=\widetilde{\aut}_{G_1}(H)$. Therefore it suffices to show that $\overline{Aut}_{G_1}(H)=Aut(H)$. Indeed, take $\alpha\in Aut(H)$ and observe that $\psi:=(id,\alpha) \in Aut(G_1)$, where $id$ is automorphism identity of $M$, i.e., $\psi|_H=\alpha$. Follows from the Theorem \ref{genus 1} item iv) that  $|g(G, \mathcal{F})|=1$, since $f\in Aut(G_1)$.
%\end{exa}


\begin{exa} Let $M$ be an $\Oe$-group and $H$ be a finite group such that $\out(H) $ is trivial (for example this occurs for symmetric groups $ S_n$ with $n \neq 6$ and Mathieu simple groups $M_{11}, M_{23}$ and $M_{24}$). %Let $M$ be a finitely generated residually finite torsion free group that is not virtually infinite cyclic such that $\hat M$ does not contain a non-abelian free pro-$p$ subgroup. 
Let $G_1=M\times H.$ From Proposition \ref{extensao de OE group} $G_1$ is an $\Oe$-group. Consider $G=\hnn(G_1,H,f,t)$ where $f\in \aut(H)$. %Observe that $\widetilde{\aut}_{\widehat{G}_1}(H)=\widetilde{\aut}_{G_1}(H)$ because $H$ is characteristic in $\hat G_1.$ 
Note that $\overline{\aut}_{G_1}(H)=\aut(H).$ Indeed, take $\alpha\in \aut(H)$ and observe that $\psi:=(id,\alpha) \in \aut(G_1)$. Thus $\overline{\aut}_{G_1}(H)=\aut(H)=\overline{\aut}_{\widehat G_1}(H)=\inn(H),$ and therefore we have $\widetilde{N}_{\widehat G}(H)=\widetilde{N}_{G_1}(H)=\widetilde{N}_{\widehat G_1}(H)=\widetilde\aut_{G_1}(H)=\widetilde{\aut}_{\widehat G_1}(H)=\out(H)=\{1\}.$ %from which it follows from Corollary \ref{normal iguais genero 1 non conjugated} that $|g(G, \mathcal{F})|=1.$ 
Thus all hypotheses of Theorem \ref{genus 1 normal} are satisfied, so $|g(G, \mathcal{F})|=1.$
\end{exa}

\begin{rem}
%\sout{The previous example is also an immediate consequence of the Corollary \ref{autG classe}, more precisely if $G=\hnn(G_1, H, f, t)$ is a normal abstract HNN-extension such that $\widetilde{\aut} _{\widehat{G}_1}(H)=\widetilde{\aut}_{G_1}(H),$ $\widetilde{N}_{\widehat G_1}(H)=\widetilde{N}_{ G_1}(H)$ (this occurs since $N_{G_1}(H)= G_1$  is dense in $N_{\widehat G_1}(H)=\widehat G_1$) and $|\out(H)| \leq 2$ then it follows from the Corollary \ref{autG classe} that $|g(G, \F)|=1.$}
In the previous example, if $G_1=M\times H,$ where $H=S_6,$ then $\overline{\aut}_{G_1}(H)=\aut(H)=\overline{\aut}_{\widehat G_1}(H)$ consequently $\widetilde{\aut}_{G_1}(H)=\out(H)=\widetilde{\aut}_{\widehat G_1}(H) \cong C_2$ and $\widetilde{N}_{\widehat G_1}(H)=\widetilde{N}_{ G_1}(H)$ (this occurs since $N_{G_1}(H)= G_1$  is dense in $N_{\widehat G_1}(H)=\widehat G_1$), therefore it follows from Corollary \ref{autG classe} that $|g(G, \mathcal{F})|=1.$ Similarly, the same result is concluded for the simple Mathieu groups $M_{12}$ and $M_{22}.$ 
\end{rem}

 
We finish this subsection with an example of HNN-extension where the base group is virtually infinite 
cyclic which is relevant by Remark \ref{finite-by-cyclic}.


\begin{exa} Let $H$ and $K$ be  isomorphic finite cyclic groups and $G_1= (H\rtimes_{\varphi}\mathbb{Z})\times K$ where  $\varphi:\mathbb{Z}\longrightarrow \aut(H)$ is a homomorphism satisfying   $\varphi(1)\neq id.$ Consider $G=\hnn(G_1,H,K,f,t)$ where $f\in \Iso(H,K)$. Observe that $\widetilde{\aut}_{G_1}(H)=\overline{\aut}_{G_1}(H)=\out(H)=\aut(H)=\overline{\aut}_{\widehat G_1}(H)=\widetilde{\aut}_{\widehat G_1}(H).$ Indeed, let $\beta$ be an automorphism of $H$. The map $\theta:G_1\longrightarrow G_1$ defined by $\theta|_{H}=\beta$, $\theta|_{K}=id$, $\theta|_{\mathbb{Z}}=id$ is an automorphism of $G_1$ that lifts $\beta$. Now there is no $\hat \psi\in \aut(\widehat G_1)$ such that $\hat \psi(H)=K$ since $K$ is central in $G_1$ but $H$ is not. Note that $H \neq K,$ thus it follows from Theorem \ref{genus 1} $ii)$ that $|g(G, \mathcal{F})|=1$.
 
 %Consider $\alpha\in \aut(H\times K)$ given for $\alpha:=(f,f^{-1}).$ The map $\psi:G_1\longrightarrow G_1$ defined by $\psi|_{H\times K}=\alpha$ and $\psi|_{\mathbb{Z}}=id$ is an automorphism of $G_1$ that lifts $\alpha$, furthermore, $\alpha|_{H}=f.$ 
\end{exa}

\section{Non-fixed associated subgroup} \label{non fixed}

Let $\A$ be the class of all finitely generated residually finite accessible groups whose vertex groups of its JSJ-decomposition are $\Oe$-groups (see  on p. 224 of \cite{BPZ} for more details). 

Let $H$ and $K$ be finite subgroups of an $\Oe$-group $G_1$ and $$G=\hnn(G_1, H,\ K, f, t)$$ be an HNN-extension. Let $B\in \A$ such that $\widehat G \cong \widehat B.$ Then by \cite[Theorem 1.1]{BPZ} $B$ is an HNN-extension of the form $\hnn(B_1, H_1, K_1, f_1, t_1)$ with $\widehat B_1\cong \widehat G_1, H \cong H_1\cong K \cong K_1,$ and furthermore $B_1 \in \mathcal{A}$. The class $\F(G,G_1,\A)$ is the subclass of all such $B \in \mathcal{A}$ with $B_1\cong G_1$. Thus we assume from now on that $B_1=G_1$, i.e. consider all HNN-extensions of the form $\hnn(G_1, H_1, K_1, f_1, t_1)$ whose associated  subgroups are finite. 

 We will denote the genus of $G$ in the class $\F(G,G_1,\A)$ as $g(G, G_1, \mathcal{A}).$ Our goal in this section is to find formulas and/or bounds for the cardinality of the genus $g(G, G_1, \mathcal{A})$ of a given HNN-extension with $G_1$ fixed in the class $\A$.

Recall that $S_H$ (respectively $P_H$) is the collection of all finite subgroups of $G_1$ isomorphic to $H$ (respectively subgroups of $\widehat{G}_1$ isomorphic to $H$).  Note that $G_1$ (resp. $\widehat{G}_1$) acts naturally on $S_H$ (resp. $P_H$) by conjugation and as defined in Subsections \ref{cota bonita non conjugados} and \ref{quotasprofnonconjugates}, respectively, denote by 
$$\overline{S}_H:=G_1 \backslash S_H\ \  (\textrm{resp.}\  \overline{P}_H:= \widehat{G}_1 \backslash P_H),$$ the orbit spaces of these actions (the sets of conjugacy classes of subgroups of $S_H$ (and of $P_H$ respectively)). %Let $\overline{S}_H=\{ \inn(G_1) \cdot L\,|\, L \in S_H\}$ and $\overline{P}_H=\{ \inn(\widehat{G} _1) \cdot L\,|\  L \in P_H\}$ the orbit spaces defined in Subsections \ref{cota bonita non conjugados} and \ref{quotasprofnonconjugates}, respectively. 

Denote by $[S_H]^2$ and $[P_H]^2$ the respective sets of unordered pairs of elements of $S_H$ and $P_H.$
The respective sets of unordered pairs of these orbits were defined as $[\overline{S}_H]^2$ and $[\overline{P}_H]^2.$ Furthermore, let us remember that $\aut(G_1)$ (resp. $\aut(\widehat{G}_1)$) acts naturally on $[\overline{S}_H]^2$ (resp. $[\overline{P}_H]^2$) (see (\ref{acao epsilon}) and (\ref{acao p})) (see Subsections \ref{cota bonita non conjugados} and \ref{quotasprofnonconjugates} for more details). 

 Their respective stabilizers will be denoted as $\aut_{G_1}(H_1, K_1)$ and $\aut_{\widehat G_1}(H_1,\\ K_1),$ respectively. Note that $\aut_{G_1}(H_1, K_1)$ is the subgroup of automorphisms of $G_1$ that leaves the conjugacy classes of $H_1$ and $K_1$ in $G_1$ invariant 
 up to a swap %=\{\psi\in \aut( G_1) \ | \ \psi(H_1)=H_1^{g_1} \mbox{ and } \psi(K_1)=K_1^{g_2} \mbox{ or } \psi(H_1)=K_1^{g_1} \mbox{ and } \psi(K_1)=H_1^{g_2}, \mbox{ where } g_i \in G_1,\\ (i=1,2) \}$ 
 and similarly $\aut_{\widehat G_1}(H_1, K_1)$ is the subgroup of automorphisms of $\widehat G_1$ that leaves the conjugacy classes of $H_1$ and $K_1$ in $\widehat G_1$ invariant 
 up to a swap.
 % =\{\hat\psi\in \aut(\widehat G_1) \ | \ \hat\psi(H_1)=H_1^{g_1} \mbox{ and } \hat\psi(K_1)\\=K_1^{g_2} \mbox{ or }\hat\psi(H_1)=K_1^{g_1} \mbox{ and } \hat\psi(K_1)=H_1^{g_2}, \mbox{ where } g_i \in \widehat G_1, (i=1,2)\}$.


%As previously discussed $\{H_1, K_1\}$ and $\{H_2, K_2\}$ represent the same $\aut(G_1)$-orbit according to action (\ref{acao epsilon}) (respectively $\aut(\widehat{G}_1)$-orbit according to action (\ref{acao p})) if and only if (respectively $\{\inn (\widehat{G}_1) \cdot H_1, \inn(\widehat{G}_1) \cdot K_1 \}$ and $\{\inn(\widehat{G}_1) \cdot H_2, \inn(\widehat{G}_1) \cdot K_2 \}$ are in the same  $\aut(\widehat{G}_1)$-orbits according to action (\ref{acao p})). 






%Denote by $\aut(G_1)\{H_1,K_1\}$ the orbit of $\{\inn (G_1) \cdot H_1, \inn(G_1) \cdot K_1 \}$ with respect to action $(\ref{acao epsilon})$ and $\aut(\widehat{G}_1)\{H_1,K_1\}$ the orbit of $\{\inn(\widehat{G}_1) \cdot H_1, \inn(\widehat{G}_1) \cdot K_1\}$ with respect to action $(\ref{acao p}).$ 
 



%Also remember that $\aut(G_1)\{H_1,K_1\}$ and $\aut(\widehat{G_1})\{H_1,K_1\}$ are the orbits of $\{H_1,K_1\}$  with respect to these actions, respectively. The stabilizer of $\{H_1, K_1\} \in S_H \times S_H$ according to action (\ref{acao epsilon}) will be denoted by $\aut_{G_1}(H_1, K_1),$ and  similarly $\aut_{\widehat G_1}(H_1, K_1)$ will be the stabilizer of $\{H_1, K_1\}\\ \in P_H \times P_H$ according to action (\ref{acao p}). 

%Note that $\aut_{G_1}(H_1, K_1)=\{ \{\psi\in \aut( G_1) \ | \ \psi(H)=H^{g_1} \mbox{ and } \psi(K)=K^{g_2} \mbox{ or } \psi(H)=K^{g_1} \mbox{ and } \psi(K)=H^{g_2}, \mbox{ where } g_i \in G_1, (i=1,2)\} \}$ and $\aut_{\widehat G_1}(H, K)=\{\hat\psi\in \aut(\widehat G_1) \ | \ \hat\psi(H)=H^{g_1} \mbox{ and } \hat\psi(K)=K^{g_2} \mbox{ or }\hat\psi(H)=K^{g_1} \mbox{ and } \hat\psi(K)=H^{g_2}, \mbox{ where } g_i \in \widehat G_1, (i=1,2)\}$. We denote by $k(G_1,H_1,K_1)$  the number of $\aut(G_1)$-orbits of $(S_H\times S_H)\cap \aut(\widehat G_1)\{H_1,K_1\}$.

%When $H_1=K_1$ (which is the case for associated subgroups conjugated in the base group), we will use more compact notations as discussed in Subsections \ref{cota bonita non conjugados} and \ref{quotasprofnonconjugates}, namely $H_1 := \{H_1, H_1\},$ $k(G_1, H_1):=k(G_1, H_1, H_1),$ \begin{equation} \label{estabilizadores}  \aut_{G_1}(H_1):=\aut_{G_1}(H_1, H_1) \mbox{ and }  \aut_{\widehat G_1}(H_1):=\aut_{\widehat G_1}(H_1, H_1).\end{equation} 


%\begin{rem} \label{non normal com non normal genus}
%From Proposition \ref{normal non iso nao normal} we know that an (profinite) abstract HNN-extension with associated finite subgroups conjugated in the base group (a normal (profinite) abstract HNN-extension) cannot be isomorphic to a non-normal (profinite) abstract HNN-extension. So when calculating the cardinality of genus of $G=\hnn(G_1, H, K, f)$ in $(G_1,\A),$ if we have $H$ and $K$ conjugated in the base group, it will be enough to consider in the calculation of the genus, only groups $B \in \mathcal{NF}$ that are of the form $B=\hnn(G_1, H_1, K_1, f_1)$ such that $H_1$ and $K_1$ are conjugated in the base group, and vice versa. 
%\end{rem}

%\begin{rem} \label{nilpotent} 
%Suppose $G=\hnn(G_1, H, f, t)$ is a normal HNN-extension. If $\aut(\widehat G_1)$-orbit of $H$ is finite,
%then  we have that the number of $\aut(G_1)$-orbits of $\aut(\widehat G_1)H$ is 
%$|\aut(G_1)\backslash \aut(\widehat{G_1})/\aut_{\widehat{G_1}}(H)| < \infty$. 
%So an estimate for $k(G_1,H)$ is given by 
%$$k(G_1,H)\leq |\aut(G_1)\backslash \aut(\widehat G_1)/\aut_{\widehat{G_1}}(H)|.$$
%Furthermore, if $S=\widehat S$ then the inequality becomes an equality. This is true for example if $G_1$ is finitely generated nilpotent because $tor(G_1)$ is finite  and $tor(G_1)=tor(\widehat{G_1})$ (see Corollary 4.7.9 in \cite{RZ}).  
%\end{rem}

\bigskip
 Let $G=\hnn(G_1, H, K, f, t)$ be an abstract HNN-extension. Denote by $\aut(\widehat G_1)\{H,K\}  \subset [\overline{P}_H]^2$ the $\aut(\widehat G_1)$-orbit of $\{ \langle\langle H\rangle\rangle^{\widehat G_1} , \langle\langle K\rangle\rangle^{\widehat G_1} \}$ (see Section \ref{quotasprofnonconjugates} for the notation). Note that $\aut(G_1)$ (seen as a subgroup of $\aut(\widehat G_1)$) acts on $\aut(\widehat G_1)\{H,K\} $. If $\aut(\widehat G_1)\{H,K\}$ is finite, then   the number of $\aut(G_1)$-orbits of $\aut(\widehat G_1)\{H,K\}$ is
$$|\aut(G_1)\backslash \aut(\widehat G_1)/\aut_{\widehat G_1}(H, K)| < \infty.$$
 
 We denote by $\mathscr{L}_{_{HK}}$ the subset of $\aut(G_1)$-orbits of $\aut(\widehat G_1)\{H,K\}$  which has representatives in $[S_H]^2,$ i.e, satisfies the following property: $$\{ \langle\langle H_1 \rangle\rangle^{\widehat G_1} , \langle\langle K_1\rangle\rangle^{\widehat G_1} \} \in \mathscr{L}_{_{HK}}$$ if and only if  there exists $g_i\in \widehat G_1 (i=1,2)$ such that $\{H_1^{g_1}, K_1^{g_2}\} \in [S_H]^2$. 
 
We denoted by $k(G_1,H,K)$ the cardinality of set  $\mathscr{L}_{_{HK}}$. The estimate for $k(G_1,H,K)$ is given by 
$$k(G_1,H,K)\leq |\aut(G_1)\backslash \aut(\widehat G_1)/\aut_{\widehat G_1}(H, K)|.$$
Furthermore, if $S=P$ (in particular $S_H = P_H$)  then the inequality becomes an equality. This is true for example if $G_1$ is finitely generated nilpotent, because $tor(G_1)$ is finite  and $tor(G_1)=tor(\widehat{G}_1)$ (see Corollary 4.7.9 in \cite{RZ}).



%\begin{rem} Let $G=\hnn(G_1, H, K, f, t)$ be an abstract HNN-extension. Denote by $ \mathscr{L} \subset \overline{P}_H\times\overline{P}_H$ the $\aut(\widehat G_1)$-orbit of $\{\inn(\widehat{G}_1) \cdot H, \inn(\widehat{G}_1) \cdot K\}$. Note that $\aut(G_1)$ (seen as a subgroup of $\aut(\widehat G_1)$) acts on $\mathscr{L}$. If $\mathscr{L}$ is finite, then  we have that the number of $\aut(G_1)$-orbits of $\mathscr{L}$ is
%$$|\aut(G_1)\backslash \aut(\widehat G_1)/\aut_{\widehat G_1}(H, K)| < \infty.$$
 %We denote by $\aut(\widehat G_1)\{H,K\}$ the subset of $\aut(G_1)$-orbits of $\mathscr{L}$  which has representatives in $S_H \times S_H,$ i.e, satisfies the following property: $\{\inn( \widehat G_1) \cdot H_1, \inn( \widehat G_1) \cdot K_1\} \in \aut(\widehat G_1)\{H,K\}$ if and only if  there exists $g_i\in \widehat G_1 (i=1,2)$ such that $\{H_1^{g_1}, K_1^{g_2}\} \in S_H\times S_H$. \textcolor{red}{Essa condição é fundamental pra gente, ela garante que só as orbitas  (veja que elas são elementos de $\overline{P}_H\times \overline{P}_H$) que tem reresentantes em $S_H\times S_H$ serão consideradas. Do contrario, você estaria contabilizando uma orbita da forma $\{H_1\inn(\widehat G_1), H_2\inn(\widehat G_1)\}$ onde nenhum dos subgrupos da forma $H_1^{g_1}$ e $K_1^{g_2}$ são subgrupos de $G_1$.  Não teria como construir uma hnn abstrata $\hnn(G_1, H_1^{g_1}, K_1^{g_2}, f,t)$ cuja completamento é $\widehat G$.}

 %\textcolor{blue}{No teorema 7.18 iremos considerar a variação do somatorio nos representantes $\{H_1,K_1\}\in S_H\times S_H$ das $\aut(\widehat G_1)\{H,K\}$ órbitas. Como esse conjunto foi construido esse representante sempre vai existir. Aqui nem precissamos fazer mais a interseção.}
 %The cardinality of set  $\aut(\widehat G_1)\{H,K\}$  will be denoted by $k(G_1,H,K)$. So an estimate for $k(G_1,H,K)$ is given by 
%$$k(G_1,H,K)\leq |\aut(G_1)\backslash \aut(\widehat G_1)/\aut_{\widehat G_1}(H, K)|.$$
%Furthermore, if $S_H = P_H$ then the inequality becomes an equality. This is true for example if $G_1$ is finitely generated nilpotent because $tor(G_1)$ is finite  and $tor(G_1)=tor(\widehat{G}_1)$ (see Corollary 4.7.9 in \cite{RZ}).
%\end{rem}


%%%%%%%%%%%%%%%%%%%%%%%%%%%%%%%%%%%
%%%%%%%%%%%%%%%%%%%%%%%%%%%%%%%%%%%
%\begin{rem}  
%Let $G=\hnn(G_1, H, K, f, t)$ be an abstract HNN-extension. If $\aut(\widehat G_1)$-orbit of $\{\inn(\widehat{G}_1) \cdot H, \inn(\widehat{G}_1) \cdot K\}$ is finite, then  we have that the number of $\aut(G_1)$-orbits of $\aut(\widehat G_1)\{H, K\}$ is
%$$|\aut(G_1)\backslash \aut(\widehat G_1)/\aut_{\widehat G_1}(H, K)| < \infty.$$ So an estimate for $k(G_1,H,K)$ is given by 
%$$k(G_1,H,K)\leq |\aut(G_1)\backslash \aut(\widehat G_1)/\aut_{\widehat G_1}(H, K)|.$$
%Furthermore, if $S_H = P_H$ then the inequality becomes an equality. This is true for example if $G_1$ is finitely generated nilpotent because $tor(G_1)$ is finite  and $tor(G_1)=tor(\widehat{G}_1)$ (see Corollary 4.7.9 in \cite{RZ}).  
%\end{rem}


\begin{rem} \label{os ls}
Note that $\aut(G_1)$ acts naturally on $S_H$ by $\psi \cdot L=\psi(L), \, \forall L \in S_H \mbox{ and } \psi \in \aut( G_1).$ Denote by $\aut(G_1) \cdot L$ the orbit of $L$ with respect to this action, similarly write $\aut(\widehat{G}_1) \cdot L$ for the orbit of the respective action when $L$ is in $P_H.$ Let $l_{H}, l_{K}$ be the cardinality of the $\aut(G_1)$-orbits of $S_H \cap (\aut(\widehat{G}_1) \cdot H)$ and $S_H \cap (\aut(\widehat{G}_1) \cdot K)$, respectively. Then $k(G_1, H, K) \leq l_{H} \cdot l_{K}.$ So when $l_H=l_K=1$ we have  $k(G_1, H, K)=1.$       
\end{rem}

\begin{rem} Note that if $G_1$ has only finitely many conjugacy classes of finite subgroups, then  $k(G_1,H,K)$ is finite. Arithmetic groups   \cite{borelarith} and many groups of geometric nature as well as  mapping class groups of any surface of finite type \cite[Theorem 6]{Brid}, the automorphism group and the outer automorphism group of a free group of finite rank (see \cite{culler}, \cite{zimm} and p. 19 in \cite{vog}),  hyperbolic groups   \cite[Theorem 3.2, p. 459]{bridcurvature},  
groups that act properly and cocompactly by isometries on a $CAT(0)$ space \cite[Corollary 2.8 (2), p. 179]{bridcurvature}, share this property.    
\end{rem}

Now  we can state the  theorem which gives the precise cardinality for the genus $g(G,G_1,\A)$, where $G=\hnn(G_1, H, K, f, t)$. If one of the indices in the following summation is infinite or if any of the summands is infinite, then we say that the genus is infinite.

\begin{thm}\label{genus hnn nao fixos} Let $G_1$ be an $\Oe$-group and $G=\hnn(G_1, H, K, f, t)$ be an HNN-extension with finite associated subgroups. Then the cardinality of genus $g(G,G_1,\A)$ of $G$  equals to $$|g(G,G_1,\A)|=\displaystyle\sum_{\{H_1,K_1\}} \left| g\left(\hnn(G_1, H_1, K_1, f_1, t_1), \F\right)\right|,$$ where $\{H_1,K_1\}$  ranges over unordered pairs of representatives of  $\mathscr{L}_{_{HK}}.$ 

%$\aut(G_1)$-orbits of $\mathscr{L}_{_{HK}}.$
%$\textcolor{red}{(\overline{S}_H\times \overline{S}_H)}\cap \aut(\widehat G_1)\{H,K\}$. 
\end{thm}
%$\mathscr{L}_{_{HK}}$

%\begin{thm}\label{genus hnn nao fixos} Let $G_1$ be an $\Oe$-group and $G=\hnn(G_1, H, K, f, t)$ be an HNN-extension with finite associated subgroups. Then the cardinality of genus $g(G,\mathcal{NF})$ of $G$  equals to $$|g(G,\mathcal{NF})|=\displaystyle\sum_{\{\inn(G_1) \cdot H_1, \inn(G_1) \cdot K_1\}} \left| g\left(\hnn(G_1, H_1, K_1, f_1, t_1), \F\right)\right|,$$ where $\{\inn(G_1) \cdot H_1, \inn(G_1) \cdot K_1\}$ ranges over unordered pairs of representatives of $\aut(G_1)$-orbits of $\textcolor{red}{(\overline{S}_H\times \overline{S}_H)}\cap \aut(\widehat G_1)\{\inn(\widehat{G}_1) \cdot H, \inn(\widehat{G}_1) \cdot K\}$. 
%\end{thm}

%Let $\mathcal{B}$ be the family of all finitely generated residually finite not one-ended groups and denote by $g(G,\mathcal{B})$ the genus of the group $G=\hnn(G_1, H, K, f, t)$ in this  family. 

%\begin{cor}\label{classe b} Suppose $G_1$ is a profinitely rigid $\Oe$-group. Then the statement of Theorem \ref{genus hnn nao fixos} is valid for $g(G,\mathcal{B})$, where $\mathcal{B}$ is the family of all finitely generated residually
%finite not one-ended groups.
%\end{cor}


%\begin{proof}
%Let $B \in \B$ be such that $\widehat{B}\cong \widehat{G}$. By Corollaries \ref{class B} and  \ref{rigidez} $B=\hnn(B_{1}, H_1, K_1, f_1, t_1)$ such that $B_1\cong G_1$ and $H\cong H_1 \cong K \cong K_1.$ Therefore the result follows from Theorem \ref{genus hnn nao fixos}.  
%\end{proof} 




\begin{exa} Let $G_1$ be a finitely generated abelian group and $H\neq K$ its finite subgroups. It follows from Corollary \ref{fin ger abelian} that 
 $$|g(\hnn(G_1, H, K, f, t), G_1,  \mathcal{A})|=1.$$ Note if $G_1$ is virtually infinite cyclic with $H \neq tor(G_1) \neq K$ the result still follows by Remarks \ref{finite-by-cyclic}. 
 This result generalizes Corollary 3.3 in \cite{BZ}. 
\end{exa} 


\begin{rem} \label{ki finitos}
If $k(G_1,H,K)=1$ then it follows from Theorems \ref{fixed H} and \ref{genus hnn nao fixos} that $|g(G, \F)|=|g(G, G_1,\A)|=\left|\overline{\Gamma}_{HK}  \backslash  \Omega  \right|$, where $\Omega$ was defined in (\ref{Omega como uniao}). This happens for example when  $G_1$ is finite or $H$ and $K$ are characteristic in $\widehat G_1$. 
\end{rem} 


\begin{exa} \label{usar genus class A}
 Let $\Gamma$ be indecomposable into a free product finitely generated of Fuchsian groups and $H \neq \{1\}$ is one of its finite subgroup. Let $G_1=K\times \Gamma$ such that $H \cong K$. As in the example \ref{grupo fuchsian} $G_1$ is an $\Oe$-group. Consider $G=\hnn(G_1,H, K, f, t)$. It follows from the fact that $K$ is a  characteristic subgroup of $G_1$ (also of $\widehat{G}_1$) and from the calculation of the orbits made in Corollary 7.4 of \cite{BPZ2} that $k(G_1,H,K)=1$ (see Remark \ref{os ls}). It follows from the Example \ref{grupo fuchsian} and Remark \ref{ki finitos} that $|g(G, G_1,\A)|=|g(G, \F)|=1$. 
\end{exa}

\begin{exa} Let $\Gamma$ be an indecomposable into a free product finitely generated  Fuchsian group and $H \neq \{1\}$ a finite abelian group. Let $G_1=H\times \Gamma$ and $G=\hnn(G_1, H, f,t)$ a normal HNN-extension. As in the previous example  $H$ is central and characteristic in $\widehat{G}_1$ and in addition $k(G_1, H, H)=1.$ Let $\alpha \in \aut(H)$ and consider the following automorphism $\beta:G_1 \longrightarrow G_1$ that satisfies $\beta|_{_{H}}=\alpha$ and $\beta|_{_{ \Gamma}}=id,$ so we have $\alpha = \beta|_{_{H}} \in \overline{\aut}_{G_1}(H),$ from where it follows that $\widetilde{\aut}_{G_1}(H)= \out(H)=\widetilde{\aut}_{\widehat G_1}(H).$ As $H$ is central it follows that $\widetilde{N}_{\widehat{G}_1}(H)=\{1\}=\widetilde{N}_{G_1}(H).$ It follows from Remark \ref{ki finitos} and Corollary \ref{cotas n} that $$|g(G, G_1,\A)|=|g(G, \mathcal{F})| \leq \left\{\begin{array}{cc}
  1,&\mbox{ if } |\tilde f| \leq 2,\\
\displaystyle\frac{\phi(|\tilde f|)}{2}, &\mbox{ if } |\tilde f| \geq 3,
\end{array}\right.$$
where $\phi$ is Euler's function and $\tilde f$ is the image of $f$ in $\out(H)$. Furthermore, if $\aut(H)$ is abelian, the inequality becomes an equality in the formula above.
\end{exa}

\begin{exa}\label{free by Cp} Let $H$ be a group of order $3$ and $M$ a  $\Z H$-lattice of $\Z$-rank $2$ with non-trivial action. Let $G_1=M\rtimes_{\varphi} H$ be the corresponding semidirect product  and $G=\hnn(G_1,H,f,t)$ with $f\in \aut(H)$. Then $G_1$ is an $\Oe$-group (by Proposition 3.1 in \cite{BPZ2}) and has just two subgroups of order $3$ up to conjugation that can be swapped by an automorphism, so $k(G_1,H,H)=1$. 
Note that $\widehat G_1=\widehat M\rtimes_{\hat \varphi} H,$ where $\hat \varphi$ is the natural extension of $\varphi$ (see Proposition 2.6 in \cite{GZ}) and $\overline{N}_{\widehat G_1}(H) \leq \inn(H)=1$ since the action of $H$ is non-trivial, so $\widetilde{N}_{\widehat G_1}(H) = \{1\}=\widetilde{N}_{G_1}(H).$ By Lemma 2.1 \cite{Die}  $\overline{\aut}_{ G_1}(H)$ contains  $\overline N_{\aut(M)}(H)$. Besides, $\aut(M)\cong GL_2(\Z)$ contains $S_3$ and a subgroup of order $3$ is unique up to conjugation in $GL_2(\Z)$. Note that $\overline{\aut}_{G_1}(H)=\aut(H).$ Indeed, take $\alpha\in \aut(H)$ and observe that $\psi:=(id,\alpha) \in \aut(G_1)$ and $\psi|_H=\alpha \in \overline{\aut}_{G_1}(H)$. So $\widetilde{\aut}_{G_1}(H)=\widetilde{\aut}_{\widehat{G}_1}(H)=\aut(H)=\out(H)\cong C_2$. Then by Remark \ref{ki finitos} and Corollary \ref{autG classe} $|g(G,G_1,\A)|=|g(G, \mathcal{F})|=1.$ 
If the action of $H$ is trivial then $H$ is normal in $\widehat{G}_1$ and since $\aut(H) \cong C_2$, its follows that $\widetilde{N}_{G_1}(H)=\widetilde{N}_{\widehat G_1}(H)=\{1\},$  and the genus is 1. 

%If the action of $H$ is trivial then $H$ is normal in $\widehat{G}_1$ and thus $\overline{N}_{G_1}(H)=\overline{N}_{\widehat G_1}(H)$ and therefore $\widetilde{N}_{G_1}(H)=\widetilde{N}_{\widehat G_1}(H).$ And as in the previous case we also have that $|g(G,G_1,\A)|=|g(G, \mathcal{F})|=1.$
\end{exa}

%%%%%%%%%%%%%%%%%

%\begin{cor} Let $G_1$ be an indecomposable into a free product finitely generated Fuchsian group and $G=\hnn(G_1, H, f, t)$ a normal HNN-extension, where $H$ is a finite subgroup of $G_1$. Then $|g(G,G_1,\A)|=|\widetilde{N}_{G_1}(H)\backslash \cR/C_2|,$ where $ %R:=\{f\tau_{n^{-1}}\,|\, n \in N^{+}\}$.
%\cR=
%\{\tilde f \tilde \tau_{n^{-1}}\,|\, n \in N^{+}\}$.
 % In particular, if $H$ is maximal then the cardinality of the genus satisfies $$|g(G,G_1,\A)|=|\mathcal{R}/C_2| \leq \left\{\begin{array}{cc}
 %\displaystyle\frac{n}{2}+1,&\mbox{ if } n=|f| \mbox{ is even, }\\
%\displaystyle\frac{n-1}{2}+1 , &\mbox{ if } n=|f| \mbox{ is odd.}
%\end{array}\right.$$ 
%\end{cor} 

%\begin{proof} 
%By Proposition 3.3 (i) in \cite{BPZ2} $G_1$ is $\Oe$-group since that indecomposable into a free product Fuchsian groups are virtually compact surface groups. By Corollary 7.4 (see proof in \cite{BPZ2}) we have that $l_H=1$ then $k(G_1, H, H)=1$ by Remark \ref{os ls}. On the other hand, it follows from the presentation of a Fuchsian group  any automorphism of $H$ extends to an automorphism of $G_1$, i.e., $\widetilde{\aut}_{G_1}(H)=\out(H)=\widetilde{\aut}_{\widehat{G}_1}(H)$. Therefore, by Corollary \ref{nova conta conjugadoss} and Theorem \ref{genus hnn nao fixos}   $|g(G,G_1,\A)|=|g(G,\mathcal{F})|= |\widetilde{N}_{G_1}(H)\backslash\mathcal{R}/C_2|.$  If $H$ is a maximal finite subgroup then $N_{G_1}(H)=H$ and by Theorem 5.1 in \cite{BCR16} we have $\overline{N}_{G_1}(H)=\overline{N}_{\widehat G_1}(H)$ and therefore $\widetilde{N}_{G_1}(H)=\{1\}=\widetilde{N}_{\widehat{G}_1}(H).$ It follows from the Corollary \ref{cotas n} that $$|g(G, G_1,\A)|=|g(G, \mathcal{F})| \leq \left\{\begin{array}{cc}
 %\displaystyle\frac{n}{2}+1,&\mbox{ if } n=|f| \mbox{ is even, }\\
%\displaystyle\frac{n-1}{2}+1 , &\mbox{ if } n=|f| \mbox{ is odd.}
%\end{array}\right.$$ 

%\end{proof}

\begin{cor} Let $G_1$ be an indecomposable into a free product finitely generated Fuchsian group and $G=\hnn(G_1, H, f, t)$ a normal HNN-extension, where $H$ is a finite subgroup of $G_1$. Then $|g(G,G_1,\A)|=|\widetilde{N}_{G_1}(H)\backslash \cR/C_2|,$ where  %R:=\{f\tau_{n^{-1}}\,|\, n \in N^{+}\}$.
$\cR=
\{\tilde f \tilde \tau_{n^{-1}}\,|\, n \in N^{+}\}$.
  In particular, if $H$ is maximal then the cardinality of the genus satisfies $$|g(G,G_1,\A)|=|\mathcal{R}/C_2| = \left\{\begin{array}{cc}
  1,&\mbox{ if } |f| \leq 2,\\
\displaystyle\frac{\phi(|f|)}{2}, &\mbox{ if } | f| \geq 3,
\end{array}\right.$$
where $\phi$ is Euler's function. %and $\tilde f$ is the image of $f$ in $\out(H)$. 
\end{cor} 

%\begin{cor} Let $G_1$ be an indecomposable into a free product finitely generated Fuchsian group and $G=\hnn(G_1, H, f, t)$ a normal HNN-extension, where $H$ is a finite subgroup of $G_1$. Then $|g(G,G_1,\A)|=|\widetilde{N}_{G_1}(H)\backslash \cR/C_2|,$ where $ %R:=\{f\tau_{n^{-1}}\,|\, n \in N^{+}\}$.
%\cR=
%\{\tilde f \tilde \tau_{n^{-1}}\,|\, n \in N^{+}\}$.
 % In particular, if $H$ is maximal then the cardinality of the genus satisfies $$|g(G,G_1,\A)|=|\mathcal{R}/C_2| \leq \left\{\begin{array}{cc}
 %\displaystyle\frac{n}{2}+1,&\mbox{ if } n=|f| \mbox{ is even, }\\
%\displaystyle\frac{n-1}{2}+1 , &\mbox{ if } n=|f| \mbox{ is odd.}
%\end{array}\right.$$ 
%\end{cor}

\begin{proof} 
By Proposition 3.3 (i) in \cite{BPZ2} $G_1$ is an $\Oe$-group since indecomposable into a free product Fuchsian groups are virtually compact surface groups. By Corollary 7.4  in \cite{BPZ2} (see also the proof) we have  $l_H=1$ and so $$k(G_1, H, H)=1$$ by Remark \ref{os ls}. On the other hand, it follows from the presentation of a Fuchsian group that any automorphism of $H$ extends to an automorphism of $G_1$, i.e., $\widetilde{\aut}_{G_1}(H)=\out(H)=\widetilde{\aut}_{\widehat{G}_1}(H)$. Furthermore, $G_1$ is conjugacy separable  (see \cite{fine}), so $\overline{N}_{\widehat{G_1}}(H)=\overline{N}_{G_1}(H)$ and consequently $$\widetilde{N}_{\widehat{G_1}}(H)=\widetilde{N}_{G_1}(H).$$ Note that $\out(H)$ is abelian since $H$ is cyclic. Therefore, by Corollary \ref{nova conta conjugadoss} %Theorems \ref{1.8 da introducao} 
and Theorem \ref{genus hnn nao fixos} $|g(G,G_1,\A)|=|g(G,\mathcal{F})|= |\widetilde{N}_{\widehat G_1}(H)\backslash\mathcal{R}/C_2|=|\widetilde{N}_{G_1}(H)\backslash \cR/C_2|.$ If $H$ is a maximal finite subgroup then $N_{G_1}(H)=H$ and therefore $\widetilde{N}_{G_1}(H)=\{1\}=\widetilde{N}_{\widehat{G}_1}(H).$ The formula for the genus in this case follows from Corollary \ref{cotas n} and Remark \ref{H subgroup cyclic}.
\end{proof}

\begin{rem}
With hypotheses of the previous corollary, assume that $G=\hnn(G_1, H, K, f,t)$ where $H, K$ are non-conjugate finite subgroups of $G_1.$ Note that $|g(G,G_1,\A)|=|g(G,\mathcal{F})|= \left|\overline{\Gamma}_{HK}  \backslash \Omega\right|\,\,(\mbox{ see formula  } (\ref{Omega como uniao}))$ by  Theorem \ref{fixed H}. If $H$ is maximal, then $K$ is also maximal and they are self-normalized, from where it follows that $$\widetilde{N}_{G_1}(H)=\{1\}=\widetilde{N}_{G_1} (K)=\widetilde{N}_{\widehat{G}_1}(H)=\widetilde{N}_{\widehat{G}_1}(K).$$ Therefore by Proposition \ref{cp1} $\widetilde{N}_{G}(H)=\{1\}=\widetilde{N}_{\widehat{G}}(H)=\widetilde{N}^{+}.$ In this case the genus is simply given by $|g(G,G_1,\A)|=|g(G,\mathcal{F})|= \left|\overline{\Gamma}_{HK} \backslash \Omega\right|,$ where $\Omega=\overline{\Gamma}_{_{\widehat{HK}}} \cdot f.$  %\sout{ So here it also $\leq \dfrac{\phi(|\tilde f|)+1}{2}.$}
\end{rem}


From Theorems \ref{genus hnn nao fixos} and \ref{1.8 da introducao} we deduce the following

\begin{prop} \label{genus classF}
Let $H$ be a finite subgroup of an $\Oe$-group $G_1$ and $$G=\hnn(G_1, H, f, t)$$ be a normal HNN-extension. Then the cardinality of genus $g(G, G_1,\A)$ satisfies $$ k(G_1, H, H) \leq |g(G, G_1,\A)| %\leq k(G_1,H) \cdot |g(G, \mathcal{F})| 
\leq k(G_1, H, H) \cdot |\out(H)|.$$ In particular, if $|\out(H)|=1$ then $|g(G, G_1,\A)|=k(G_1, H, H).$ % when the associated subgroups are not conjugate in the base group and $$ k_{G_1}^{H} \leq |g(G, (G_1,\A))|\leq k_{G_1}^{H} \cdot |g(G, \mathcal{F})| \leq k_{G_1}^{H}  \cdot |Out(H)|,$$ when they are conjugate (in particular for HNN-normal extensions).
\end{prop}




%The next corollary shows that that in general genus is not 1. 

%\begin{cor}\label{in normalizer}  Suppose the  $f$ belongs to the $Aut(G)$-conjugacy class of  $\overline N_{G}(H)$.  Then $HNN(G,H,f,t)\cong HNN(G,H,f',t')$, where $t'$ centralizes $H$ and $g=k_H$. In particular, this is the case if $f\in \overline N_{G}(H)$.
%\end{cor}

%\begin{cor} \label{f conjugacy familia F} Let $H$ be a finite subgroup of a $\Oe$-group $G_1$ and $G=HNN (G_1, H, f, t)$ be an HNN-extension with $f \in \aut(H).$  Suppose  $f$ belongs to the $\overline{Aut}_{G_1}(H)$-conjugacy class of  $\overline{N}_{G_1}(H)$. Then $|g(G, \mathcal{NF})|=k(G_1,H).$  In particular, this is the case if $f\in \overline{N}_{G_1}(H)$.  
%\end{cor}

By Remarks \ref{rem dos normalizadores e densidade}, \ref{tilde implica overline}, \ref{ki finitos} and Corollary \ref{autG classe}, we have the following

\begin{cor}\label{densohnn} Let $H$ be a finite subgroup of an $\Oe$-group $G_1$ and $$G=\hnn(G_1, H, f, t)$$ be a normal HNN-extension. Suppose that  $$\widetilde{\aut}_{G_1}(H) = \widetilde{\aut}_{\widehat G_1}(H), \widetilde{N}_{\widehat G_1}(H)=\widetilde{N}_{ G_1}(H)$$ and $k(G_1, H, H)=1.$ Then $$|g(G, G_1,\A)|=|g(G, \mathcal{F})|\leq |\tilde \kappa(\widetilde{N}_{ G}(H))|.$$ In particular, if $\widetilde{N}_{\widehat G}(H)=\widetilde{N}_{G_1}(H)$ or $|\out(H)|\leq 2$ then $|g(G, G_1,\A)|=1$. 
\end{cor}


\begin{exa} Let $P=\langle x_1, x_2\rangle$ be a  free nilpotent group of rank $2$ and \break $H\cong \Z/p \Z$ be  a group of odd prime order $p$. Note that $$\aut(\Z/p \Z)=\langle c\rangle\cong C_{p-1}$$ is a  cyclic group of order  $p-1$.   We look at $c$ as the unit of $\Z/p\Z$ and define $G_1=H\rtimes_{\psi_1} P$ with $h^{x_j}= h^{c}$, for $i,j=1,2$ and by Proposition 3.1 in \cite{BPZ2} $G_1$ is an $\Oe$-group. Note that $\widehat G_1=H\rtimes_{\hat \psi_1} \widehat P,$ where $\hat \psi_1$ is the natural extension of $\psi_1$ to $\widehat P$ (see Proposition 2.6 in \cite{GZ}). Let $G=\hnn(G_1,H,f,t)$ be an HNN-extension with $f\in \aut(H)$. Since $H$ is characteristic in $\widehat G_1$, by Remark \ref{ki finitos} $k(G_1, H, H)=1$. Furthermore, $\widetilde{N}_{G_1}(H)=\out(H)=\widetilde{N}_{\widehat G_1}(H)=\widetilde{N}_{\widehat{G}}(H)$. By Example 7.11 in \cite{BPZ2} we have $\widetilde{\aut}_{ G_1}(H)=\widetilde{\aut}_{\widehat{G}_1}(H)=\aut(H)=\out(H).$ It follows from Corollary \ref{densohnn} that $|g(G,\mathcal{F})|=|g(G,G_1,\A)|=1$. 
\end{exa}

\begin{exa} Let $G_1 = GL_n(\mathbb{F}_p)$ and consider $H$ to be the Borel subgroup of the upper triangular matrices of $G_1$. Let $G=\hnn(G_1, H, f, t)$ be an HNN-extension with $f \in \aut(H).$ Note that $H$ is self-normalizing in $G_1$ so $\widetilde{N}_{G_1}(H)=\{1\},$ moreover $k(G_1, H, H)=1$ since $G_1$ is finite. By Proposition \ref{cp1} we have that $\widetilde{N}_{\widehat{G}}(H)=\left \langle \tilde{f} \right \rangle,$ where $\tilde{ f}$ is the image of $f$ in $\out(H).$ It follows from Corollaries \ref{nova conta conjugadoss}, \ref{cotas n} and  Remark \ref{ki finitos} %Theorem \ref{genus hnn nao fixos} 
that  $$|g(G,G_1,\A)|=|g(G,\mathcal{F})|= |\mathcal{R}/(\widetilde{\aut}_{ G_1}(H) \times C_2)| \leq \left| \left \langle \tilde{f} \right \rangle / (\widetilde{\aut}_{ G_1}(H) \times C_2) \right|,$$ and therefore %$$|g(G,G_1,\A)| \leq \dfrac{\phi(|\tilde f|)+1}{2}.$$ 

$$|g(G,G_1,\A)| \leq \left\{\begin{array}{cc}
  1,&\mbox{ if } |\tilde f| \leq 2,\\
\displaystyle\frac{\phi(|\tilde f|)}{2}, &\mbox{ if } |\tilde f| \geq 3,
\end{array}\right.$$
where $\phi$ is Euler's function and $\tilde f$ is the image of $f$ in $\out(H)$.
In particular, if $f \in \inn(H),$  then $|g(G,G_1,\A)|=1.$    
\end{exa}

\begin{exa}\label{dieadra ove klein} Let $G_1= \langle r, c \ | \ r^4=c^2=1, crc^{-1}=r^{-1}\rangle \cong D_8$  and $H=\langle c, r^2\rangle \cong C_2 \times C_2$ be one of the Klein subgroups of $G_1.$ Let $f_1$ be an automorphism of $H$ that fixes $c$ and maps $r^2$ to $cr^2$ and $f_2$ an automorphism that fixes $r^2$ and maps $c$ to $cr^2$. Let $B=\hnn(G_1, H,  f_1,t_1)$ and $D=\hnn(G_1, H, f_2,t_2)$  be normal HNN-extensions. Note that $$\aut(H)=\out(H)\cong S_3, \overline{N}_{G_1}(H)=\{id, f_2=\tau_{cr}|_{_{H}}\} = \overline{\aut}_{G_1}(H) \cong C_2.$$ By Proposition \ref{nonnormalhnnabs}  $B$ and $D$ are not isomorphic 
and by Corollary \ref{corhnn}  the number of the isomorphism classes of HNN-extensions $G=\hnn(G_1, H, f,t)$ with $f\in \aut(H)$ is equal to 2 since the two distinct orbits are given by the orbits of the representatives $f_1$ and $f_2.$ By Theorem \ref{genus 1 normal}  $$|g(D, \F)|=1$$ because $f_2 \in \overline{N}_{G_1}(H)$. Consequently, $\widehat B$ and $\widehat D$
are non-isomorphic and $|g(B,\F)|=1$.
As $G_1$ is finite it follows from Remark \ref{ki finitos} that $$k(G_1, H, H)=1,$$ therefore the genus satisfies $$|g(D, \F)|=|g(B,\F)|=|g(D,G_1,\A)|=|g(B,G_1,\A)|=|g(G,G_1,\A)|=1.$$
\end{exa}


\begin{exa} In the previous example take $H=\langle r^2,  c \rangle \cong C_2 \times C_2$ and $K=\langle r^2, r c \rangle \cong C_2 \times C_2$ the distinct Klein subgroups of $G_1$ and let $$G=\hnn(G_1, H, K, f, t)$$ be a non-normal HNN-extension where $f \in \Iso(H, K)$ fixes $r^2$ and sends $c$ to $rc.$ Note that $f$ extends to an automorphism $\varphi \in \aut(G_1)$ where $\varphi(r)=r$ and $\varphi(c) = rc,$ therefore it follows from  Theorem \ref{genus 1} i) and Remark \ref{ki finitos} that $|g(G, G_1,\A)|=1.$ 

Similarly let $G_2=Q_8=\left \langle i, j\,|\, i^4=j^4=1, i^j=j^{3}=-j \right \rangle$ be the quaternion group of order 8, $H=\left \langle i \right \rangle$ and $K=\left \langle j \right \rangle $ two of its non-conjugate subgroups of order 4. Let $W=\hnn(G_2, H, K, f, t)$ such that $f(i)=\pm j.$ Note that $f$ extends to an automorphism $\psi: G_2 \longrightarrow G_2$ such that $\psi(i)=j$ and $\psi(j)=i$ or $\psi(i)=-j$ and $\psi(j)=-i.$ Then for the same reason as in the dihedral case we have  $|g(W, G_2,\A)|=1,$ and similarly for any permutation between three distinct subgroups of order 4 of the quaternion group.
\end{exa} 

\begin{exa}\label{diedral over cyclic}
In the previous example, let $H$ be any of the proper subgroups of $G_1 \cong D_8$ that are not Klein subgroups. Then $H$ is cyclic of order $\leq 4,$ from where it follows that $|\aut(H)|=|\out(H) | \leq 2,$ so $$|g(\hnn(D_8, H, H, f,t), D_8,\A)|=1$$ by Corollary \ref{autG classe} and Remark \ref{ki finitos}.  Therefore, for any HNN-extension $$G=\hnn(D_8, H, H, f,t)$$ the cardinality of  genus $g(G, G_1,\A)$ is always equal to 1 regardless of the chosen associated subgroup. If $G_2=Q_8$ is the Quaternion group of order $8,$ the same follows because all its proper subgroups are cyclic of order $\leq 4,$ i.e., $$|g(\hnn(Q_8, H, H, f,t), Q_8,\A)|=1$$ independently of the chosen associated subgroups.     
\end{exa}

\begin{rem} Note that $$G=\hnn(D_8, H, K, f,t),\  {\rm and}\  G=\hnn(Q_8, H, K, f,t)$$ are  profinitely rigid if the Remeslennikov's question has a positive answer. Indeed,  a finitely generated residually finite group $B$  having the same profinite completion as $G$ would be virtually free and by \cite[Theorem 1.1]{BPZ} $B$ will be an HNN-extension with finite base group.
\end{rem}



\subsection{Genus in a more general class} \label{general section}


 In this subsection we give formulas  for the genus of HNN-extensions in class $\mathcal{A}$ with $G_1$ not fixed. When $G$ is a group in $\mathcal{A}$,  we will denote the genus of $G$ in class $\mathcal{A}$, as $g(G, \mathcal{A}).$

Let $H$ and $K$ be finite subgroups of an $\Oe$-group $G_1$ and $$G=\hnn(G_1, H, K, f, t)$$ an HNN-extension. The objective  is to find a formula to calculate the cardinality of the set $g(G, \mathcal{A}).$
To do this, consider $B\in \mathcal{A}$ such that $\widehat B\cong \widehat G$. Then by \cite[Theorem 1.1]{BPZ} $B$ is an HNN-extension of the form $B=\hnn(B_1, H_1, K_1, f_1, t_1)$ with $\widehat B_1\cong \widehat G_1, H \cong H_1\cong K \cong K_1,$ and furthermore $B_1 \in \mathcal{A}$.


%\begin{deff} Let $G_1$ be a group and $\widehat G_1$ be its profinite completion. Let $X$ and $Y$ be finite subgroups of $\widehat G_1$. We say that the pair $(X,Y)$ is \textbf{permutable}, if there exists $\psi\in\aut(\widehat G_1)$ that  permutes the conjugacy classes of $X$ and $Y$.\end{deff}


 \begin{deff} Let $G_1$ be a group and $\widehat G_1$ be its profinite completion. Let $X_1, X_2, Y_1, Y_2$ be finite subgroups of $\widehat G_1$. We say that the pairs $(X_1,Y_1)$ and $(X_2,Y_2)$ are \textbf{permutable}, if there exists $\psi\in\aut(\widehat G_1)$ that  permutes the conjugacy classes of $X_i$ and $Y_i,$ for each  $i=1,2$.\end{deff} 
 
Let $B_1 \in g(G_1,\A)$. We will identify $\widehat G_1$ with $\widehat B_1$ and consider $G_1$ and $B_1$ as dense subgroups of $\widehat G_1=\widehat B_1.$ 
  Let $B=\hnn(B_1,H_1,K_1,f_1,t_1)$ be an HNN-extension with finite associated subgroups. We say that $B$ is \textbf{ compatible }  with $G$ if the pairs $(H,H_1)$ and $(K,K_1)$ are permutable in $\widehat G_1$ (i.e. exists $\psi\in\aut(\widehat G_1)$ such that $\psi(H)=H_1$ and $\psi(K)=K_1^{g_1}$ for some $g_1\in \widehat G_1$) and $f_1 = \tau_{g_1} f^{\psi^{-1}}$.
By Proposition \ref{2a}, $B$ is compatible with $G$ if and only if $\widehat B\cong \widehat G.$ 
%If $B$ is compatible to $G$, then by Proposition \ref{2a} we have $\widehat B\cong \widehat G$.
With these considerations, we state the following result


\begin{thm}\label{fixed H2} The cardinality of the genus $g(G, \mathcal{A})$ of an HNN-extension $G=\hnn(G_1, H, K, f, t)$ is equal to $$|g(G,\A)| = \sum_{B_1\in g(G_1,\A)} |g\left(\hnn(B_1, H_1, K_1, f_1, t_1  ), B_1,\A\right)|,$$ where $B_1$ ranges over the representatives of isomorphism classes in $g(G_1,\A)$ such that $G$ is compatible with $\hnn(B_1, H_1, K_1, f_1, t_1 ).$   
\end{thm}



%\textcolor{red}{Observe que se $B_1 \in \mathfrak{g}(G_1,\A)$   mas não existirem subgrupos finitos $H_1, K_1$ de $B_1$ satisfazendo a condição $A$ então esse $B_1$ deve ser desconsiderado no somatorio, pois não haverá nenhuma HNN extensão com grupo base $B_1$ e subgrupo associado finito isomorfo a $H$ cujo completamento é $\widehat G$. Para ver isso, basta supor por contradição que existe essa HNN e observar a contradição com a Proposição \ref{2a}}

\begin{cor} \label{genus 1 gerall}
If $| g(G_1,\A)|=1$ then $|g(G, \mathcal{A})|=|g(G, G_1,\A)|.$    
\end{cor}

\begin{exa} 
 Let $\Gamma$ be an indecomposable into a free product finitely generated  Fuchsian group and $H \neq \{1\}$ is one of its finite subgroup. Let $G_1=K\times \Gamma$ such that $H \cong K$. Consider $G=\hnn(G_1,H, K, f, t)$. If $\Gamma$ is profinitely rigid (see \cite{BMRS20} for examples of such) then so is $G_1$ and by Example \ref{usar genus class A} and Corollary \ref{genus 1 gerall}  $|g(G,\mathcal{A})|=1.$  
\end{exa}






\nocite{*}

\begin{thebibliography}{99}
%\section*{References}
%\bibliographystyle{line}
%\bibliography{JAMS-paper}

%\begin{thebibliography}{50}
%\bibitem[BAUM]{BAUM} B. Baumslag and M. Tretkoff, {\it Residually finite HNN-extensions.} Communications in Algebra, 6:2, (1973) 179-194.

\bibitem{Aka} M. Aka. {\it Arithmetic groups with isomorphic finite quotients.} Journal of Algebra, 352 (2012) 322-240.

\bibitem{A} S. Andreadakis. {\it On the automorphisms of free groups and free nilpotent groups.} Proc. London Math. Soc. 15 (1965) 239-268.

\bibitem{BAU74} G. Baumslag. {\it Residually finite groups with the same finite images.} Compositio Mathematica, 29 (\textbf{3}) (1974) 249-252.

\bibitem{borelarith} A. Borel. \textit{ Arithmetic properties of linear algebraic groups.}
Arithmetic properties of linear algebraic groups.   Proc. ICM Stockholm. (1962) 10-22. 
%\bibitem[BF]{BF} M. Bestvina and M. Feighn. Bounding the complexity of simplicial group actions on trees. Inventiones Mathematicae, vol. 103 (1991), no. 3, pp. 449–469.

%\bibitem[BL]{BL} David El-Chai Ben-Ezra, Alexander Lubotzky, The Congruence Subgroup Problemfor finitely generated Nilpotent Groups. arxiv.org/pdf/2005.03263.pdf.


\bibitem{BPZ2}  V.R. de Bessa, A.L.P. Porto, P.A. Zalesskii. {\it Profinite genus of free products with finite amalgamation.} Journal of Algebra, {\bf 643} (2024) 11-48.

\bibitem{BPZ}  V.R. de Bessa, A.L.P. Porto, P.A. Zalesskii. {\it The profinite  completion of accessible groups.} Monatsh. Math. v. 202 {\bf 2} (2022) 217-227.

\bibitem{BZ} V.R. Bessa, P.A. Zalesskii. {\it The genus for HNN-extensions.} Mathematische Nachrichten, v. 286, (2013) 817-831.

\bibitem{BCR16}  M.R. Bridson, M.D.E. Conder, A.W. Reid. {\it Determining Fuchsian groups by their finite quotients.} Israel Journal of Mathematics, \textbf{214} (2016) 1-41.
	
\bibitem{BMRS18} M.R. Bridson, D.B. McReynolds, A.W. Reid, R. Spitler. \textit{ Absolute profinite rigidity and hyperbolic geometry}. Annals of Mathematics,  {\bf 192} (2020) 679--719. arXiv preprintarXiv:  1811.04394, 2018.
	
\bibitem{BMRS20} M.R. Bridson, D.B. McReynolds, A.W. Reid, R. Spitler. \textit{On the profinite rigidity of triangle groups}, Bulletin of the London Math. Society, {\bf 53} (2021) 1849--1862. arXiv: 2004.07137, 2020.



\bibitem{Brid} M.R. Bridson. \textit{ Finiteness properties for subgroups of $GL(n, \mathbb{Z})$.} Math. Ann. 317 {\bf(4)} (2000) 629-633.

\bibitem{bridcurvature} M.R. Bridson, A. Haefliger. \textit{Metric spaces of non-positive curvature.} Grundlehren der math Wiss. 319, Springer-Verlag, Berlin. Heidelberg. New York (1999).


\bibitem{Dav} D. Carolillo, G. Paolini. \textit{Profinite rigidity of crystallographic groups arising from Lie theory.} arXiv:2506.15494

\bibitem{Sam} S. M. Corson, S. Hughes, P. M\"oller, O. Varghese. \textit{Profinite rigidity of affine Coxeter groups.} Mathematische Zeitschrift, 311 \textbf{(11)} (2025) 1--9. 

\bibitem{Sam2} S. M. Corson, S. Hughes, P. M\"oller, O. Varghese. \textit{Higman-Thompson groups and profinite properties
of right-angled Coxeter groups.} (2023) arXiv:2309.06213


%[BO]? J. Boler and B. Evans, {\it The free product of residually finite groups amalgamated along retracts is residually finite.} Proceedings of the American Mathematic Society, v. 37, Number 1, 1973.

%\bibitem{bourbaki} N. Bourbaki. {\it Elements of Mathematics: General Topology, Pt.2.} Addison-Wesley Educational Publishers Inc.,U.S., (1967).  

%\bibitem[CR]{CR}. W. Curtis and I. Reiner, Methods  of representation theory  with applications  to finitegroups and ordes, John Wiley and Sons, New York (1981).


%\bibitem[D]{D} M.J. Dunwoody, {\it The accessibility of finitely presented groups.}  Inventiones mathematicae, 81 (\textbf{3}) (1985) 449-457.

\bibitem{culler} M. Culler. {\it Finite groups of automorphisms of free groups.} In: Contributions to Group Theory, Contemp. Math., 33 (1984) 197-207. Amer. Math. Soc., Providence, R.I.

\bibitem{DI} W. Dicks, M.J. Dunwoody. Groups acting on graphs. Cambridge Studies in Advanced Math. {\bf 17}, Cambridge University Press, 1989.

\bibitem{Die} J. Dietz. {\it Automorphism of groups of semi-direct products.}  Communications in Algebra {\bf 40} (2012) 3308-3316.  

\bibitem{fine} B. Fine, G. Rosenberger. {\it Conjugacy separability of Fuchsian groups and related questions}.  Contemporary Mathematics {\bf 109} (American Mathematical Society, Providence), (1990) 11-18.

\bibitem{GPS} F.J. Grunewald, P.F. Pickel, D. Segal. {\it Finiteness theorems for polycyclic groups.} Bull. Amer. Math. Soc. (N.S.) 1 (1979), no. 3, 575-578.

\bibitem{GS} F. Grunewald, R. Scharlau. {\it A note on finitely generated torsion-free nilpotent groups of class 2.} J. Algebra 58 (\textbf{1}) (1979) 162-175.

\bibitem{GZ} F.J. Grunewald, P.A. Zalesskii. {\it Genus for groups.} Journal of Algebra 326, (2011) 130-168.

%\bibitem{haran} D. Haran, M. Jarden, F. Pop. {\it $P$-adically Projective Groups as Absolute Galois Groups.} International Mathematics Research Notices 32, (2005). 

\bibitem{J} A. Jaikin-Zapirain. The finite and solvable genus of finitely generated free and surface groups. Res. Math. Sci. 10, 44 (2023). 

\bibitem{J2} A. Jaikin-Zapirain, I. Morales. {\it Prosolvable rigidity of surface groups.} Sel. Math. New Ser. 31, 86 (2025). https://doi.org/10.1007/s00029-025-01082-1

\bibitem{J3} A. Jaikin-Zapirain, A. Lubotzky. {\it Some remarks on Grothendieck pairs.} Groups, Geometry, and Dynamics, 19 (2), (2025) 597-616. https://doi.org/10.4171/GGD/889


%\bibitem{lennox} J. C. Lennox, S. E.  Stonehewer. { \textit Subnormal Subgroups of Groups. } Oxford Mathematical Monographs. New York (1987).

%\bibitem[KAR]{KAR} J.I. Merzljakov and M.I. Kargapolov, {\it Fundamentals of the theory of groups.} Second Edition. Springer-Verlag, New York Inc. (1979).

%\bibitem[KAR]{KAR} J.I. Merzljakov and M.I. Kargapolov, {\it Fundamentals of the theory of groups.} Second Edition. Springer-Verlag, New York Inc. (1979).

%\bibitem[LIM]{LIM} C. K. Lim, {\it Free pro-C Groups.} PhD Thesis, McGill University, Montreal, Canada (1971).

\bibitem{K}  E. I. Khukhro, V. D. Mazurov (eds.). The Kourovka notebook. Unsolved problems in group theory, 18th ed., Ross. Akad. Sci. Sib. Div., Inst. Math., Novosibirsk 2014, 227 pp.

\bibitem{L} A. Lubotzky. {\it Combinatorial group theory for pro-p-groups.} J. Pure and Appl. Algebra 25 (1982) 311-325.

%\bibitem[LI]{LI} J. Landin, I. Reiner, {\it Automorphisms of the general %linear group over a principal ideal domain.} Annal of Mathematics %\textbf{65} (1957) 519-526.

\bibitem{LS} R.C. Lyndon, P.E. Schupp. Combinatorial group theory. Reprint of the (1977) edition. Classics in Mathematics. Springer-Verlag, Berlin, 2001.

\bibitem{MKS} W. Magnus, A. Karrass, D. Solitar. Combinatorial group theory: Presentations of groups in terms of generators and relations. Second edition, Dover Publication, INC. New York, 1976.

\bibitem{M} I.  Morales, On the profinite rigidity of free and surface groups. Mathematische Annalen, Volume 390, pages 1507–1540, (2024).

%\bibitem{MM} J.M. Masley, H.L. Montgomery. {\it Cyclotomic fields with unique factorization.} J. reine angew. Math. (\textbf{286-287}) (1976) pp. 248-256.

\bibitem{Ner19} G.J. Nery, \textit{Profinite genus of fundamental groups of torus bundles.} Communications in Algebra, \textbf{48} (2019) 1567-1576.

\bibitem{Ner20} G.J. Nery. {\it Profinite genus of fundamental groups of compact flat manifolds with holonomy group of prime order.} Journal of Group Theory, \textbf{24\,(6)}, 1135-1148. 

\bibitem{Ner24} G.J. Nery. {\it Profinite genus of fundamental groups of compact flat manifolds with cyclic holonomy group of square-free order.} Forum Mathematicum, 2024. 

%\bibitem{N} W.K. Nicholson, Introduction to abstract algebra. John Wiley and Sons, 2012.

%\bibitem[PET]{PET} M. R. Pettet, {\it Virtually free groups with finitely many outer automorphisms.} Transactions of the American Mathematical Society, v. 349, {\bf 11} (1997) 4565-4587.

%\bibitem[PIC]{PIC} P. F. Pickel, {\it Fitting subgroups and profinite completions of polycyclic groups.} Journal of Algebra {\bf 42} (1976) 41-45.

%[RE] A. W. Reid, {\it Profinite properties of discrete groups.} Conference: Groups St Andrews, St Andrews (2013).

%\bibitem{NS-07} N. Nikolov, D. Segal.  {\it On finitely generated profinite groups, II. Products in quasisimple groups.} Ann. of Math. (2) {\bf 165} (2007), 239-273.
\bibitem{newman} M. Newman. {\it A note on Fuchsian groups.} Illinois Journal of Mathematics, v. 29, {\bf 
 4} (1985) 682-686.

\bibitem{R} L. Ribes. Profinite graphs and groups. Springer-Verlag {\bf 66}, 2017.

\bibitem{Sta} J. R. Stallings. Group theory and three-dimensional manifolds, Yale University Press, New Haven, Conn.-London, 1971, A James K. Whittemore Lecture in Mathematics given at Yale University, 1969, Yale Mathematical Monographs, 4. 

\bibitem{RO} D.J.S. Robinson. A course on the theory of groups. Springer-Verlag, New York, Berlin-Heidelberg, 1982.

\bibitem{RSZ} L. Ribes, D. Segal, P.A. Zallesskii. {\it Conjugacy separability and free products of groups with cyclic amalgamation.} J. London Math. Soc. 57 (\textbf{2}), (1998) 609-628.

\bibitem{RZ} L. Ribes, P.A. Zalesskii. Profinite groups. Second edition, Springer-Verlag, Berlin Heidelberg, 2010.

%\bibitem{Sc} B. Schmidt. Characters and cyclotomic fields in finite geometry.   Berlin, Springer, 2002 (Lecture notes in mathematics, 1797) ISBN 3-540-44243-X. 

\bibitem{Se} D. Segal, Polycyclic groups (Cambridge University Press, 1983).

\bibitem{S} J.P. Serre. Trees. Springer-Verlag, 2003.

\bibitem{Ste72} P.F. Stebe. {\it Conjugacy separability of groups of integer matrices.} Proceedings of the American Mathematical Society, 32 (\textbf{1}) (1972)  1-7.


\bibitem{vog} K. Vogtmann. {\it Automorphisms of Free Groups and Outer Space.}, Geometriae Dedicata 94, (2002) 1-31. 
%\bibitem{T} O. Taussky. {\it On matrix classes corresponding to an ideal and its inverse.} Illinois J. Math. \textbf{1} (1957), no. 1, 108-113.

\bibitem{Wil17} G. Wilkes. \textit{Profinite rigidity for Seifert fibre spaces.} Geometriae Dedicata, \textbf{188} (2017) 141-163.

%\bibitem{Z} P.A. Zalesskii. {\it Profinite groups that act on trees and do not have free nonabelian pro-$p$-subgroups.} {\empty Math. USSR Sbornik} {\bf 69} (1991) 57-67.

\bibitem{ZM} P.A. Zalesskii, O.V. Melnikov. {\it Fundamental Groups of Graphs of Profinite Groups.} {\empty Algebra i Analiz} {\bf 1} (1989); translated in: {\empty Leningrad Math. J.} {\bf 1} (1990) 921-940.

\bibitem{ZM1} P.A. Zalesskii, O.V. Melnikov. {\it Subgroups of profinite groups acting on trees.} {\empty Math. USSR Sbornik} {\bf 63} (1989) 405-424.


\bibitem{zimm} B. Zimmermann. {\it Finite groups of outer automorphisms of free groups.} Glasgow Mathematical Journal,  Vol. 38, Issue 3 (1996) 275-282. 





\end{thebibliography}
%\end{thebibliography}


\end{document}
